% Les secteurs d'activité — Allemand 1ère A — Cours signé OnBuch+
% Programme officiel MINESEC. Compiler avec : tectonic ALA07-les-secteurs-d-activite.tex
\newcommand{\DOCMATIERE}{Allemand}
\newcommand{\DOCNIVEAU}{1ère A}
\newcommand{\DOCMODULE}{Chapitre 2 : Vie économique}
\newcommand{\DOCLECON}{ALA07}
\newcommand{\DOCTITRE}{Les secteurs d'activité}
% OnBuch+ — préambule « cours signés » ENRICHI (compilé avec tectonic / XeLaTeX)
% Système visuel « Pop » d'OnBuch+ : fond crème (jamais blanc), bordures encre
% épaisses, ombres « dures » décalées, titres Archivo Black en capitales.
% Version MATHÉMATIQUES : graphes (pgfplots/tikz), boîtes pédagogiques
% (définition, propriété, méthode, attention, le savais-tu, exercice résolu,
% exercice, TP, bilan), page de garde et signature OnBuch+ ancrée en bas.
%
% Macros attendues dans main.tex AVANT \input{preamble} :
%   \DOCMATIERE  \DOCNIVEAU  \DOCMODULE  \DOCLECON (numéro)  \DOCTITRE
\documentclass[11pt]{article}

\usepackage{fontspec}
\usepackage[ngerman,french]{babel} % français = langue principale ; l'allemand via \all{..} / textebox
\usepackage[a4paper,margin=2cm,top=3.4cm,bottom=2.8cm,headheight=1.4cm,footskip=1.3cm]{geometry}
\usepackage{amsmath,amssymb,amsfonts}
\usepackage{mathtools} % \xmapsto, \coloneqq... souvent utilisés par les modèles
\usepackage{cancel} % simplifications barrées (bilans de forces)
% \abs{z} (module, valeur absolue) et \norm{u} (norme), utilisés par les modèles
\providecommand{\abs}{}\let\abs\relax\DeclarePairedDelimiter{\abs}{\lvert}{\rvert}
\providecommand{\norm}{}\let\norm\relax\DeclarePairedDelimiter{\norm}{\lVert}{\rVert}
\usepackage[table]{xcolor} % [table] : fournit aussi \rowcolors (alternance de lignes)
\usepackage{pagecolor}
\usepackage{tikz}
\usetikzlibrary{arrows.meta,positioning,calc,shapes.geometric,shapes.misc,shapes.arrows,decorations.pathmorphing,decorations.pathreplacing,decorations.markings,patterns,shadows,mindmap,trees,fit,backgrounds,matrix,intersections,angles,quotes}
\usepackage{pgfplots}
\usepackage[version=4]{mhchem} % équations nucléaires \ce{^{235}_{92}U}
\usepackage[european,siunitx]{circuitikz} % schémas électriques
\pgfplotsset{compat=1.18}
\usepgfplotslibrary{fillbetween,groupplots} % aires entre courbes (calcul intégral)
\usepackage{enumitem}
\usepackage{fancyhdr}
% [most] charge toutes les bibliothèques tcolorbox (listings, minted, raster,
% documentation...) et gonfle inutilement la mémoire TeX sur un document long
% (~300 boîtes) : on ne charge que ce qui est réellement utilisé.
\usepackage[skins,breakable]{tcolorbox}
\usepackage{graphicx}
\usepackage{colortbl}
\usepackage{booktabs}
\usepackage{tabularx}
\usepackage{diagbox} % case d'en-tête diagonale des tableaux à double entrée
% Colonnes extensibles centrée / gauche / droite (C, L, R), souvent utilisées par les modèles
\newcolumntype{C}{>{\centering\arraybackslash}X}
\newcolumntype{L}{>{\raggedright\arraybackslash}X}
\newcolumntype{R}{>{\raggedleft\arraybackslash}X}
\usepackage{array}
\usepackage{multicol}
\usepackage{multirow}
\usepackage{siunitx}
% parse-numbers=false : accepte aussi \SI{2,5 \times 10^{-3}}{...}
\sisetup{locale=FR,output-decimal-marker={,},group-minimum-digits=4,parse-numbers=false}
\DeclareSIUnit\ohm{\ensuremath{\symup{\Omega}}} % U+2126 absent de la police
\DeclareSIUnit\division{div} % réglages d'oscilloscope (ms/div, V/div)
\DeclareSIUnit\tr{tr} % tours (vitesse de rotation, ex. \SI{1500}{\tr\per\minute})
\DeclareSIUnit\molar{M} % concentration molaire (ex. \SI{3}{\molar}, \SI{1}{\micro\molar})
\DeclareSIUnit\an{an} % année (le modèle confond parfois avec \year, un compteur TeX) — ex. \SI{5}{\kilo\meter\per\an}
\DeclareSIUnit\FCFA{FCFA} % franc CFA (ex. \SI{500}{\FCFA\per\kilo\gram})
\DeclareSIUnit\degreeBrix{\textdegree Brix} % teneur en sucre d'un sirop/jus (réfractométrie)
\DeclareSIUnit\month{mois} % le modèle utilise parfois \month, un compteur TeX, comme unité de durée
\usepackage{qrcode}
% \degree : commande fréquemment utilisée par les modèles pour un angle ou une
% température en dehors de \SI{}{} (non fournie par les paquets chargés).
\providecommand{\degree}{^{\circ}}
% \qed (fin de démonstration), utilisé par les modèles sans amsthm
\providecommand{\qed}{\hfill$\square$}
% \barre : double barre des valeurs interdites dans un tableau de variations (tkz-tab)
\providecommand{\barre}{\ensuremath{\|}}
% environnement proof (amsthm non chargé) utilisé par les modèles
\ifdefined\proof\else\newenvironment{proof}[1][Démonstration]{\par\noindent\textit{#1.}\ }{\qed\par}\fi
\usepackage{lastpage}
\usepackage{hyperref}
\hypersetup{hidelinks}

% ============================================================
% Palette « Pop » OnBuch+ — mêmes hex que la classe P (plus_ui.dart)
% ============================================================
\definecolor{poporange}{HTML}{F59321}
\definecolor{popdark}{HTML}{E07A0C}
\definecolor{popcream}{HTML}{FFF6E8}
\definecolor{popink}{HTML}{1C1714}
\definecolor{poppurple}{HTML}{7A5AE0}
\definecolor{popgreen}{HTML}{2E9E6A}
\definecolor{poppink}{HTML}{E0457B}
\definecolor{popblue}{HTML}{2F7DE1}
\definecolor{popgold}{HTML}{C98A12}
\definecolor{popmuted}{HTML}{8A7F76}
\definecolor{popline}{HTML}{EADFD2}
\definecolor{popgray}{HTML}{8A7F76}
\colorlet{popgrey}{popgray}
% Alias tolérants (le modèle nomme parfois les couleurs différemment)
\colorlet{popwhite}{white}
\colorlet{popblack}{popink}
\colorlet{popred}{poppink}
\colorlet{popyellow}{popgold}
% Teintes claires (fonds de tableaux/encadrés) — le modèle nomme parfois
% les couleurs avec un suffixe L (ex. popblueL) sans les avoir définies.
\colorlet{poporangeL}{poporange!20}
\colorlet{popdarkL}{popdark!20}
\colorlet{poppurpleL}{poppurple!20}
\colorlet{popgreenL}{popgreen!20}
\colorlet{poppinkL}{poppink!20}
\colorlet{popblueL}{popblue!20}
\colorlet{popgoldL}{popgold!20}
\colorlet{popredL}{popred!20}
\colorlet{popgrayL}{popgray!20}
% Teintes claires pour fonds de tableaux / zones
\colorlet{popblueL}{popblue!10!white}
\colorlet{poporangeL}{poporange!14!white}
\colorlet{popgreenL}{popgreen!12!white}
\colorlet{poppurpleL}{poppurple!12!white}
\colorlet{poppinkL}{poppink!12!white}
\colorlet{popgoldL}{popgold!14!white}

\pagecolor{popcream}

% ============================================================
% Typographie
% ============================================================
\defaultfontfeatures{Ligatures=TeX}
\setmainfont{PlusJakartaSans-Medium.ttf}[Path=../../../fonts/,BoldFont=PlusJakartaSans-Bold.ttf]
% Maths en sans-serif (Fira Math) avec chiffres et lettres droites de Plus
% Jakarta, pour que les formules se fondent dans le texte.
\usepackage[math-style=ISO,bold-style=ISO]{unicode-math}
\setmathfont{FiraMath-Regular.otf}[Path=../../../fonts/]
\setmathfont{PlusJakartaSans-Medium.ttf}[Path=../../../fonts/,range={up/num,up/{Latin,latin},\mathup}]
\usepackage{icomma} % virgule décimale sans espace en mode math
% Caractères mathématiques Unicode absents des polices de texte (Plus Jakarta,
% Archivo Black) : rendus par leur équivalent mathématique, en texte comme en
% formule — sinon ils s'affichent en carré vide (ex. « Arithmétique dans ℤ »).
\usepackage{newunicodechar}
\newunicodechar{ℝ}{\ensuremath{\mathbb{R}}}
\newunicodechar{ℤ}{\ensuremath{\mathbb{Z}}}
\newunicodechar{ℂ}{\ensuremath{\mathbb{C}}}
\newunicodechar{ℕ}{\ensuremath{\mathbb{N}}}
\newunicodechar{ℚ}{\ensuremath{\mathbb{Q}}}
\newunicodechar{𝔻}{\ensuremath{\mathbb{D}}}
\newunicodechar{α}{\ensuremath{\alpha}}
\newunicodechar{β}{\ensuremath{\beta}}
\newunicodechar{γ}{\ensuremath{\gamma}}
\newunicodechar{δ}{\ensuremath{\delta}}
\newunicodechar{ε}{\ensuremath{\varepsilon}}
\newunicodechar{θ}{\ensuremath{\theta}}
\newunicodechar{λ}{\ensuremath{\lambda}}
\newunicodechar{μ}{\ensuremath{\mu}}
\newunicodechar{σ}{\ensuremath{\sigma}}
\newunicodechar{τ}{\ensuremath{\tau}}
\newunicodechar{φ}{\ensuremath{\varphi}}
\newunicodechar{ψ}{\ensuremath{\psi}}
\newunicodechar{ω}{\ensuremath{\omega}}
\newunicodechar{Σ}{\ensuremath{\Sigma}}
% majuscules produites par \MakeUppercase dans les titres (α-aminés -> Α)
\newunicodechar{Α}{\ensuremath{\alpha}}
\newunicodechar{Β}{\ensuremath{\beta}}
\newunicodechar{Γ}{\ensuremath{\Gamma}}
\newunicodechar{Δ}{\ensuremath{\Delta}}
\newunicodechar{Ε}{\ensuremath{\varepsilon}}
\newunicodechar{Θ}{\ensuremath{\Theta}}
\newunicodechar{Λ}{\ensuremath{\Lambda}}
\newunicodechar{Μ}{\ensuremath{\mu}}
\newunicodechar{Τ}{\ensuremath{\tau}}
\newunicodechar{Φ}{\ensuremath{\Phi}}
\newunicodechar{Ψ}{\ensuremath{\Psi}}
\newunicodechar{Ω}{\ensuremath{\Omega}}
\newunicodechar{↦}{\ensuremath{\mapsto}}
\newunicodechar{∈}{\ensuremath{\in}}
\newunicodechar{∉}{\ensuremath{\notin}}
\newunicodechar{∧}{\ensuremath{\wedge}}
\newunicodechar{⇔}{\ensuremath{\Leftrightarrow}}
\newunicodechar{⟺}{\ensuremath{\Longleftrightarrow}}
\newunicodechar{⇒}{\ensuremath{\Rightarrow}}
\newunicodechar{⟹}{\ensuremath{\Longrightarrow}}
\newunicodechar{∘}{\ensuremath{\circ}}
\newunicodechar{∪}{\ensuremath{\cup}}
\newunicodechar{∩}{\ensuremath{\cap}}
\newunicodechar{≡}{\ensuremath{\equiv}}
\newunicodechar{⊂}{\ensuremath{\subset}}
\newunicodechar{⊥}{\ensuremath{\perp}}
\newunicodechar{∃}{\ensuremath{\exists}}
\newunicodechar{∀}{\ensuremath{\forall}}
\newunicodechar{∛}{\ensuremath{\sqrt[3]{\,}}}
\newunicodechar{∜}{\ensuremath{\sqrt[4]{\,}}}
\newunicodechar{⃗}{\ensuremath{{}^{\to}}}
\newunicodechar{ⁿ}{\ifmmode^{n}\else\textsuperscript{\ensuremath{n}}\fi}
\newunicodechar{ˣ}{\ifmmode^{x}\else\textsuperscript{\ensuremath{x}}\fi}
\newunicodechar{ᵇ}{\ifmmode^{b}\else\textsuperscript{\ensuremath{b}}\fi}
\newunicodechar{ʳ}{\ifmmode^{r}\else\textsuperscript{\ensuremath{r}}\fi}
\newunicodechar{ᵘ}{\ifmmode^{u}\else\textsuperscript{\ensuremath{u}}\fi}
\newunicodechar{ʸ}{\ifmmode^{y}\else\textsuperscript{\ensuremath{y}}\fi}
\newunicodechar{ᵃ}{\ifmmode^{a}\else\textsuperscript{\ensuremath{a}}\fi}
\newunicodechar{ᵉ}{\ifmmode^{e}\else\textsuperscript{\ensuremath{e}}\fi}
\newunicodechar{ᵏ}{\ifmmode^{k}\else\textsuperscript{\ensuremath{k}}\fi}
\newunicodechar{ᵖ}{\ifmmode^{p}\else\textsuperscript{\ensuremath{p}}\fi}
\newunicodechar{ᴺ}{\ifmmode^{N}\else\textsuperscript{\ensuremath{N}}\fi}
\newunicodechar{ᶜ}{\ifmmode^{c}\else\textsuperscript{\ensuremath{c}}\fi}
\newunicodechar{ʰ}{\ifmmode^{h}\else\textsuperscript{\ensuremath{h}}\fi}
\newunicodechar{ᵠ}{\ifmmode^{q}\else\textsuperscript{\ensuremath{q}}\fi}
\newunicodechar{ₙ}{\ifmmode_{n}\else\textsubscript{\ensuremath{n}}\fi}
\newunicodechar{ₐ}{\ifmmode_{a}\else\textsubscript{\ensuremath{a}}\fi}
\newunicodechar{ᵢ}{\ifmmode_{i}\else\textsubscript{\ensuremath{i}}\fi}
\newunicodechar{ᵣ}{\ifmmode_{r}\else\textsubscript{\ensuremath{r}}\fi}
\newunicodechar{ₖ}{\ifmmode_{k}\else\textsubscript{\ensuremath{k}}\fi}
\newunicodechar{ᵦ}{\ifmmode_{\beta}\else\textsubscript{\ensuremath{\beta}}\fi}
\newunicodechar{ₓ}{\ifmmode_{x}\else\textsubscript{\ensuremath{x}}\fi}
\newfontfamily\popdisplay{ArchivoBlack-Regular.ttf}[Path=../../../fonts/]
\newfontfamily\poplogo{SpaceGrotesk-Bold.ttf}[Path=../../../fonts/]
\newfontfamily\popsemi{PlusJakartaSans-SemiBold.ttf}[Path=../../../fonts/]
\newfontfamily\popxbold{PlusJakartaSans-ExtraBold.ttf}[Path=../../../fonts/]
\newfontfamily\popxboldls{PlusJakartaSans-ExtraBold.ttf}[Path=../../../fonts/,LetterSpace=3.0]

\setlist[enumerate]{leftmargin=1.7em,itemsep=5pt,topsep=4pt}
\setlist[itemize]{leftmargin=1.7em,itemsep=4pt,topsep=4pt,label={\color{poporange}\textbullet}}
\setlength{\parindent}{0pt}
\setlength{\parskip}{5pt}
\renewcommand{\arraystretch}{1.25}

% pgfplots : style Pop par défaut
\pgfplotsset{
  every axis/.append style={axis line style={popink,line width=1pt},tick style={popink},
    grid style={popline,line width=0.5pt},label style={font=\small},tick label style={font=\footnotesize}},
  popcurve/.style={poporange,line width=1.8pt,no markers,smooth},
}
% Même style disponible dans un \draw TikZ simple (hors pgfplots)
\tikzset{popcurve/.style={poporange,line width=1.8pt}}
\tikzset{popgrid/.style={popline,very thin}}
\colorlet{popcurve}{poporange} % le nom du style est parfois utilisé comme couleur

% ============================================================
% Pill / squiggle
% ============================================================
\newcommand{\poppill}[3]{%
  \tikz[baseline=-0.55ex]\node[draw=popink,line width=1.2pt,rounded corners=2.6pt,%
    fill=#2,inner xsep=6.5pt,inner ysep=3pt]{%
    \popxboldls\fontsize{7}{7}\selectfont\color{#3}\MakeUppercase{#1}};%
}
\newcommand{\popsquiggle}[1]{%
  \tikz[baseline=-0.4ex]{
    \draw[#1,line width=2.6pt,line cap=round]
      plot [smooth] coordinates {(0,0.09) (0.34,0.01) (0.68,0.21) (1.02,0.01) (1.36,0.10)};
  }%
}
% Mot-clé surligné (marqueur orange sous le texte)
\newcommand{\cle}[1]{{\popxbold\color{popdark}#1}}

% ============================================================
% En-tête / pied
% ============================================================
\pagestyle{fancy}
\fancyhf{}
\renewcommand{\headrulewidth}{0pt}
\renewcommand{\footrulewidth}{0pt}
\lhead{\raisebox{-0.15ex}{\poplogo\fontsize{13}{13}\selectfont\color{popink}OnBuch\color{poporange}+}}
\rhead{\poppill{\DOCMATIERE\ \textbullet\ \DOCNIVEAU\ \textbullet\ Leçon \DOCLECON}{white}{popink}}
\lfoot{\popxbold\fontsize{7.6}{7.6}\selectfont\color{popmuted}\MakeUppercase{Cours signé OnBuch+}\ \textbullet\ {\popsemi\DOCTITRE}}
\rfoot{\tikz[baseline=-0.6ex]\node[rounded corners=8.5pt,fill=popink,minimum height=17pt,minimum width=17pt,inner xsep=5pt,inner sep=0]{\popxbold\fontsize{8}{8}\selectfont\color{popcream}\thepage\,/\,\pageref*{LastPage}};}

% ============================================================
% Titres
% ============================================================
\newcommand{\chaptitre}[2]{%
  \begin{tcolorbox}[enhanced,colback=poporange,colframe=popink,boxrule=2.6pt,arc=11pt,
      drop shadow=popink,left=15pt,right=15pt,top=12pt,bottom=11pt,before skip=3pt,after skip=18pt]
    \poppill{#2}{popink}{popcream}\par\vspace{9pt}
    {\raggedright\hyphenpenalty=10000\exhyphenpenalty=10000\popdisplay\fontsize{22}{24}\selectfont\color{popcream}\MakeUppercase{#1}\par}
  \end{tcolorbox}
}
% Section numérotée : pastille orange avec le numéro + titre Archivo
\newcounter{popsec}
\newcommand{\coursec}[1]{%
  \stepcounter{popsec}\setcounter{popsub}{0}%
  \par\vspace{16pt}\noindent
  \begin{minipage}{\linewidth}
  \tikz[baseline=-0.7ex]\node[draw=popink,line width=1.6pt,fill=poporange,rounded corners=4pt,minimum size=22pt,inner sep=0]{\popdisplay\fontsize{12}{12}\selectfont\color{popcream}\thepopsec};%
  \hspace{8pt}{\popdisplay\fontsize{15}{17}\selectfont\color{popink}\MakeUppercase{#1}}\par
  \vspace{4pt}\noindent{\color{poporange}\rule{36pt}{3.6pt}}
  \end{minipage}\par\nopagebreak\vspace{8pt}
}
\newcounter{popsub}
\newcommand{\courssub}[1]{%
  \stepcounter{popsub}%
  \par\vspace{9pt}\noindent{\popxbold\fontsize{11.5}{13}\selectfont\color{popink}{\color{poporange}\thepopsec.\thepopsub}\hspace{6pt}#1}\par\nopagebreak\vspace{4pt}
}

% ============================================================
% Boîtes pédagogiques — même recette : fond blanc, bordure épaisse colorée,
% ombre dure encre, badge plein. Titre optionnel [#1].
% ============================================================
\tcbset{popbase/.style={breakable,enhanced,colback=white,boxrule=2.3pt,arc=9pt,
  shadow={3pt}{-3pt}{0pt}{popink},left=14pt,right=14pt,top=11pt,bottom=11pt,before skip=11pt,after skip=15pt,
  fontupper=\popsemi\fontsize{10.6}{14.5}\selectfont\color{popink},
  fontlower=\popsemi\fontsize{10.6}{14.5}\selectfont\color{popink}}}
% Coche dessinée (le glyphe ✓ n'existe pas dans les polices du document)
\newcommand{\popcheck}[1]{\tikz[baseline=-0.5ex]\draw[#1,line width=1.6pt,line cap=round,line join=round](-0.13,0.0)--(-0.04,-0.09)--(0.13,0.1);}
\providecommand{\coche}{\popcheck{popgreen}}
\providecommand{\resultat}[1]{\par\smallskip\noindent\fcolorbox{popgreen}{popgreenL}{\parbox{\dimexpr\linewidth-2\fboxsep-2\fboxrule\relax}{\textbf{Résultat.} #1}}\par\smallskip}
\newcommand{\popboxhead}[3]{\poppill{#1}{#2}{white}\ifx\relax#3\relax\else\hspace{6pt}{\popxbold\fontsize{10.6}{12}\selectfont\color{popink}#3}\fi\par\vspace{7pt}}

\newtcolorbox{aretenir}[1][]{popbase,colframe=popgreen,before upper={\popboxhead{À retenir}{popgreen}{#1}}}
% Texte en allemand (document de lecture, dialogue, extrait) : cadre
% d'étude. Un titre optionnel (source, auteur) s'affiche dans l'en-tête.
\newtcolorbox{textebox}[1][]{popbase,colframe=popblue,before upper={\popboxhead{Text}{popblue}{#1}\begin{otherlanguage}{ngerman}},after upper={\end{otherlanguage}}}
% Marques ✓ / ✗ (forme correcte / fautive) souvent utilisées par les modèles
\providecommand{\cross}{\textcolor{poppink}{\ensuremath{\times}}}
\providecommand{\xmark}{\cross}\providecommand{\crossmark}{\cross}\providecommand{\cmark}{\ensuremath{\checkmark}}
% Ligne de réponse / justification à compléter (inventée par les modèles)
\providecommand{\justification}[1]{\par\noindent\textit{Justification~:} \dotfill\par}
% \textipa (package tipa non chargé) : la police ne couvre pas l'API -> simple passage
\providecommand{\textipa}[1]{#1}
% Soulignement (ulem non chargé) : \uline, \ul
\providecommand{\uline}[1]{\underline{#1}}\providecommand{\ul}[1]{\underline{#1}}
% \alert (beamer) inventée par les modèles : texte mis en évidence
\providecommand{\alert}[1]{\textbf{\textcolor{poppink}{#1}}}
% variantes ASCII d'ordinaux écrites en macro
\providecommand{\iere}{\textsuperscript{re}}\providecommand{\ieme}{\textsuperscript{e}}\providecommand{\ere}{\textsuperscript{re}}\providecommand{\eme}{\textsuperscript{e}}\providecommand{\er}{\textsuperscript{er}}\providecommand{\ie}{\textsuperscript{e}}
% Ordinaux français écrits en macro par les modèles : 1\ière, 3\ième
\newcommand{\ière}{\textsuperscript{re}}\newcommand{\ième}{\textsuperscript{e}}
% \exemple (inventée par les modèles) : étiquette « Exemple : »
\providecommand{\exemple}{\textbf{Exemple~:}\ }
% \square utilisable aussi hors mode math (cases à cocher)
\AtBeginDocument{\let\squareMath\square\renewcommand{\square}{\ensuremath{\squareMath}}}
% Mot ou phrase en allemand à l'intérieur d'un paragraphe français
\newcommand{\all}[1]{\ifmmode\text{\textit{#1}}\else\begingroup\selectlanguage{ngerman}\itshape #1\endgroup\fi}
\let\esp\all % alias toléré
\newtcolorbox{exemplebox}[1][]{popbase,colframe=poporange,before upper={\popboxhead{Exemple}{poporange}{#1}}}
\newtcolorbox{definition}[1][]{popbase,colframe=popblue,before upper={\popboxhead{Définition}{popblue}{#1}}}
\newtcolorbox{propriete}[1][]{popbase,colframe=poppurple,before upper={\popboxhead{Propriété}{poppurple}{#1}}}
\newtcolorbox{methode}[1][]{popbase,colframe=popgold,before upper={\popboxhead{Méthode}{popgold}{#1}}}
\newtcolorbox{attention}[1][]{popbase,colframe=poppink,before upper={\popboxhead{Attention}{poppink}{#1}}}
\newtcolorbox{experience}[1][]{popbase,colframe=popblue,colback=popblueL,before upper={\popboxhead{Expérience}{popblue}{#1}}}
% « Le savais-tu ? » : carte violette pleine (contexte, culture, Cameroun)
\newtcolorbox{savaistu}[1][]{popbase,colback=poppurple,colframe=popink,
  fontupper=\popsemi\fontsize{10.4}{14.2}\selectfont\color{white},
  before upper={\poppill{Le savais-tu\,?}{white}{poppurple}\ifx\relax#1\relax\else\hspace{6pt}{\popxbold\color{white}#1}\fi\par\vspace{7pt}}}
% Exercice résolu : énoncé en haut, \tcblower puis la solution sur fond crème
\newcounter{popexo}\newcounter{popexr}
\newtcolorbox{exoresolu}[1][]{popbase,colframe=poporange,colbacklower=popcream,
  segmentation style={popink,line width=1.2pt,dashed},
  before upper={\stepcounter{popexr}\popboxhead{Exercice résolu \thepopexr}{poporange}{#1}},
  before lower={\poppill{Solution}{popink}{popcream}\par\vspace{6pt}}}
% Exercice (non résolu en place) — la correction est dans la partie Corrigés
\newtcolorbox{exercice}[1][]{popbase,colframe=popink,
  before upper={\stepcounter{popexo}\popboxhead{Exercice \thepopexo}{popink}{#1}}}
% Corrigé
\newtcolorbox{corrige}[1][]{popbase,colframe=popgreen,colback=popgreenL,
  before upper={\popboxhead{Corrigé}{popgreen}{#1}}}
% Objectifs / prérequis / bilan
\newtcolorbox{objectifs}{popbase,colframe=popink,colback=poporangeL,
  before upper={\popboxhead{Objectifs de la leçon}{popink}{}}}
\newtcolorbox{prerequis}{popbase,colframe=popink,colback=popblueL,
  before upper={\popboxhead{Prérequis}{popblue}{}}}
\newtcolorbox{bilan}{popbase,colframe=popink,colback=white,boxrule=2.8pt,
  before upper={\popboxhead{Fiche bilan}{poporange}{L'essentiel à connaître pour l'examen}}}
% Figure encadrée avec légende
\newtcolorbox{popfigure}[1][]{enhanced,breakable=false,colback=white,colframe=popline,boxrule=1.4pt,arc=8pt,
  left=10pt,right=10pt,top=10pt,bottom=8pt,before skip=10pt,after skip=12pt,halign=center,
  lower separated=false,halign lower=center,
  fontlower=\popsemi\fontsize{9}{11.5}\selectfont\color{popmuted},
  #1}
\newcommand{\legende}[1]{\par\vspace{6pt}{\popsemi\fontsize{9}{11.5}\selectfont\color{popmuted}#1}}

% ============================================================
% Signature OnBuch+
% ============================================================
\newcommand{\poplogoinline}[1]{{\poplogo\fontsize{#1}{#1}\selectfont\color{popink}OnBuch\color{poporange}+}}
\newcommand{\signatureonbuch}{%
  \begin{tcolorbox}[enhanced,colback=popink,colframe=popink,boxrule=2.2pt,arc=10pt,
      drop shadow=poporange,left=14pt,right=14pt,top=10pt,bottom=10pt,before skip=0pt,after skip=0pt]
    \tikz[baseline=-0.7ex]\node[circle,fill=popgreen,draw=popcream,line width=1.3pt,minimum size=22pt,inner sep=0]{\popcheck{white}};%
    \hspace{9pt}\begin{minipage}[c]{0.62\linewidth}
      {\popxbold\fontsize{12}{13}\selectfont\color{popcream}Signé\ \textbullet\ {\poplogo OnBuch\color{poporange}+}}\par\vspace{2pt}
      {\raggedright\popsemi\fontsize{8}{10.5}\selectfont\color{popline}\MakeUppercase{Équipe pédagogique OnBuch+}\\\MakeUppercase{Conforme au programme officiel MINESEC}\par}
    \end{minipage}\hfill
    {\poplogo\fontsize{22}{22}\selectfont\color{popcream}OnBuch\color{poporange}+}
  \end{tcolorbox}%
}

% ============================================================
% Page de garde
% #1 titre · #2 matière · #3 niveau · #4 module · #5 n° leçon · #6 durée ·
% #7 description / objectifs (texte ou itemize)
% ============================================================
\newcommand{\pagegarde}[7]{%
  \thispagestyle{empty}
  \newgeometry{a4paper,margin=1.7cm,top=1.6cm,bottom=1.4cm}%
  \begin{tikzpicture}[remember picture,overlay]
    \fill[poporange!18] ([xshift=-3.2cm,yshift=-3.6cm]current page.north east) circle (4.6cm);
    \fill[poppurple!14] ([xshift=2.4cm,yshift=7.6cm]current page.south west) circle (3.8cm);
  \end{tikzpicture}%
  {\poplogo\fontsize{30}{30}\selectfont\color{popink}OnBuch\color{poporange}+}\hfill
  \raisebox{0.6ex}{\poppill{Cours signé \textbullet\ Premium}{popink}{popcream}}\par
  \vspace{1pt}\popsquiggle{poppurple}\par
  \vspace{8pt}
  \begin{tcolorbox}[enhanced,colback=poporange,colframe=popink,boxrule=2.8pt,arc=13pt,
      drop shadow=popink,left=18pt,right=18pt,top=12pt,bottom=13pt,after skip=10pt]
    \poppill{#2 \textbullet\ #3}{popink}{popcream}\hspace{5pt}\poppill{#4}{white}{popink}\par\vspace{8pt}
    {\popdisplay\fontsize{14}{14}\selectfont\color{popink}LEÇON #5}\par\vspace{4pt}
    {\raggedright\hyphenpenalty=10000\exhyphenpenalty=10000\popdisplay\fontsize{25}{27}\selectfont\color{popcream}\MakeUppercase{#1}\par}
    \vspace{8pt}
    {\popxbold\fontsize{9.5}{9.5}\selectfont\color{popink}\MakeUppercase{Durée conseillée : #6}}
  \end{tcolorbox}
  \begin{tcolorbox}[enhanced,colback=white,colframe=popink,boxrule=2.2pt,arc=10pt,
      drop shadow=popink,left=16pt,right=16pt,top=10pt,bottom=10pt,after skip=8pt]
    \poppill{Dans cette leçon}{poppurple}{white}\par\vspace{5pt}
    \popsemi\fontsize{9.3}{12.6}\selectfont\color{popink}#7
  \end{tcolorbox}
  \noindent\begin{minipage}[t]{0.487\linewidth}
    \begin{tcolorbox}[enhanced,colback=popgreen,colframe=popink,boxrule=2.2pt,arc=10pt,
        drop shadow=popink,left=11pt,right=11pt,top=10pt,bottom=10pt,height=3.1cm,valign=center]
      \tcbox[colback=white,colframe=popink,boxrule=1.3pt,arc=5pt,boxsep=0pt,nobeforeafter,
        left=3pt,right=3pt,top=3pt,bottom=3pt,valign=center]{\qrcode[height=1.9cm]{https://play.google.com/store/apps/details?id=com.luvvix.onbuch&hl=fr}}%
      \hspace{8pt}\begin{minipage}[c]{0.5\linewidth}
        {\popdisplay\fontsize{11}{12.5}\selectfont\color{white}\MakeUppercase{Télécharge OnBuch}}\par\vspace{4pt}
        {\raggedright\popsemi\fontsize{8.4}{11}\selectfont\color{white}Gratuit \textbullet\ Google Play\\Tuteur IA, annales, concours, résultats\par}
      \end{minipage}
    \end{tcolorbox}
  \end{minipage}\hfill
  \begin{minipage}[t]{0.487\linewidth}
    \begin{tcolorbox}[enhanced,colback=poppurple,colframe=popink,boxrule=2.2pt,arc=10pt,
        drop shadow=popink,left=11pt,right=11pt,top=10pt,bottom=10pt,height=3.1cm,valign=center]
      \tikz[baseline=-0.7ex]\node[circle,fill=white,draw=popink,line width=1.3pt,minimum size=1.6cm,inner sep=0]{\popdisplay\fontsize{24}{24}\selectfont\color{poppurple}+};%
      \hspace{8pt}\begin{minipage}[c]{0.6\linewidth}
        {\popdisplay\fontsize{11}{12.5}\selectfont\color{white}\MakeUppercase{Passe à OnBuch+}}\par\vspace{4pt}
        {\raggedright\popsemi\fontsize{8.4}{11}\selectfont\color{white}Tous les cours signés, Tuteur IA illimité, fiches PDF — onglet OnBuch+ de l'app\par}
      \end{minipage}
    \end{tcolorbox}
  \end{minipage}
  \vspace{8pt}
  \popcontenu
  \vfill
  \signatureonbuch
  \restoregeometry
  \newpage
}

% Rangée « Ce cours contient » (page de garde)
\newcommand{\poptile}[3]{% #1 couleur #2 gros texte #3 légende
  \begin{tcolorbox}[enhanced,colback=#1,colframe=popink,boxrule=2pt,arc=9pt,shadow={2.5pt}{-2.5pt}{0pt}{popink},
      left=6pt,right=6pt,top=9pt,bottom=9pt,halign=center,valign=center,height=1.75cm,width=\linewidth,nobeforeafter]
    {\popdisplay\fontsize{12.5}{13}\selectfont\color{white}\MakeUppercase{#2}}\par\vspace{3pt}
    {\centering\popsemi\fontsize{7.8}{9.5}\selectfont\color{white}#3\par}
  \end{tcolorbox}}
\newcommand{\popcontenu}{%
  \par\noindent{\popxboldls\fontsize{7.5}{7.5}\selectfont\color{popmuted}CE COURS SIGNÉ CONTIENT}\par\vspace{7pt}
  \noindent\begin{minipage}[t]{0.235\linewidth}\poptile{popblue}{Cours}{complet, illustré, pas à pas}\end{minipage}\hfill
  \begin{minipage}[t]{0.235\linewidth}\poptile{poporange}{Méthodes}{et exercices résolus}\end{minipage}\hfill
  \begin{minipage}[t]{0.235\linewidth}\poptile{poppink}{Exercices}{type examen + corrigés}\end{minipage}\hfill
  \begin{minipage}[t]{0.235\linewidth}\poptile{popgreen}{Fiche bilan}{l'essentiel à retenir}\end{minipage}\par}

% Sommaire
\newtcolorbox{sommaire}{enhanced,colback=white,colframe=popink,boxrule=2.2pt,arc=10pt,
  drop shadow=popink,left=18pt,right=18pt,top=13pt,bottom=13pt,before skip=6pt,after skip=20pt,
  before upper={\poppill{Sommaire}{poppurple}{white}\par\vspace{9pt}\popsemi\fontsize{10.6}{14.5}\selectfont\color{popink}}}

% Signature de fin de cours — poussée tout en bas de la dernière page
\newcommand{\findecours}{%
  \par\vfill\nopagebreak
  \begin{center}\popsquiggle{poporange}\end{center}\vspace{4pt}
  \signatureonbuch
}
\AtBeginDocument{\def\llbracket{\mathopen{[\mkern-3.5mu[}}\def\rrbracket{\mathclose{]\mkern-3.5mu]}}} % intervalles d'entiers (glyphe ⟦ absent de la police)
% Glyphes absents de Fira Math : redéfinis à partir de glyphes disponibles
\AtBeginDocument{%
  \def\setminus{\mathbin{\backslash}}\let\smallsetminus\setminus
  \def\vdots{\mathord{\vbox{\baselineskip3.2pt\lineskiplimit0pt\kern4pt\hbox{.}\hbox{.}\hbox{.}}}}%
  \def\ddots{\mathinner{\mkern1mu\raise6.5pt\hbox{.}\mkern2mu\raise3.6pt\hbox{.}\mkern2mu\raise0.7pt\hbox{.}\mkern1mu}}%
  \def\triangle{\mathord{\scalebox{0.9}{$\symup{\Delta}$}}}%
  \def\nabla{\mathord{\raisebox{\depth}{\scalebox{1}[-1]{$\symup{\Delta}$}}}}%
  \def\checkmark{\popcheck{popdark}}%
  \def\star{\mathbin{\ast}}\def\bigstar{\mathord{\ast}}%
  \def\diamond{\mathbin{\scalebox{0.6}{$\Diamond$}}}%
  \def\bigcup{\mathop{\mathchoice{\raisebox{-0.3em}{\scalebox{1.7}{$\cup$}}}{\scalebox{1.25}{$\cup$}}{\cup}{\cup}}\displaylimits}%
  \def\bigcap{\mathop{\mathchoice{\raisebox{-0.3em}{\scalebox{1.7}{$\cap$}}}{\scalebox{1.25}{$\cap$}}{\cap}{\cap}}\displaylimits}%
  \def\lesseqqgtr{\mathrel{\lessgtr}}\def\bigoplus{\mathop{\oplus}\displaylimits}%
}
\providecommand{\ssi}{si et seulement si} % abréviation fréquente
\AtBeginDocument{\providecommand{\wideparen}{\overparen}} % arc de cercle

\begin{document}

\pagegarde{Les secteurs d'activité}{Allemand}{1ère A}{Chapitre 2 : Vie économique}{ALA07}{2 h}{Cette leçon propose la lecture et le commentaire d'un texte allemand sur les trois secteurs d'activité (agriculture, industrie, services), l'acquisition du vocabulaire économique thématique et l'étude des compléments de temps et des propositions subordonnées en wenn. Elle prépare à la compréhension écrite et à la production guidée exigées au Bac.
\vspace{4pt}
\begin{multicols}{2}\small
\begin{itemize}
\item À la fin de cette leçon, je sais comprendre globalement puis en détail un texte allemand sur les secteurs d'activité.
\item À la fin de cette leçon, je sais employer le vocabulaire des trois secteurs économiques avec les articles et les pluriels corrects.
\item À la fin de cette leçon, je sais reconnaître et construire des compléments de temps (Wann? Wie lange? Wie oft?).
\item À la fin de cette leçon, je sais former et utiliser des propositions subordonnées introduites par wenn avec le verbe en fin de proposition.
\item À la fin de cette leçon, je sais distinguer wenn (condition/temps répété) et als (passé unique).
\item À la fin de cette leçon, je sais répondre par écrit à des questions de compréhension en phrases complètes.
\end{itemize}\end{multicols}}

\begin{sommaire}
\begin{multicols}{2}
\begin{enumerate}
\item Texte support : Die Wirtschaft in Deutschland
\item Compréhension du texte
\item Lexique thématique : les secteurs d'activité
\item Grammaire 1 : les compléments de temps
\item Grammaire 2 : les propositions en wenn
\item Méthode du commentaire et production guidée
\item Activité d'intégration
\item Méthodes et exercices résolus
\item Exercices
\item Corrigés des exercices
\item Fiche bilan
\end{enumerate}
\end{multicols}
\end{sommaire}

\begin{objectifs}
\begin{itemize}
\item À la fin de cette leçon, je sais comprendre globalement puis en détail un texte allemand sur les secteurs d'activité.
\item À la fin de cette leçon, je sais employer le vocabulaire des trois secteurs économiques avec les articles et les pluriels corrects.
\item À la fin de cette leçon, je sais reconnaître et construire des compléments de temps (Wann? Wie lange? Wie oft?).
\item À la fin de cette leçon, je sais former et utiliser des propositions subordonnées introduites par wenn avec le verbe en fin de proposition.
\item À la fin de cette leçon, je sais distinguer wenn (condition/temps répété) et als (passé unique).
\item À la fin de cette leçon, je sais répondre par écrit à des questions de compréhension en phrases complètes.
\item À la fin de cette leçon, je sais rédiger un court paragraphe guidé sur l'économie du Cameroun ou de l'Allemagne.
\item À la fin de cette leçon, je sais résumer un texte en 3-4 phrases avec mes propres mots.
\end{itemize}
\end{objectifs}

\begin{prerequis}
\begin{itemize}
\item Les articles définis et indéfinis (der, die, das) et le pluriel des noms (Seconde)
\item La conjugaison des verbes réguliers et irréguliers au présent (Seconde)
\item La place du verbe en position 2 dans la phrase déclarative (Seconde)
\item Les subordonnées en dass et weil avec verbe final (début de 1ère)
\item Le vocabulaire de base des métiers (der Lehrer, der Arzt...) vu en Seconde
\end{itemize}
\end{prerequis}

\newpage

\coursec{Texte support : Die Wirtschaft in Deutschland}

Au Cameroun, ton oncle cultive le cacao à Ebolowa, ta cousine travaille dans une usine de transformation à Douala et ton grand frère est informaticien dans une banque à Yaoundé : sans le savoir, ils appartiennent à trois mondes économiques différents ! En Allemagne aussi, on classe les métiers en trois grands secteurs : \all{die Landwirtschaft} (l'agriculture), \all{die Industrie} (l'industrie) et \all{der Dienstleistungssektor} (le secteur des services). Mais comment les Allemands décrivent-ils leur économie ? Quels mots utilisent-ils pour parler des paysans, des usines et des bureaux ? Cette leçon te permettra de lire, de comprendre et de commenter un texte économique allemand, et de parler des métiers qui t'entourent.

\emph{Problématique :} Comment l'économie allemande est-elle organisée en trois secteurs, et comment lire efficacement un texte allemand sur ce sujet ?

\courssub{Lecture globale : dégager le thème général}

Avant de lire, observe le titre : \all{Die Wirtschaft in Deutschland}. Le mot \all{die Wirtschaft} signifie « l'économie ». Tu sais donc déjà de quoi parle le texte : de l'économie de l'Allemagne. C'est la première étape de toute bonne lecture.

\begin{methode}[Comment lire un texte allemand]
\begin{enumerate}
\item Lis le \cle{titre} et repère les \cle{mots transparents} (mots qui ressemblent au français).
\item Fais une \cle{première lecture} silencieuse, sans dictionnaire, pour dégager le sens global.
\item Fais une \cle{deuxième lecture} avec le glossaire pour comprendre les détails.
\item \cle{Souligne} les mots-clés de chaque paragraphe (qui ? quoi ? où ? combien ?).
\end{enumerate}
\end{methode}

\textbf{Repérage des mots transparents.} Avant même de lire, cherche dans le texte les mots que tu comprends sans dictionnaire :

\begin{center}
\begin{tabularx}{\linewidth}{|X|X|X|}
\hline
\rowcolor{popblueL} Mot allemand & Mot français & Famille \\
\hline
\all{die Industrie} & l'industrie & industriel \\
\hline
\all{die Fabrik} & la fabrique, l'usine & fabriquer \\
\hline
\all{der Sektor} & le secteur & sectoriel \\
\hline
\all{die Produktion} & la production & produire \\
\hline
\all{exportieren} & exporter & l'exportation \\
\hline
\all{die Chemikalien} & les produits chimiques & la chimie \\
\hline
\all{die Banken} & les banques & bancaire \\
\hline
\end{tabularx}
\end{center}

\begin{attention}[Ne traduis pas mot à mot !]
L'erreur classique du francophone est de traduire chaque mot dans l'ordre. En allemand, le verbe conjugué est à la \cle{deuxième place} de la phrase principale, mais à la \cle{fin} dans une subordonnée. Cherche d'abord le \cle{sens global} de la phrase, puis consulte le glossaire seulement pour les mots vraiment inconnus. Exemple : \all{Die deutschen Produkte exportiert man in die ganze Welt.} ne se traduit pas « Les produits allemands exporte on... » mais « On exporte les produits allemands dans le monde entier. »
\end{attention}

\courssub{Le texte : Deutschlands Wirtschaft}

Lis maintenant le texte deux fois en silence. À la première lecture, cherche seulement le thème général. À la deuxième, aide-toi du glossaire.

\begin{textebox}[Deutschlands Wirtschaft]
Deutschland hat eine starke Wirtschaft. Nur wenige Menschen arbeiten in der Landwirtschaft: Die Bauern produzieren Getreide, Kartoffeln und Milch. Viele Menschen arbeiten in der Industrie: In den Fabriken produziert man Autos, Maschinen und Chemikalien. Wenn man in Deutschland arbeitet, arbeitet man meistens im Dienstleistungssektor: in Büros, Banken, Schulen, Krankenhäusern oder Geschäften. Jeden Tag fahren die Arbeiter früh zur Arbeit. Die deutschen Produkte exportiert man in die ganze Welt.
\end{textebox}

\textbf{Glossaire (mots difficiles) :}

\begin{itemize}
\item \all{die Wirtschaft} (pluriel rare) : l'économie
\item \all{die Landwirtschaft} : l'agriculture
\item \all{der Bauer}, \all{die Bauern} : le paysan, l'agriculteur
\item \all{das Getreide}, \all{die Getreide} : les céréales
\item \all{die Kartoffel}, \all{die Kartoffeln} : la pomme de terre
\item \all{die Fabrik}, \all{die Fabriken} : l'usine
\item \all{die Maschine}, \all{die Maschinen} : la machine
\item \all{das Krankenhaus}, \all{die Krankenhäuser} : l'hôpital
\item \all{das Geschäft}, \all{die Geschäfte} : le magasin, le commerce
\item \all{beschäftigen} : employer ; \all{der Arbeiter}, \all{die Arbeiter} : l'ouvrier
\end{itemize}

\courssub{Lecture détaillée : paragraphe par paragraphe}

Reprenons le texte phrase par phrase avec la méthode \emph{qui ? quoi ? où ? combien ?}

\textbf{Phrase 1 :} \all{Deutschland hat eine starke Wirtschaft.} « L'Allemagne a une économie forte. » C'est l'idée générale du texte : le verbe \all{haben} est à la deuxième place, comme en français.

\textbf{Phrase 2 :} \all{Nur wenige Menschen arbeiten in der Landwirtschaft.} « Seulement peu de gens travaillent dans l'agriculture. » \emph{Combien ?} Très peu (environ 1 \%). \emph{Quoi ?} \all{Die Bauern produzieren Getreide, Kartoffeln und Milch.} « Les paysans produisent des céréales, des pommes de terre et du lait. »

\textbf{Phrase 3 :} \all{Viele Menschen arbeiten in der Industrie.} « Beaucoup de gens travaillent dans l'industrie. » \emph{Où ?} \all{In den Fabriken produziert man Autos, Maschinen und Chemikalien.} « Dans les usines, on produit des voitures, des machines et des produits chimiques. » Remarque le pronom indéfini \all{man} = « on ».

\textbf{Phrase 4 :} \all{Wenn man in Deutschland arbeitet, arbeitet man meistens im Dienstleistungssektor.} « Quand on travaille en Allemagne, on travaille le plus souvent dans le secteur des services. » Observe bien : dans la subordonnée introduite par \all{wenn}, le verbe conjugué \all{arbeitet} va \cle{à la fin} ; la phrase principale commence alors par le verbe. C'est la règle que nous étudierons dans la partie grammaire.

\textbf{Phrases 5 et 6 :} \all{Jeden Tag fahren die Arbeiter früh zur Arbeit. Die deutschen Produkte exportiert man in die ganze Welt.} « Chaque jour, les ouvriers vont tôt au travail. On exporte les produits allemands dans le monde entier. » Remarque : quand un complément (\all{Jeden Tag}, \all{Die deutschen Produkte}) ouvre la phrase, le verbe reste à la deuxième place et le sujet passe après le verbe : c'est l'\cle{inversion}.

\begin{exemplebox}[Les trois secteurs en une phrase chacun]
\all{Der Bauer arbeitet auf dem Feld.} « Le paysan travaille dans le champ. »

\all{Der Arbeiter arbeitet in der Fabrik.} « L'ouvrier travaille dans l'usine. »

\all{Die Verkäuferin arbeitet im Geschäft.} « La vendeuse travaille dans le magasin. »
\end{exemplebox}

\courssub{Les trois secteurs en chiffres}

Le texte oppose trois mondes du travail. Voici leur poids approximatif dans l'emploi en Allemagne :

\begin{popfigure}
\begin{tikzpicture}[scale=1.1]
\fill[popgreen] (0,0) -- (90:2) arc (90:93.6:2) -- cycle;
\fill[poporange] (0,0) -- (93.6:2) arc (93.6:183.6:2) -- cycle;
\fill[popblue] (0,0) -- (183.6:2) arc (183.6:450:2) -- cycle;
\node[font=\small, text=popgreen!60!black] at (91.8:2.6) {1 \%};
\node[font=\small, text=poporange!80!black] at (138:2.6) {25 \%};
\node[font=\small, text=white] at (300:1.2) {74 \%};
\fill[popgreen] (3.2,1.4) rectangle (3.7,1.8);
\node[right, font=\small] at (3.7,1.6) {die Landwirtschaft};
\fill[poporange] (3.2,0.7) rectangle (3.7,1.1);
\node[right, font=\small] at (3.7,0.9) {die Industrie};
\fill[popblue] (3.2,0) rectangle (3.7,0.4);
\node[right, font=\small] at (3.7,0.2) {die Dienstleistungen};
\end{tikzpicture}
\legende{Die Beschäftigten in Deutschland nach Sektoren (répartition simplifiée des actifs par secteur).}
\end{popfigure}

\begin{savaistu}[Le savais-tu ?]
Moins de 2 \% des Allemands travaillent dans l'agriculture, mais l'Allemagne exporte aussi des produits agricoles ! Au Cameroun, c'est presque l'inverse : l'agriculture (cacao, café, bananes, manioc) occupe la majorité de la population active. Comparer les deux pays est un excellent sujet de production écrite : \all{In Kamerun arbeiten viele Menschen in der Landwirtschaft, in Deutschland nur wenige.}
\end{savaistu}

\courssub{Vérification de la compréhension}

\begin{exoresolu}[Questions de compréhension]
Réponds en allemand par des phrases complètes : 1) \all{Wo arbeiten nur wenige Menschen?} 2) \all{Was produziert man in den Fabriken?} 3) \all{Wo arbeitet man meistens in Deutschland?}
\tcblower
1) \all{Nur wenige Menschen arbeiten in der Landwirtschaft.} « Peu de gens travaillent dans l'agriculture. » 2) \all{In den Fabriken produziert man Autos, Maschinen und Chemikalien.} « Dans les usines, on produit des voitures, des machines et des produits chimiques. » 3) \all{Man arbeitet meistens im Dienstleistungssektor.} « On travaille le plus souvent dans le secteur des services. » Méthode : reprendre les mots de la question, garder le verbe à la deuxième place, et recopier l'information du texte.
\end{exoresolu}

\begin{exercice}[À toi de jouer : reformulation orale]
Par groupes de deux, reformulez le contenu du texte en français en trois phrases (une par secteur), puis traduisez cette phrase en allemand : « Au Cameroun, beaucoup de paysans produisent du cacao. » Indice : \all{der Kakao} = le cacao ; attention à la place du verbe !
\end{exercice}

\begin{corrige}[Reformulation]
\all{In Kamerun produzieren viele Bauern Kakao.} Le complément de lieu \all{In Kamerun} est en tête, donc le verbe \all{produzieren} reste à la deuxième place et le sujet \all{viele Bauern} vient après : c'est l'inversion, obligatoire en allemand.
\end{corrige}

\begin{aretenir}[L'essentiel de cette section]
\begin{itemize}
\item L'économie allemande repose sur trois secteurs : \all{die Landwirtschaft} (environ 1 \%), \all{die Industrie} (environ 25 \%) et \all{der Dienstleistungssektor} (environ 74 \%).
\item Pour lire un texte : titre, mots transparents, lecture globale, lecture détaillée avec glossaire, mots-clés soulignés.
\item Le verbe conjugué est toujours à la deuxième place de la phrase principale ; après \all{wenn}, il va à la fin de la subordonnée.
\end{itemize}
\end{aretenir}

\coursec{Compréhension du texte}

Dans la section précédente, nous avons lu et découvert le texte \all{Die Wirtschaft in Deutschland}, qui présente les trois grands secteurs d'activité : \all{die Landwirtschaft} (l'agriculture), \all{die Industrie} (l'industrie) et \all{der Dienstleistungssektor} (le secteur des services). Nous avons également écouté le texte et relevé les mots nouveaux. Maintenant, il faut vérifier que nous avons vraiment \emph{compris} le texte. C'est l'étape de la \cle{compréhension écrite} (\all{das Leseverstehen}) : répondre à des questions, repérer des informations précises et résumer le texte avec nos propres mots. C'est exactement l'exercice demandé à l'évaluation et au baccalauréat.

\courssub{La démarche de compréhension : du global au détail}

Comprendre un texte, c'est avancer en trois étapes, comme un escalier : d'abord le sens général, puis les détails, enfin le résumé.

\begin{popfigure}
\begin{center}
\begin{tikzpicture}[
  node distance=0.55cm and 0.9cm,
  box/.style={draw=popblue, fill=popblueL, rounded corners=3pt, thick, align=center, minimum height=0.9cm, text width=3.1cm, font=\small},
  arr/.style={-{Stealth[length=3mm]}, thick, popdark}
]
\node[box, fill=poporangeL, draw=poporange, align=center] (t){\textbf{Text}\\ \all{Die Wirtschaft in Deutschland}};
\node[box, below=of t, align=center] (g){\textbf{Fragen global}\\ \all{Worum geht es?}};
\node[box, below=of g, align=center] (d){\textbf{Fragen detailliert}\\ \all{Wer? Was? Wo? Wie viele?}};
\node[box, below=of d, fill=popgreenL, draw=popgreen, align=center] (z){\textbf{Zusammenfassung}\\ résumé en 3 phrases};
\draw[arr] (t) -- (g);
\draw[arr] (g) -- (d);
\draw[arr] (d) -- (z);
\end{tikzpicture}
\end{center}
\legende{La frise de la démarche de compréhension : du texte vers le résumé.}
\end{popfigure}

\begin{methode}[Répondre à une question de compréhension en 3 étapes]
\begin{enumerate}
\item \textbf{Repérer le mot interrogatif} : \all{wer} (qui), \all{was} (que/quoi), \all{wo} (où), \all{wann} (quand), \all{wie viele} (combien de), \all{warum} (pourquoi). Il annonce le type d'information attendu.
\item \textbf{Localiser la ligne du texte} qui contient la réponse : souligne-la au crayon.
\item \textbf{Rédiger une phrase complète en allemand} : reprends les mots de la question, mets le verbe en deuxième position, et ne recopie \emph{jamais} tout le paragraphe. Une seule phrase suffit.
\end{enumerate}
\end{methode}

\courssub{Questions de compréhension globale}

La première question, toujours, porte sur le sens général du texte.

\begin{exemplebox}[La question globale type]
\all{Worum geht es im Text?} « De quoi parle le texte ? »\par
Réponse attendue : \all{Der Text handelt von der Wirtschaft in Deutschland und den drei Sektoren: Landwirtschaft, Industrie und Dienstleistungssektor.} « Le texte traite de l'économie en Allemagne et des trois secteurs : l'agriculture, l'industrie et le secteur des services. »
\end{exemplebox}

Formules utiles pour répondre à une question globale :

\begin{center}
\begin{tabularx}{\linewidth}{|X|X|}
\hline
\rowcolor{popblueL} \textbf{Deutsch} & \textbf{Français} \\
\hline
\all{Der Text handelt von ...} & Le texte traite de ... (+ datif) \\
\hline
\all{Es geht um ...} & Il s'agit de ... (+ accusatif) \\
\hline
\all{Der Text beschreibt ...} & Le texte décrit ... (+ accusatif) \\
\hline
\all{Der Text informiert über ...} & Le texte informe sur ... (+ accusatif) \\
\hline
\end{tabularx}
\end{center}

\begin{attention}[handelt von / es geht um : attention au cas]
Après \all{handeln von}, on met le \cle{datif} : \all{Der Text handelt von der Wirtschaft.} Après \all{es geht um}, on met l'\cle{accusatif} : \all{Es geht um die Wirtschaft.} Ne confonds pas les deux constructions !
\end{attention}

\courssub{Questions de détail : Wer? Was? Wo?}

Passons maintenant aux questions de détail. Chaque mot interrogatif appelle un type de réponse précis.

\begin{center}
\begin{tabularx}{\linewidth}{|l|X|X|}
\hline
\rowcolor{popblueL} \textbf{Mot interrogatif} & \textbf{Demande} & \textbf{Réponse attendue} \\
\hline
\all{Wer?} & qui ? & une personne, un groupe \\
\hline
\all{Was?} & quoi ? & une chose, une action \\
\hline
\all{Wo?} & où ? & un lieu (avec datif : \all{in der Fabrik}) \\
\hline
\all{Wann?} & quand ? & un moment (\all{am Montag}, \all{im Sommer}) \\
\hline
\all{Wie viele?} & combien de ? & un \textbf{nombre} obligatoirement \\
\hline
\all{Warum?} & pourquoi ? & une cause avec \all{weil} + verbe final \\
\hline
\end{tabularx}
\end{center}

\begin{exemplebox}[Questions de détail sur notre texte, avec réponses modèles]
\all{Wer arbeitet in der Landwirtschaft?} « Qui travaille dans l'agriculture ? »\par
→ \all{Die Bauern arbeiten in der Landwirtschaft.} « Les agriculteurs travaillent dans l'agriculture. »\par\medskip
\all{Was produziert man in den Fabriken?} « Que produit-on dans les usines ? »\par
→ \all{In den Fabriken produziert man Autos, Maschinen und Chemikalien.} « Dans les usines, on produit des voitures, des machines et des produits chimiques. »\par\medskip
\all{Wo arbeiten die meisten Menschen?} « Où travaillent la plupart des gens ? »\par
→ \all{Die meisten Menschen arbeiten im Dienstleistungssektor.} « La plupart des gens travaillent dans le secteur des services. »
\end{exemplebox}

Observe la technique : la réponse \emph{reprend les mots de la question}. \all{Wo arbeiten die meisten Menschen?} → \all{Die meisten Menschen arbeiten im Dienstleistungssektor.} Le sujet de la question devient le sujet de la réponse, le verbe reste en deuxième position, et on ajoute seulement l'information trouvée dans le texte. C'est la méthode la plus sûre pour construire une phrase allemande correcte.

\begin{attention}[« Wie viele » et « Warum » : deux pièges classiques]
\begin{itemize}
\item \all{Wie viele} exige un \textbf{nombre} dans la réponse. Répondre \all{viele} (« beaucoup ») est insuffisant si le texte donne un chiffre : écris \all{Ungefähr zwei Prozent der Menschen arbeiten in der Landwirtschaft.}
\item \all{Warum} exige une réponse avec \all{weil}, et le verbe conjugué va \textbf{à la fin} de la proposition : \all{Warum arbeiten viele Menschen im Dienstleistungssektor? — Weil die Industrie weniger Arbeiter braucht.} (et non \all{weil die Industrie braucht weniger Arbeiter}, faute typique des francophones).
\end{itemize}
\end{attention}

\courssub{Recherche d'informations chiffrées}

Un texte économique contient souvent des nombres et des pourcentages. L'examinateur adore les questions du type \all{Wie viele Prozent ...?}. Il faut donc savoir lire et écrire les chiffres en allemand.

\begin{center}
\begin{tabularx}{\linewidth}{|l|X|X|}
\hline
\rowcolor{popblueL} \textbf{Deutsch} & \textbf{Français} & \textbf{Exemple} \\
\hline
\all{Prozent} (das, -) & pour cent & \all{zwei Prozent} = 2 \% \\
\hline
\all{ungefähr / etwa} & environ & \all{etwa 25 Prozent} \\
\hline
\all{fast} & presque & \all{fast 70 Prozent} \\
\hline
\all{die Hälfte} & la moitié & \all{die Hälfte der Menschen} \\
\hline
\all{ein Drittel / ein Viertel} & un tiers / un quart & \all{ein Viertel der Betriebe} \\
\hline
\all{die Mehrheit} & la majorité & \all{die Mehrheit arbeitet im Büro} \\
\hline
\end{tabularx}
\end{center}

\begin{attention}[La virgule et le point]
En allemand, la décimale s'écrit avec une \textbf{virgule} comme en français : \all{2,5 Prozent} se lit « deux virgule cinq pour cent » (\all{zwei Komma fünf Prozent}). En revanche, les milliers se séparent par un point : \all{1.000.000 Menschen} = 1 000 000 personnes. Ne traduis pas le point allemand par une virgule !
\end{attention}

\courssub{Vrai ou faux : justifier par une citation du texte}

L'exercice \all{Richtig oder Falsch?} est un classique. La règle d'or : une réponse sans justification ne rapporte que la moitié des points. Il faut \textbf{citer la phrase du texte} qui prouve ta réponse, entre guillemets allemands „...“.

\begin{exoresolu}[Richtig oder Falsch? — Exercice résolu]
Juge les affirmations suivantes d'après le texte \all{Die Wirtschaft in Deutschland} et justifie par une citation :\par
1. \all{In Deutschland arbeiten viele Menschen in der Landwirtschaft.}\par
2. \all{Die Industrie produziert Autos und Maschinen.}\par
3. \all{Die meisten Menschen arbeiten in Fabriken.}
\tcblower
\textbf{Solution détaillée :}\par
1. \textbf{Falsch.} Citation : „Nur ungefähr zwei Prozent der Menschen arbeiten heute in der Landwirtschaft.“ Deux pour cent, ce n'est pas « beaucoup » : l'affirmation est fausse.\par
2. \textbf{Richtig.} Citation : „In den Fabriken produziert man Autos, Maschinen und Chemikalien.“ L'affirmation reprend exactement l'information du texte.\par
3. \textbf{Falsch.} Citation : „Die meisten Menschen, etwa 70 Prozent, arbeiten im Dienstleistungssektor.“ La majorité travaille dans les services, pas dans les usines.
\end{exoresolu}

\begin{methode}[Réussir le Richtig/Falsch]
\begin{enumerate}
\item Lis l'affirmation et souligne le mot-clé (le chiffre, le lieu, la personne).
\item Retrouve la phrase correspondante dans le texte.
\item Compare mot à mot : si un seul détail change (le chiffre, le lieu), l'affirmation est \all{falsch}.
\item Écris \all{Richtig} ou \all{Falsch}, puis recopie la citation exacte entre guillemets „...“.
\end{enumerate}
\end{methode}

\courssub{Le résumé oral en trois phrases}

Dernière étape de la démarche : résumer. Au baccalauréat, la consigne type est la suivante :

\begin{textebox}[Questions types du Bac]
\all{Fassen Sie den Text in drei Sätzen zusammen.} « Résumez le texte en trois phrases. »\par\medskip
\all{Beantworten Sie die Fragen mit vollständigen Sätzen.} « Répondez aux questions par des phrases complètes. »\par\medskip
\all{Was sagen die Zahlen im Text?} « Que disent les chiffres du texte ? »
\end{textebox}

Pour construire un mini-résumé de trois phrases, on utilise trois \cle{connecteurs} simples qui organisent le temps du récit : \all{zuerst} (d'abord), \all{dann} (ensuite), \all{schließlich} (finalement). Attention : après ces connecteurs placés en tête de phrase, le verbe vient tout de suite après (position 2), puis le sujet.

\begin{exemplebox}[Modèle de résumé oral de notre texte]
\all{Zuerst beschreibt der Text die Landwirtschaft: nur wenige Bauern arbeiten noch auf dem Land.} « D'abord, le texte décrit l'agriculture : seuls quelques agriculteurs travaillent encore à la campagne. »\par
\all{Dann spricht er über die Industrie: in den Fabriken produziert man Autos und Maschinen.} « Ensuite, il parle de l'industrie : dans les usines, on produit des voitures et des machines. »\par
\all{Schließlich erklärt er, dass die meisten Menschen im Dienstleistungssektor arbeiten, zum Beispiel in Büros, Schulen und Krankenhäusern.} « Finalement, il explique que la plupart des gens travaillent dans le secteur des services, par exemple dans les bureaux, les écoles et les hôpitaux. »
\end{exemplebox}

\begin{propriete}[Les connecteurs et la place du verbe]
\all{Zuerst}, \all{dann}, \all{schließlich} occupent la \cle{première position} de la phrase : le verbe conjugué reste en \textbf{deuxième position} et le sujet suit le verbe (inversion).\par
\all{Zuerst beschreibt der Text die Landwirtschaft.} (et non \all{Zuerst der Text beschreibt...}).\par
Exception utile : dans une subordonnée avec \all{dass}, le verbe va à la fin : \all{..., dass die meisten Menschen im Dienstleistungssektor arbeiten.}
\end{propriete}

\courssub{Entraînement guidé}

\begin{exercice}[À toi de jouer !]
Réponds par des phrases complètes, puis fais le vrai/faux :\par
1. \all{Worum geht es im Text „Die Wirtschaft in Deutschland“?}\par
2. \all{Wie viele Prozent der Menschen arbeiten in der Landwirtschaft?}\par
3. \all{Wo arbeitet ein Bauer?}\par
4. \all{Richtig oder Falsch? „Der Dienstleistungssektor ist der größte Sektor.“} Justifie par une citation.\par
5. Résume le texte oralement en trois phrases avec \all{zuerst}, \all{dann}, \all{schließlich}.
\end{exercice}

\begin{corrige}[Exercice « À toi de jouer ! »]
1. \all{Der Text handelt von den drei Wirtschaftssektoren in Deutschland.}\par
2. \all{Ungefähr zwei Prozent der Menschen arbeiten in der Landwirtschaft.}\par
3. \all{Ein Bauer arbeitet auf dem Land, auf seinem Hof.}\par
4. \textbf{Richtig.} Citation : „Die meisten Menschen, etwa 70 Prozent, arbeiten im Dienstleistungssektor.“\par
5. Exemple : \all{Zuerst spricht der Text über die Landwirtschaft. Dann beschreibt er die Industrie mit den Fabriken. Schließlich erklärt er, dass die meisten Menschen im Dienstleistungssektor arbeiten.}
\end{corrige}

\begin{aretenir}[L'essentiel de la compréhension de texte]
\begin{itemize}
\item Toujours commencer par la question globale : \all{Worum geht es im Text?} → \all{Der Text handelt von ...}
\item Le mot interrogatif commande la réponse : \all{wer} → personne, \all{wo} → lieu, \all{wie viele} → nombre, \all{warum} → \all{weil} + verbe final.
\item Réponse = phrase complète qui reprend les mots de la question, verbe en position 2.
\item Vrai/faux : toujours justifier par une citation exacte entre guillemets „...“.
\item Résumé en trois phrases : \all{zuerst} — \all{dann} — \all{schließlich}, avec inversion du sujet après le connecteur.
\end{itemize}
\end{aretenir}

\coursec{Lexique thématique : les secteurs d'activité}

Après avoir lu et compris le texte \all{Die Wirtschaft in Deutschland}, nous avons repéré de nombreux mots qui tournent autour du monde du travail : \all{die Landwirtschaft}, \all{die Industrie}, \all{der Dienstleistungssektor}, \all{der Bauer}, \all{die Fabrik}, \all{der Beruf}. Cette troisième partie a pour but de fixer ce vocabulaire : nous allons le classer en trois grands champs lexicaux, apprendre les articles et les pluriels par cœur, découvrir comment l'allemand fabrique les familles de mots et les féminins de métiers, puis nous entraîner à employer ces mots dans des phrases complètes. À la fin de la séance, vous devrez être capables de présenter oralement le métier d'un membre de votre famille en trois ou quatre phrases.

\courssub{Observation : un court texte pour faire émerger le vocabulaire}

Avant d'apprendre des listes de mots, observons-les d'abord en situation. Lis le texte suivant deux fois : une première fois pour le sens général, une deuxième fois en soulignant tous les mots qui désignent un métier, un lieu de travail ou une activité économique.

\begin{textebox}[Drei Berufe in meiner Familie]
Meine Familie ist groß, und wir arbeiten in drei verschiedenen Sektoren. Mein Großvater ist Bauer. Er hat einen Bauernhof in der Nähe von Bafoussam. Auf seinen Feldern baut er Mais und Gemüse an. Im September ist die Ernte : die ganze Familie hilft mit, denn es gibt viel Arbeit. Mein Vater ist Arbeiter in einer Fabrik in Douala. Die Fabrik produziert Schokolade aus Kakao. Er arbeitet jeden Tag von sieben Uhr bis fünf Uhr nachmittags. Die Maschinen sind sehr laut, aber er mag seinen Beruf. Die Produkte der Fabrik exportiert die Firma auch nach Europa. Meine Schwester ist Angestellte in einer Bank in Yaoundé. Sie arbeitet in einem Büro im Zentrum der Stadt. Jeden Tag bedient sie viele Kunden : sie öffnet Konten, gibt Informationen und verkauft auch Produkte der Bank. Sie verdient gutes Geld und ist mit ihrer Arbeit zufrieden. Am Wochenende sprechen wir oft über unsere Berufe. Mein Großvater sagt : „Die Landwirtschaft ist die wichtigste Arbeit, denn ohne Essen kann niemand leben.“ Mein Vater antwortet : „Aber ohne Industrie gibt es keine Waren und keine Arbeitsplätze !“ Und meine Schwester lacht : „Und ohne Banken gibt es kein Geld für euch beide !“ Ich finde : alle drei Sektoren sind wichtig für die Wirtschaft unseres Landes. Und ich ? Ich möchte später Lehrerin werden und in einer Schule arbeiten.
\end{textebox}

\textbf{Glossaire} : \all{verschieden} : différent ; \all{der Mais} : le maïs ; \all{das Gemüse} : les légumes ; \all{laut} : bruyant ; \all{die Firma} : l'entreprise ; \all{das Konto (Konten)} : le compte ; \all{zufrieden} : satisfait ; \all{ohne} : sans ; \all{niemand} : personne, ne... personne ; \all{später} : plus tard ; \all{helfen mit} : aider, participer.

\textbf{Questions rapides} (réponds en allemand, phrase complète) :
\begin{enumerate}
\item \all{Wo arbeitet der Großvater ?}
\item \all{Was produziert die Fabrik in Douala ?}
\item \all{Wo arbeitet die Schwester und was macht sie jeden Tag ?}
\item \all{Was möchte die Erzählerin später werden ?}
\end{enumerate}

\courssub{Le concept clé : \cle{der Sektor}}

\begin{definition}[der Sektor]
\all{Der Sektor} (pluriel : \all{die Sektoren}) désigne un grand domaine de l'économie. Les économistes distinguent trois secteurs d'activité : \all{der primäre Sektor} (le secteur primaire : l'agriculture et l'élevage), \all{der sekundäre Sektor} (le secteur secondaire : l'industrie et la transformation) et \all{der tertiäre Sektor} (le secteur tertiaire : les services, le commerce, l'éducation, la santé). Ces trois mots sont très utiles pour le commentaire de texte au Bac.
\end{definition}

\begin{exemplebox}[Les trois secteurs en une phrase]
\all{Der Bauer arbeitet im primären Sektor, der Arbeiter im sekundären Sektor und der Angestellte im tertiären Sektor.} « Le paysan travaille dans le secteur primaire, l'ouvrier dans le secteur secondaire et l'employé dans le secteur tertiaire. »
\end{exemplebox}

\courssub{Champ lexical 1 : \cle{die Landwirtschaft} (l'agriculture)}

Voici les mots essentiels du premier secteur. Apprends chaque nom avec son article et son pluriel : en allemand, c'est indissociable.

\begin{center}
\begin{tabularx}{\linewidth}{|l|l|X|}\hline
\rowcolor{popblueL} Deutsch & Pluriel & Français \\ \hline
die Landwirtschaft & (rare) & l'agriculture \\ \hline
der Bauer & die Bauern & le paysan, l'agriculteur \\ \hline
die Bäuerin & die Bäuerinnen & la paysanne, l'agricultrice \\ \hline
das Feld & die Felder & le champ \\ \hline
der Bauernhof & die Bauernhöfe & la ferme \\ \hline
das Getreide & (rare) & les céréales \\ \hline
die Ernte & die Ernten & la récolte \\ \hline
anbauen & -- & cultiver (verbe à particule séparable) \\ \hline
ernten & -- & récolter \\ \hline
\end{tabularx}
\end{center}

\begin{exemplebox}[Exemples traduits]
\all{Der Bauer arbeitet auf dem Feld.} « Le paysan travaille au champ. »\par
\all{Die Bäuerin baut Getreide an.} « La paysanne cultive des céréales. » (attention : \all{anbauen} se sépare : \all{baut ... an})\par
\all{Im Herbst ernten die Bauern den Mais.} « En automne, les paysans récoltent le maïs. »\par
\all{Auf dem Bauernhof gibt es viele Tiere.} « À la ferme, il y a beaucoup d'animaux. »
\end{exemplebox}

\begin{attention}[der Bauer : un nom faible (N-Deklination) !]
\all{Der Bauer} a un pluriel irrégulier : \all{die Bauern} (et non *\all{Bauer}). C'est un \textbf{nom faible} : il prend un \textbf{-n} à tous les cas obliques du singulier (génitif, datif, accusatif) : \all{des Bauern} (génitif), \all{dem Bauern} (datif), \all{den Bauern} (accusatif). Seul le nominatif singulier reste \all{der Bauer}. Même phénomène avec \all{der Kunde} : \all{des Kunden}, \all{dem Kunden}, \all{den Kunden}, \all{die Kunden} (pluriel).
\end{attention}

\courssub{Champ lexical 2 : \cle{die Industrie} (l'industrie)}

\begin{center}
\begin{tabularx}{\linewidth}{|l|l|X|}\hline
\rowcolor{popblueL} Deutsch & Pluriel & Français \\ \hline
die Industrie & die Industrien & l'industrie \\ \hline
die Fabrik & die Fabriken & l'usine \\ \hline
der Arbeiter & die Arbeiter & l'ouvrier \\ \hline
die Arbeiterin & die Arbeiterinnen & l'ouvrière \\ \hline
die Maschine & die Maschinen & la machine \\ \hline
das Produkt & die Produkte & le produit \\ \hline
die Ware & die Waren & la marchandise \\ \hline
produzieren & -- & produire, fabriquer \\ \hline
exportieren & -- & exporter \\ \hline
\end{tabularx}
\end{center}

\begin{exemplebox}[Exemples traduits]
\all{Die Fabrik produziert Autos.} « L'usine produit des voitures. »\par
\all{Der Arbeiter repariert die Maschine.} « L'ouvrier répare la machine. »\par
\all{Deutschland exportiert viele Waren nach Afrika.} « L'Allemagne exporte beaucoup de marchandises vers l'Afrique. »\par
\all{Die Arbeiterin kontrolliert die Produkte.} « L'ouvrière contrôle les produits. »
\end{exemplebox}

\courssub{Champ lexical 3 : \cle{der Dienstleistungssektor} (les services)}

\begin{center}
\begin{tabularx}{\linewidth}{|l|l|X|}\hline
\rowcolor{popblueL} Deutsch & Pluriel & Français \\ \hline
der Dienstleistungssektor & die -sektoren & le secteur des services \\ \hline
der Beruf & die Berufe & le métier, la profession \\ \hline
der Angestellte & die Angestellten & l'employé \\ \hline
die Angestellte & die Angestellten & l'employée \\ \hline
das Büro & die Büros & le bureau \\ \hline
die Bank & die Banken & la banque \\ \hline
das Geschäft & die Geschäfte & le magasin, le commerce \\ \hline
die Schule & die Schulen & l'école \\ \hline
das Krankenhaus & die Krankenhäuser & l'hôpital \\ \hline
der Kunde & die Kunden & le client \\ \hline
die Kundin & die Kundinnen & la cliente \\ \hline
bedienen & -- & servir (un client) \\ \hline
verkaufen & -- & vendre \\ \hline
\end{tabularx}
\end{center}

\begin{exemplebox}[Exemples traduits]
\all{Sie verdient Geld in einer Bank.} « Elle gagne de l'argent dans une banque. »\par
\all{Die Verkäuferin bedient die Kunden im Geschäft.} « La vendeuse sert les clients dans le magasin. »\par
\all{Der Arzt arbeitet im Krankenhaus.} « Le médecin travaille à l'hôpital. »\par
\all{Was ist dein Beruf ? — Ich bin Lehrerin an einer Schule.} « Quel est ton métier ? — Je suis professeure dans une école. »
\end{exemplebox}

\begin{savaistu}[L'allemand, champion des mots composés]
\all{Die Dienstleistung} (la prestation de service) se décompose en deux mots : \all{der Dienst} (le service) + \all{die Leistung} (la prestation, la performance). L'allemand adore coller les mots entre eux ! Autre exemple dans notre champ lexical : \all{das Krankenhaus} = \all{krank} (malade) + \all{das Haus} (la maison), littéralement « la maison des malades ». Et \all{der Bauernhof} = \all{der Bauer} + \all{der Hof} (la cour) : « la cour du paysan », c'est-à-dire la ferme. Quand tu rencontres un long mot allemand, découpe-le : tu comprendras souvent son sens tout seul !
\end{savaistu}

\courssub{Mots transversaux : le vocabulaire du travail en général}

Ces mots s'emploient dans les trois secteurs :

\begin{center}
\begin{tabularx}{\linewidth}{|l|l|X|}\hline
\rowcolor{popblueL} Deutsch & Pluriel & Français \\ \hline
die Wirtschaft & die Wirtschaften & l'économie \\ \hline
der Sektor & die Sektoren & le secteur \\ \hline
die Arbeit & die Arbeiten & le travail \\ \hline
der Betrieb & die Betriebe & l'entreprise, l'exploitation \\ \hline
arbeiten & -- & travailler \\ \hline
verdienen & -- & gagner (de l'argent) \\ \hline
\end{tabularx}
\end{center}

L'expression \all{Geld verdienen} (gagner de l'argent) est indispensable : \all{Mein Bruder verdient sein Geld als Taxifahrer.} « Mon frère gagne sa vie comme chauffeur de taxi. » On dit aussi \all{bei einer Firma arbeiten} (travailler dans une entreprise) : \all{Sie arbeitet bei einer deutschen Firma.} « Elle travaille dans une entreprise allemande. »

\courssub{Les familles de mots : un radical, plusieurs mots}

L'allemand construit des familles entières à partir d'un verbe ou d'un radical. Les connaître multiplie ton vocabulaire sans effort.

\begin{center}
\begin{tabularx}{\linewidth}{|X|X|X|}\hline
\rowcolor{popblueL} Verbe & Nom d'action & Personne / lieu \\ \hline
arbeiten (travailler) & die Arbeit (le travail) & der Arbeiter / die Arbeiterin (l'ouvrier / l'ouvrière) \\ \hline
arbeiten & der Arbeitsplatz (le poste de travail) & der Arbeitgeber (l'employeur) \\ \hline
produzieren (produire) & die Produktion (la production) & das Produkt (le produit) \\ \hline
verkaufen (vendre) & der Verkauf (la vente) & der Verkäufer / die Verkäuferin (le vendeur / la vendeuse) \\ \hline
\end{tabularx}
\end{center}

\begin{exemplebox}[Une famille en action]
\all{Die Arbeiter arbeiten in der Fabrik. Ihre Arbeit ist schwer, aber die Produktion ist wichtig für die Wirtschaft.} « Les ouvriers travaillent dans l'usine. Leur travail est dur, mais la production est importante pour l'économie. »
\end{exemplebox}

\courssub{Le féminin des noms de métiers : le suffixe \cle{-in}}

Observons d'abord les paires du texte : \all{der Bauer} / \all{die Bäuerin}, \all{der Arbeiter} / \all{die Arbeiterin}, \all{der Lehrer} / \all{die Lehrerin}. Que remarques-tu ? Le féminin ajoute toujours la même syllabe.

\begin{propriete}[Formation du féminin des métiers]
Le féminin des noms de métiers et de personnes se forme en ajoutant le suffixe \textbf{-in} au masculin : \all{der Lehrer} $\rightarrow$ \all{die Lehrerin}, \all{der Kunde} $\rightarrow$ \all{die Kundin}, \all{der Verkäufer} $\rightarrow$ \all{die Verkäuferin}. L'article devient bien sûr \all{die}. Souvent, la voyelle du radical prend un \textbf{Umlaut} : \all{der Bauer} $\rightarrow$ \all{die Bäuerin}, \all{der Arzt} $\rightarrow$ \all{die Ärztin}. Le pluriel de ces féminins est toujours \textbf{-innen} : \all{die Bäuerinnen}, \all{die Lehrerinnen}, \all{die Kundinnen}.
\end{propriete}

\begin{center}
\begin{tabularx}{\linewidth}{|X|X|X|}\hline
\rowcolor{popblueL} Masculin & Féminin & Pluriel féminin \\ \hline
der Bauer (le paysan) & die Bäuerin & die Bäuerinnen \\ \hline
der Arbeiter (l'ouvrier) & die Arbeiterin & die Arbeiterinnen \\ \hline
der Lehrer (le professeur) & die Lehrerin & die Lehrerinnen \\ \hline
der Kunde (le client) & die Kundin & die Kundinnen \\ \hline
der Verkäufer (le vendeur) & die Verkäuferin & die Verkäuferinnen \\ \hline
\end{tabularx}
\end{center}

\begin{attention}[Pièges pour les francophones]
\begin{itemize}
\item Ne dis jamais *\all{die Bauerin} : il faut l'Umlaut, \all{die Bäuerin}.
\item Au pluriel féminin, c'est \textbf{-innen} avec double n : \all{die Lehrerinnen}, et non *\all{die Lehrerinen}.
\item Attention au faux ami : \all{die Fabrik} signifie « l'usine », pas « la fabrique » au sens abstrait ; et \all{der Betrieb} signifie « l'entreprise », pas « l'exploitation » au sens agricole uniquement.
\item \all{der Angestellte} se décline comme un adjectif (déclinaison faible) : \all{ein Angestellter} (un employé), \all{die Angestellten} (les employés). Au féminin, la forme est identique : \all{die Angestellte}.
\end{itemize}
\end{attention}

\courssub{Carte mentale : visualiser les trois secteurs}

\begin{popfigure}
\begin{center}
\begin{tikzpicture}[
  mindmap,
  grow cyclic,
  every node/.style={concept},
  root concept/.append style={concept color=popblue, font=\large\bfseries},
  level 1/.append style={level distance=3.6cm, sibling angle=120},
  level 2/.append style={level distance=2.6cm, sibling angle=45},
  text=white]
\node[root concept]{die Wirtschaft}
  child[concept color=popgreen]{ node {die Landwirtschaft}
    child { node[font=\small] {der Bauer} }
    child { node[font=\small] {das Feld} }
    child { node[font=\small] {die Ernte} }
  }
  child[concept color=poporange]{ node {die Industrie}
    child { node[font=\small] {die Fabrik} }
    child { node[font=\small] {die Maschine} }
    child { node[font=\small] {das Produkt} }
  }
  child[concept color=poppurple]{ node {die Dienst\-leistung}
    child { node[font=\small] {das Büro} }
    child { node[font=\small] {die Bank} }
    child { node[font=\small] {der Kunde} }
  };
\end{tikzpicture}
\end{center}
\legende{Carte mentale du vocabulaire : l'économie et ses trois secteurs d'activité.}
\end{popfigure}

Reproduis cette carte mentale dans ton cahier et ajoute deux mots supplémentaires par branche : c'est une excellente méthode de mémorisation.

\courssub{Entraînement}

\begin{exoresolu}[Complète avec le mot qui convient]
Énoncé — Complète les phrases avec : \all{Bäuerin}, \all{Fabrik}, \all{Kunden}, \all{Ernte}, \all{Büro}.
\begin{enumerate}
\item \all{Die \ldots\ arbeitet auf dem Bauernhof und baut Gemüse an.}
\item \all{In der \ldots\ produzieren die Arbeiter Schokolade.}
\item \all{Die Angestellte bedient jeden Tag viele \ldots\ in der Bank.}
\item \all{Im September ist die \ldots\ : die Bauern ernten den Mais.}
\item \all{Mein Onkel arbeitet in einem \ldots\ im Zentrum von Yaoundé.}
\end{enumerate}
\tcblower
Solution détaillée :
\begin{enumerate}
\item \all{Die \textbf{Bäuerin} arbeitet auf dem Bauernhof und baut Gemüse an.} (féminin de \all{der Bauer}, avec Umlaut : la phrase parle d'une femme, article \all{die}.)
\item \all{In der \textbf{Fabrik} produzieren die Arbeiter Schokolade.} (les ouvriers travaillent dans une usine ; datif : \all{in der Fabrik}.)
\item \all{Die Angestellte bedient jeden Tag viele \textbf{Kunden} in der Bank.} (pluriel de \all{der Kunde} : \all{die Kunden}, accusatif après \all{bedienen}.)
\item \all{Im September ist die \textbf{Ernte} : die Bauern ernten den Mais.} (le verbe \all{ernten} donne le nom \all{die Ernte}.)
\item \all{Mein Onkel arbeitet in einem \textbf{Büro} im Zentrum von Yaoundé.} (datif neutre : \all{in einem Büro}.)
\end{enumerate}
\end{exoresolu}

\begin{exercice}[À toi de jouer]
\begin{enumerate}
\item Traduis en allemand : a) Le paysan cultive des céréales dans son champ. b) L'usine exporte ses produits en Europe. c) Ma sœur est employée dans une banque et gagne bien sa vie.
\item Forme le féminin et son pluriel : \all{der Lehrer}, \all{der Verkäufer}, \all{der Arbeiter}, \all{der Kunde}.
\item Classe ces mots dans les trois secteurs : \all{das Krankenhaus}, \all{die Maschine}, \all{das Feld}, \all{die Bank}, \all{das Getreide}, \all{die Ware}.
\item Production orale : présente en quatre phrases le métier d'un membre de ta famille (secteur, lieu de travail, activité principale).
\end{enumerate}
\end{exercice}

\begin{corrige}[Exercice : À toi de jouer]
\begin{enumerate}
\item a) \all{Der Bauer baut Getreide auf seinem Feld an.} (verbe séparable : \all{an} en fin de phrase ; \all{seinem} datif neutre après \all{auf} indiquant la position.) b) \all{Die Fabrik exportiert ihre Produkte nach Europa.} (\all{nach} + datif pour la direction vers un pays/continent.) c) \all{Meine Schwester ist Angestellte in einer Bank und verdient gutes Geld.} (pas d'article devant le métier après \all{sein} ; \all{gutes Geld} accusatif.)
\item \all{der Lehrer} $\rightarrow$ \all{die Lehrerin}, \all{die Lehrerinnen} ; \all{der Verkäufer} $\rightarrow$ \all{die Verkäuferin}, \all{die Verkäuferinnen} ; \all{der Arbeiter} $\rightarrow$ \all{die Arbeiterin}, \all{die Arbeiterinnen} ; \all{der Kunde} $\rightarrow$ \all{die Kundin}, \all{die Kundinnen}.
\item Landwirtschaft : \all{das Feld}, \all{das Getreide} ; Industrie : \all{die Maschine}, \all{die Ware} ; Dienstleistung : \all{das Krankenhaus}, \all{die Bank}.
\item Modèle : \all{Mein Onkel ist Arbeiter. Er arbeitet in einer Fabrik in Douala. Die Fabrik produziert Seife. Er arbeitet von Montag bis Samstag und verdient gutes Geld.}
\end{enumerate}
\end{corrige}

\begin{aretenir}[L'essentiel du lexique]
\begin{itemize}
\item L'économie se divise en trois secteurs : \all{der primäre Sektor} (Landwirtschaft), \all{der sekundäre Sektor} (Industrie), \all{der tertiäre Sektor} (Dienstleistungen).
\item Chaque nom s'apprend avec article et pluriel : \all{der Bauer, die Bauern} ; \all{die Fabrik, die Fabriken} ; \all{das Geschäft, die Geschäfte} ; \all{das Krankenhaus, die Krankenhäuser}.
\item Le féminin des métiers se forme avec \textbf{-in} (souvent avec Umlaut) et le pluriel féminin avec \textbf{-innen} : \all{der Bauer} $\rightarrow$ \all{die Bäuerin} $\rightarrow$ \all{die Bäuerinnen}.
\item \all{der Bauer} et \all{der Kunde} sont des noms faibles (N-Deklination) : \all{den Bauern}, \all{dem Bauern}, \all{des Bauern} / \all{den Kunden}, \all{dem Kunden}, \all{des Kunden}.
\item Expressions à réutiliser : \all{Geld verdienen}, \all{bei einer Firma arbeiten}, \all{Was ist dein Beruf ?}, \all{arbeiten als} + métier (sans article).
\end{itemize}
\end{aretenir}

Maintenant que le vocabulaire est en place, passons à la grammaire : nous allons apprendre à situer ces activités dans le temps grâce aux compléments de temps.

\coursec{Grammaire 1 : les compléments de temps}

Dans la section précédente, nous avons enrichi notre vocabulaire : nous savons maintenant nommer les trois secteurs d'activité (\all{die Landwirtschaft}, \all{die Industrie}, \all{der Dienstleistungssektor}) et les métiers qui s'y rattachent (\all{der Bauer}, \all{der Arbeiter}, \all{der Verkäufer}). Mais pour raconter la vie économique — quand les ouvriers partent au travail, depuis combien de temps une employée travaille à la banque, combien de fois par semaine le marché a lieu —, il faut savoir situer les actions dans le temps. C'est le rôle des \cle{compléments de temps}, en allemand \all{die Zeitangaben}.

\courssub{Observation : repérer les compléments de temps dans le texte}

Reprenons des phrases inspirées de notre texte support \all{Die Wirtschaft in Deutschland} et soulignons ce qui indique le temps :

\begin{exemplebox}[Observation]
\all{Jeden Tag fahren die Arbeiter früh zur Arbeit.} « Chaque jour, les ouvriers vont tôt au travail. »

\all{Um 7 Uhr beginnt die Arbeit in der Fabrik.} « À 7 heures, le travail commence à l'usine. »

\all{Im Sommer arbeiten viele Bauern auf dem Feld.} « En été, beaucoup de paysans travaillent aux champs. »

\all{Seit einem Jahr arbeitet sie in einer Bank.} « Elle travaille dans une banque depuis un an. »

\all{Am Montag findet der Markt in Douala statt.} « Lundi, le marché a lieu à Douala. »
\end{exemplebox}

Observons : les groupes de mots \all{jeden Tag}, \all{um 7 Uhr}, \all{im Sommer}, \all{seit einem Jahr}, \all{am Montag} ne sont ni le sujet ni le verbe. Ils répondent à des questions comme « quand ? », « depuis combien de temps ? », « combien de fois ? ». Ce sont des \cle{compléments de temps}.

\begin{definition}[Le complément de temps]
Le \cle{complément de temps} (\all{die Zeitangabe, -n}) est un mot ou un groupe de mots qui situe l'action dans le temps. Il répond à trois questions :
\begin{itemize}
\item \all{Wann?} — quand ? \all{Am Montag, heute, um 8 Uhr.}
\item \all{Wie lange?} — combien de temps ? \all{Seit einem Jahr, den ganzen Tag.}
\item \all{Wie oft?} — combien de fois ? \all{Jeden Tag, oft, manchmal, zweimal pro Woche.}
\end{itemize}
\end{definition}

\courssub{Les formes des compléments de temps}

Les compléments de temps prennent trois formes principales.

\textbf{1. Les adverbes de temps}\par
Ce sont des mots invariables : \all{heute} (aujourd'hui), \all{jetzt} (maintenant), \all{morgen} (demain), \all{gestern} (hier), \all{oft} (souvent), \all{immer} (toujours), \all{manchmal} (parfois), \all{früh} (tôt), \all{spät} (tard).

\all{Die Arbeiter kommen heute früh.} « Les ouvriers arrivent tôt aujourd'hui. »

\textbf{2. Les groupes nominaux à l'accusatif (sans préposition)}\par
Pour exprimer la répétition, la durée ou un moment précis, l'allemand utilise l'accusatif, sans préposition : \all{jeden Tag} (chaque jour), \all{jede Woche} (chaque semaine), \all{jedes Jahr} (chaque année), \all{letztes Jahr} (l'année dernière), \all{nächsten Monat} (le mois prochain), \all{den ganzen Tag} (toute la journée).

\all{Jede Woche verkauft der Bauer Kakao auf dem Markt.} « Chaque semaine, le paysan vend du cacao au marché. »

\textbf{3. Les groupes prépositionnels}\par
C'est la forme la plus fréquente : une préposition temporelle + un nom.

\begin{propriete}[Les prépositions temporelles essentielles]
\begin{itemize}
\item \all{am} + jour ou date : \all{am Montag} (le lundi), \all{am 1. Mai} (le 1\up{er} mai), \all{am Wochenende} (le week-end).
\item \all{im} + mois ou saison : \all{im Januar} (en janvier), \all{im Sommer} (en été), \all{im Winter} (en hiver).
\item \all{um} + heure : \all{um 8 Uhr} (à 8 heures), \all{um 17.30 Uhr} (à 17 h 30).
\item \all{seit} + point de départ (datif) : \all{seit einem Jahr} (depuis un an), \all{seit Montag} (depuis lundi).
\item \all{von ... bis} (de ... à) : \all{von Montag bis Freitag} (du lundi au vendredi), \all{von 8 bis 16 Uhr} (de 8 h à 16 h).
\end{itemize}
\end{propriete}

\begin{attention}[Trois prépositions à mémoriser : la logique n'est pas celle du français]
Le français dit « le lundi, en été, à 8 heures » avec trois constructions différentes ; l'allemand aussi, mais ce ne sont pas les mêmes découpages ! Il faut mémoriser par cœur : \all{am} + jour, \all{im} + mois/saison, \all{um} + heure. Ne dites jamais *\all{in Montag} ou *\all{am 8 Uhr}.

Autre piège : \all{seit} se construit avec le \textbf{présent} allemand là où le français emploie le présent aussi, mais jamais avec le passé composé : \all{Sie arbeitet seit einem Jahr hier.} « Elle travaille ici depuis un an » (et non *\all{Sie hat seit einem Jahr gearbeitet}).
\end{attention}

\begin{center}
\begin{tabularx}{\linewidth}{|X|X|X|}
\hline
\rowcolor{popblueL} Français & Deutsch & Remarque \\
\hline
le lundi & am Montag & \all{am} + jour \\
\hline
en été & im Sommer & \all{im} + saison \\
\hline
en janvier & im Januar & \all{im} + mois \\
\hline
à 8 heures & um 8 Uhr & \all{um} + heure \\
\hline
depuis un an & seit einem Jahr & \all{seit} + datif \\
\hline
chaque jour & jeden Tag & accusatif sans préposition \\
\hline
du lundi au vendredi & von Montag bis Freitag & \all{von ... bis} \\
\hline
\end{tabularx}
\end{center}

\courssub{La place du complément de temps et l'inversion}

Deux positions sont possibles pour le complément de temps.

\textbf{Position 1 : après le verbe conjugué.}\par
\all{Der Bauer arbeitet am Montag auf dem Feld.} « Le paysan travaille le lundi aux champs. »

\textbf{Position 2 : en tête de phrase.}\par
C'est très fréquent en allemand, pour insister sur le moment. Mais attention : le verbe conjugué reste \textbf{toujours en deuxième position}, donc le sujet passe après le verbe. C'est la fameuse \cle{inversion} :

\all{Am Montag arbeitet der Bauer auf dem Feld.} « Le lundi, le paysan travaille aux champs. »

\begin{popfigure}
\begin{tikzpicture}[font=\small]
\node[draw=popblue, fill=popblueL, rounded corners, align=center] (t) at (0,0) {\textbf{Am Montag}\\ \footnotesize complément de temps};
\node[draw=poporange, fill=poporangeL, rounded corners, align=center] (v) at (3.4,0) {\textbf{arbeitet}\\ \footnotesize verbe (2\up{e} position)};
\node[draw=popgreen, fill=popgreenL, rounded corners, align=center] (s) at (6.6,0) {\textbf{der Bauer}\\ \footnotesize sujet};
\node[draw=poppurple, fill=poppurpleL, rounded corners, align=center] (l) at (9.8,0) {\textbf{auf dem Feld}\\ \footnotesize lieu};
\draw[-{Stealth}, thick, popdark] (t) -- (v);
\draw[-{Stealth}, thick, popdark] (v) -- (s);
\draw[-{Stealth}, thick, popdark] (s) -- (l);
\draw[-{Stealth}, thick, poppink, bend right=35] (s.south) to node[below, align=center, font=\footnotesize]{inversion sujet--verbe} (v.south);
\end{tikzpicture}
\legende{Quand le complément de temps ouvre la phrase, le verbe garde la 2\up{e} position et le sujet se place après lui.}
\end{popfigure}

\begin{attention}[L'erreur numéro 1 des francophones]
En français, on dit « Chaque jour, il va à l'usine » : le sujet reste avant le verbe. En allemand, c'est impossible ! Après un complément de temps en tête, le verbe vient immédiatement, puis le sujet :

\all{Jeden Tag fährt er zur Fabrik.} « Chaque jour, il va à l'usine. »

Jamais *\all{Jeden Tag er fährt zur Fabrik.} — le verbe doit rester en deuxième position, quoi qu'il y ait en premier.
\end{attention}

\begin{methode}[L'ordre TeKaMoLo]
Quand une phrase contient plusieurs compléments, l'allemand les range dans un ordre fixe. Retenez le mot-valise \cle{TeKaMoLo} :
\begin{enumerate}
\item \textbf{Te}mporal : quand ? \all{am Montag}
\item \textbf{Ka}usal : pourquoi ? \all{wegen der Arbeit}
\item \textbf{Mo}dal : comment ? \all{mit dem Bus}
\item \textbf{Lo}kal : où ? \all{zur Fabrik}
\end{enumerate}
Exemple : \all{Die Arbeiter fahren am Montag (Te) mit dem Bus (Mo) zur Fabrik (Lo).} « Le lundi, les ouvriers vont à l'usine en bus. »
\end{methode}

\courssub{Texte d'application : Ein Arbeitstag in Yaoundé}

\begin{textebox}[Ein Arbeitstag in Yaoundé]
Frau Mbarga arbeitet seit fünf Jahren in einer Bank in Yaoundé. Jeden Tag steht sie um 5.30 Uhr auf, denn der Verkehr in der Stadt ist am Morgen sehr stark. Um 6.15 Uhr fährt sie mit dem Taxi zur Arbeit. Von Montag bis Freitag beginnt ihre Arbeit um 8 Uhr und endet um 16 Uhr. Am Mittwoch hat sie oft eine Besprechung mit ihren Kollegen. In der Mittagspause, von 12 bis 13 Uhr, isst sie mit einer Freundin in einem kleinen Restaurant. Am Wochenende arbeitet sie nicht : am Samstag geht sie manchmal auf den Markt, und am Sonntag besucht sie ihre Familie. Im Juli macht sie immer Ferien. Letztes Jahr ist sie nach Deutschland geflogen und hat eine Fabrik in München besucht. Nächstes Jahr möchte sie nach Österreich reisen. „Meine Arbeit ist interessant“, sagt sie, „aber die Wochenenden sind die schönsten Tage der Woche !“
\end{textebox}

\textbf{Glossaire :} \all{die Bank, -en} (la banque) ; \all{aufstehen (steht auf, stand auf, ist aufgestanden)} (se lever) ; \all{der Verkehr} (la circulation) ; \all{die Besprechung, -en} (la réunion) ; \all{die Mittagspause, -n} (la pause de midi) ; \all{besuchen} (rendre visite à, visiter) ; \all{Ferien machen} (prendre des vacances) ; \all{reisen} (voyager).

\textbf{Questions de compréhension :}
\begin{enumerate}
\item \all{Seit wann arbeitet Frau Mbarga in der Bank?}
\item \all{Wann beginnt ihre Arbeit am Morgen?}
\item \all{Was macht sie am Samstag?}
\item \all{Wohin ist sie letztes Jahr geflogen?}
\end{enumerate}

\textbf{Réponses attendues :} 1. \all{Sie arbeitet seit fünf Jahren in der Bank.} 2. \all{Ihre Arbeit beginnt um 8 Uhr.} 3. \all{Am Samstag geht sie manchmal auf den Markt.} 4. \all{Letztes Jahr ist sie nach Deutschland geflogen.}

\courssub{Entraînement}

\begin{exoresolu}[Placer le complément de temps en tête]
\textbf{Énoncé.} Transformez les phrases en plaçant le complément de temps souligné en tête de phrase. Attention à l'inversion !

1. \all{Die Arbeiter fahren \underline{jeden Tag} früh zur Arbeit.}

2. \all{Der Bauer verkauft \underline{am Montag} Kakao auf dem Markt.}

3. \all{Wir besuchen \underline{im Sommer} eine Fabrik.}

\tcblower
\textbf{Solution détaillée.} Étape 1 : placer le complément de temps en premier. Étape 2 : mettre le verbe conjugué juste après (2\up{e} position). Étape 3 : faire suivre le sujet.

1. \all{Jeden Tag fahren die Arbeiter früh zur Arbeit.} « Chaque jour, les ouvriers vont tôt au travail. »

2. \all{Am Montag verkauft der Bauer Kakao auf dem Markt.} « Le lundi, le paysan vend du cacao au marché. »

3. \all{Im Sommer besuchen wir eine Fabrik.} « En été, nous visitons une usine. »
\end{exoresolu}

\begin{exercice}[À vous de jouer]
1. Complétez avec la bonne préposition (\all{am}, \all{im}, \all{um}, \all{seit}, \all{von ... bis}) :

a) \all{... Freitag habe ich Deutsch.} b) \all{... Winter ist es kalt.} c) \all{Der Unterricht beginnt ... 7.30 Uhr.} d) \all{Sie lernt ... zwei Jahren Deutsch.} e) \all{Die Schüler arbeiten ... 8 ... 14 Uhr.}

2. Transformez : mettez le complément de temps en tête.

a) \all{Er arbeitet seit einem Jahr in einer Bank.} b) \all{Wir fahren nächsten Monat nach Douala.}

3. Traduisez en allemand : « Chaque semaine, les paysans vendent du cacao au marché de Douala. »
\end{exercice}

\begin{corrige}[Exercice « À vous de jouer »]
1. a) \all{Am Freitag habe ich Deutsch.} b) \all{Im Winter ist es kalt.} c) \all{Der Unterricht beginnt um 7.30 Uhr.} d) \all{Sie lernt seit zwei Jahren Deutsch.} e) \all{Die Schüler arbeiten von 8 bis 14 Uhr.}

2. a) \all{Seit einem Jahr arbeitet er in einer Bank.} b) \all{Nächsten Monat fahren wir nach Douala.}

3. \all{Jede Woche verkaufen die Bauern Kakao auf dem Markt in Douala.} (Accusatif sans préposition \all{jede Woche} en tête, donc inversion : verbe puis sujet.)
\end{corrige}

\begin{aretenir}[L'essentiel sur les compléments de temps]
\begin{itemize}
\item Ils répondent à \all{Wann?}, \all{Wie lange?}, \all{Wie oft?}.
\item Trois formes : adverbes (\all{heute, oft}), accusatif sans préposition (\all{jeden Tag}), groupes prépositionnels (\all{am Montag, im Sommer, um 8 Uhr, seit einem Jahr, von ... bis}).
\item En tête de phrase : inversion obligatoire — le verbe reste en 2\up{e} position : \all{Am Montag arbeitet der Bauer...}
\item \all{seit} + datif, avec le verbe au présent : \all{Sie arbeitet seit einem Jahr hier.}
\item Ordre des compléments : \cle{TeKaMoLo} (temps, cause, manière, lieu).
\end{itemize}
\end{aretenir}

\coursec{Grammaire 2 : les propositions en wenn}

Dans la section précédente, nous avons appris à placer les \cle{compléments de temps} (\all{am Morgen}, \all{im Sommer}, \all{seit 2020}...) dans la phrase allemande. Ces compléments répondent à la question « quand ? ». Mais on peut aussi exprimer le moment ou la condition d'une action avec une \cle{phrase entière} : « \emph{quand} la récolte arrive, les paysans sont très occupés », « \emph{si} je gagne de l'argent, j'achète une voiture ». En allemand, ce type de phrase s'organise autour d'une conjonction essentielle : \all{wenn}. C'est l'une des subordonnées les plus fréquentes de la langue, et elle obéit à une règle capitale : le verbe conjugué part à la fin.

\courssub{Observation : repérons wenn dans le texte}

Reprenons une phrase du texte support \emph{Die Wirtschaft in Deutschland} :

\begin{exemplebox}[Phrase du texte]
\all{Wenn man in Deutschland arbeitet, arbeitet man meistens im Dienstleistungssektor.}

« Quand on travaille en Allemagne, on travaille le plus souvent dans le secteur des services. »
\end{exemplebox}

Observons attentivement la construction :

\begin{itemize}
\item La phrase contient \textbf{deux verbes conjugués} : \all{arbeitet} (deux fois).
\item Le premier \all{arbeitet} se trouve \textbf{à la toute fin} de la première partie, juste avant la virgule.
\item Le second \all{arbeitet} se trouve \textbf{immédiatement après la virgule}, en première position de la deuxième partie.
\item Les deux verbes se « serrent » donc autour de la virgule : \all{..., arbeitet, arbeitet ...}
\end{itemize}

Voici un court texte qui multiplie les exemples. Lisez-le et soulignez mentalement tous les verbes placés en fin de subordonnée.

\begin{textebox}[Ein Tag auf dem Land und in der Stadt]
\all{Wenn der Sommer kommt, beginnt für die Bauern eine sehr harte Zeit. Wenn die Sonne früh aufgeht, steht Familie Mballa schon um fünf Uhr auf. Wenn sie auf dem Feld arbeitet, hilft auch ihr Sohn nach der Schule. Wenn die Ernte gut ist, verkauft die Familie ihren Kakao auf dem Markt in Yaoundé.}

\all{Ihr Bruder lebt dagegen in der Stadt. Wenn er zur Arbeit fährt, nimmt er immer den Bus, denn er hat kein Auto. Als er früher in Douala wohnte, arbeitete er in einer Fabrik. Als er 2019 nach Yaoundé zog, fand er eine Stelle in einer Bank. Wenn er am Abend nach Hause kommt, ist er oft müde, aber zufrieden. Wenn er Geld spart, möchte er später ein kleines Geschäft eröffnen.}

\all{Wenn man die Stadt und das Land vergleicht, sieht man: beide Sektoren sind wichtig. Wenn die Bauern gut arbeiten, hat die Stadt genug zu essen, und wenn die Industrie produziert, finden die jungen Leute Arbeit.}
\end{textebox}

\textbf{Glossaire :} \all{die Ernte, -n} (la récolte) ; \all{der Kakao, -s} (le cacao) ; \all{die Fabrik, -en} (l'usine) ; \all{die Stelle, -n} (le poste) ; \all{das Geschäft, -e} (le commerce, le magasin) ; \all{eröffnen} (ouvrir) ; \all{vergleichen} (comparer) ; \all{genug} (assez) ; \all{aufstehen (steht auf, stand auf, ist aufgestanden)} (se lever).

\textbf{Questions de compréhension :}
\begin{enumerate}
\item \all{Wann steht Familie Mballa im Sommer auf?}
\item \all{Wo verkauft die Familie ihren Kakao?}
\item \all{Wo arbeitete der Bruder früher?}
\item \all{Was möchte der Bruder später eröffnen?}
\end{enumerate}

\courssub{La règle fondamentale : le verbe conjugué à la fin}

\begin{propriete}[Position du verbe dans la subordonnée en wenn]
\all{Wenn} est une \cle{conjonction de subordination}. Dans la subordonnée qu'elle introduit, le \textbf{verbe conjugué se place en dernière position} :

\begin{center}
\all{Wenn + Sujet + compléments + \textbf{Verbe conjugué (fin)} , ...}
\end{center}

Quand la subordonnée est placée \textbf{en tête de phrase}, elle occupe la position 1 : le verbe de la \textbf{proposition principale} vient donc \textbf{immédiatement après la virgule} (c'est l'inversion sujet-verbe) :

\begin{center}
\all{Wenn er kommt, \textbf{ARBEITEN} wir.} \hspace{1cm} \all{Wenn ich Zeit habe, \textbf{BESUCHE} ich die Fabrik.}
\end{center}

Quand la subordonnée est placée \textbf{après} la principale, celle-ci garde son ordre normal (verbe en position 2) :

\begin{center}
\all{Wir arbeiten, wenn er kommt.} \hspace{1cm} \all{Ich besuche die Fabrik, wenn ich Zeit habe.}
\end{center}
\end{propriete}

\begin{popfigure}
\begin{tikzpicture}[font=\small]
% Subordonnée
\node[draw=popblue, fill=popblueL, rounded corners, minimum height=0.9cm, align=center] (wenn) at (0,0) {\all{Wenn}};
\node[draw=popblue, fill=popblueL, rounded corners, minimum height=0.9cm, align=center] (suj1) at (2.2,0) {Sujet\\ \all{man}};
\node[draw=popblue, fill=popblueL, rounded corners, minimum height=0.9cm, align=center] (compl) at (4.6,0) {Compléments\\ \all{in Deutschland}};
\node[draw=popgreen, fill=popgreenL, rounded corners, minimum height=0.9cm, align=center] (verb1) at (7.2,0) {\textbf{Verbe 1}\\ \all{arbeitet}};
\node[align=center] (virg) at (8.3,0) {\textbf{,}};
\node[draw=poporange, fill=poporangeL, rounded corners, minimum height=0.9cm, align=center] (verb2) at (9.5,0) {\textbf{Verbe 2}\\ \all{arbeitet}};
\node[draw=poppurple, fill=poppurpleL, rounded corners, minimum height=0.9cm, align=center] (suj2) at (11.6,0) {Sujet\\ \all{man}};
\node[draw=poppurple, fill=poppurpleL, rounded corners, minimum height=0.9cm, align=center] (rest) at (13.8,0) {Suite\\ \all{meistens...}};
% Accolades et flèches
\draw[decorate, decoration={brace, amplitude=5pt}, popblue, thick] (-0.6,0.75) -- (8.0,0.75) node[midway, above=7pt, popblue, align=center] {Subordonnée : le verbe va à la fin};
\draw[decorate, decoration={brace, amplitude=5pt}, poppurple, thick] (14.9,-0.75) -- (8.8,-0.75) node[midway, below=7pt, poppurple, align=center] {Principale : inversion, verbe en position 1};
\draw[<->, >=Stealth, popgreen, very thick] (verb1.north) .. controls (7.9,1.6) and (8.9,1.6) .. (verb2.north) node[midway, above=2pt, popdark, align=center] {les deux verbes\\ se touchent};
\end{tikzpicture}
\legende{Structure de la phrase avec \all{wenn} en tête : les deux verbes conjugués se « serrent » autour de la virgule.}
\end{popfigure}

\begin{aretenir}[Le réflexe à acquérir]
Dès que vous voyez ou entendez \all{wenn}, préparez-vous à « envoyer » le verbe conjugué \textbf{à la fin} de la subordonnée. Et si la subordonnée ouvre la phrase, placez le verbe de la principale \textbf{juste après la virgule}. Les deux verbes se touchent : c'est le meilleur contrôle visuel pour vérifier votre phrase.
\end{aretenir}

\courssub{Emploi 1 : wenn exprime la condition}

\all{Wenn} introduit une \cle{condition réelle} : l'action de la principale dépend de la réalisation de l'action de la subordonnée. En français, on traduit par « si ».

\begin{exemplebox}[La condition]
\all{Wenn ich Geld verdiene, kaufe ich ein Auto.} — « Si je gagne de l'argent, j'achète une voiture. »

\all{Wenn die Ernte gut ist, verkaufen die Bauern ihren Kakao auf dem Markt.} — « Si la récolte est bonne, les paysans vendent leur cacao au marché. »

\all{Wenn du fleißig lernst, findest du später eine gute Stelle.} — « Si tu travailles dur, tu trouveras plus tard un bon poste. »

\all{Die Fabrik produziert mehr, wenn die Nachfrage steigt.} — « L'usine produit davantage si la demande augmente. »
\end{exemplebox}

\begin{attention}[Pas de futur après « si »... ni après wenn]
Le français interdit le futur après « si » (« si je gagnerai » est faux) ; on dit « si je gagne ». L'allemand fonctionne exactement pareil : pour parler du futur, on utilise le \textbf{présent} dans les deux propositions : \all{Wenn ich Geld verdiene, kaufe ich ein Auto} — les deux verbes sont au présent, le sens est futur. Ne cherchez donc pas à construire un futur compliqué : le présent suffit.
\end{attention}

\courssub{Emploi 2 : wenn exprime le temps (présent, futur, habitude)}

\all{Wenn} introduit aussi une subordonnée \cle{temporelle} : « quand », « lorsque », « chaque fois que ». C'est le cas pour les faits présents, futurs ou les \textbf{habitudes} (actions répétées).

\begin{exemplebox}[Temps et habitude]
\all{Wenn die Ernte kommt, sind die Bauern sehr beschäftigt.} — « Quand la récolte arrive, les paysans sont très occupés. »

\all{Wenn das Wetter gut ist, arbeiten die Bauern auf dem Feld.} — « Si le temps est beau / Quand il fait beau, les paysans travaillent aux champs. »

\all{Wenn ich Zeit habe, besuche ich die Fabrik.} — « Quand j'ai le temps, je visite l'usine. »

\all{Wenn der Markt öffnet, kommen viele Kunden.} — « Quand le marché ouvre, beaucoup de clients arrivent. »
\end{exemplebox}

\courssub{Emploi 3 : wenn pour une habitude au passé}

Pour une action \textbf{répétée dans le passé} (« chaque fois que... », « quand... [souvent] »), l'allemand utilise encore \all{wenn}, avec les verbes au \textbf{Präteritum}.

\begin{exemplebox}[Habitude au passé]
\all{Wenn er früher zur Arbeit fuhr, nahm er den Bus.} — « Quand il allait au travail autrefois, il prenait le bus. » (chaque jour, régulièrement)

\all{Wenn wir Ferien hatten, besuchten wir unsere Großeltern im Dorf.} — « Quand nous avions des vacances, nous rendions visite à nos grands-parents au village. »

\all{Wenn die Schule zu Ende war, half ich meinem Vater auf dem Feld.} — « Quand l'école était finie, j'aidais mon père aux champs. »
\end{exemplebox}

\courssub{La distinction capitale : wenn ou als ?}

Pour un événement passé \textbf{unique} (une seule fois, un fait ponctuel, une période passée délimitée), l'allemand n'utilise \textbf{jamais} \all{wenn}, mais \all{als}. La règle de construction est la même (verbe conjugué à la fin), seul le choix de la conjonction change.

\begin{exemplebox}[als : événement unique au passé]
\all{Als ich Kind war, besuchte ich meinen Großvater auf dem Bauernhof.} — « Quand j'étais enfant, je rendais visite à mon grand-père à la ferme. »

\all{Als ich in Yaoundé war, besuchte ich den Markt.} — « Quand j'étais à Yaoundé, j'ai visité le marché. »

\all{Als er 2019 nach Yaoundé zog, fand er eine Stelle in einer Bank.} — « Quand il a déménagé à Yaoundé en 2019, il a trouvé un poste dans une banque. »
\end{exemplebox}

\begin{attention}[Le piège numéro 1 des francophones]
En français, « quand » sert pour tout : habitude, passé unique, futur. En allemand, il faut \textbf{choisir} :

\begin{itemize}
\item \textbf{Habitude au passé} (« quand j'étais petit, j'allais \emph{souvent}... ») : \all{wenn} + Präteritum.
\item \textbf{Événement unique au passé} (« quand je suis arrivé \emph{ce jour-là}... ») : \all{als} + Präteritum, \textbf{jamais} \all{wenn}.
\end{itemize}

\all{Als ich in Yaoundé war, besuchte ich den Markt.} est correct ; \emph{\all{Wenn ich in Yaoundé war...}} pour un séjour unique est une faute typique de francophone. Autre piège : ne pas oublier l'\textbf{inversion} après la virgule quand la subordonnée est en tête : on dit \all{Wenn er kommt, arbeiten wir} et non \emph{\all{..., wir arbeiten}}.
\end{attention}

\begin{center}
\begin{tabularx}{\linewidth}{|X|X|X|}
\hline
\rowcolor{popblueL} Situation & Conjonction & Exemple \\
\hline
Condition (si) & \all{wenn} & \all{Wenn ich Geld verdiene, kaufe ich ein Auto.} \\
\hline
Présent / futur (quand) & \all{wenn} & \all{Wenn die Ernte kommt, sind die Bauern beschäftigt.} \\
\hline
Habitude au passé (chaque fois que) & \all{wenn} + Präteritum & \all{Wenn er zur Arbeit fuhr, nahm er den Bus.} \\
\hline
Événement unique au passé & \all{als} + Präteritum & \all{Als ich Kind war, besuchte ich meinen Großvater.} \\
\hline
\end{tabularx}
\end{center}

\begin{center}
\begin{tabularx}{\linewidth}{|X|X|X|}
\hline
\rowcolor{popblueL} Français & Deutsch & Remarque \\
\hline
Si je gagne de l'argent, j'achète une voiture. & \all{Wenn ich Geld verdiene, kaufe ich ein Auto.} & « si » de condition = \all{wenn} ; présent des deux côtés \\
\hline
Quand la récolte arrive, les paysans sont occupés. & \all{Wenn die Ernte kommt, sind die Bauern beschäftigt.} & habitude présente = \all{wenn} \\
\hline
Quand j'étais enfant, je visitais la ferme. & \all{Als ich Kind war, besuchte ich den Bauernhof.} & passé unique = \all{als}, jamais \all{wenn} \\
\hline
Je visite l'usine quand j'ai le temps. & \all{Ich besuche die Fabrik, wenn ich Zeit habe.} & subordonnée après : pas d'inversion \\
\hline
\end{tabularx}
\end{center}

\courssub{Méthode : traduire une phrase avec « si » ou « quand »}

\begin{methode}[Traduire « si / quand » en quatre étapes]
\begin{enumerate}
\item \textbf{Identifier le temps et le sens} : est-ce une condition (« si »), une habitude (présente ou passée) ou un événement passé unique ?
\item \textbf{Choisir la conjonction} : condition, présent, futur, habitude (même passée) $\rightarrow$ \all{wenn} ; événement passé unique $\rightarrow$ \all{als}.
\item \textbf{Construire la subordonnée} : conjonction + sujet + compléments + \textbf{verbe conjugué à la fin}.
\item \textbf{Construire la principale} : si la subordonnée est en tête, verbe de la principale \textbf{juste après la virgule} (inversion) ; sinon, ordre normal. Vérifier que les deux verbes « se touchent » autour de la virgule.
\end{enumerate}
\end{methode}

\begin{exoresolu}[Thème guidé : appliquer la méthode]
Traduisez en allemand : « Si le temps est beau, les paysans travaillent aux champs. »
\tcblower
\textbf{Étape 1} : « si » exprime ici une condition réelle au présent. \textbf{Étape 2} : conjonction = \all{wenn}. \textbf{Étape 3} : subordonnée : \all{wenn das Wetter gut ist} (verbe \all{ist} à la fin). \textbf{Étape 4} : la subordonnée est en tête, donc inversion dans la principale : le verbe \all{arbeiten} vient juste après la virgule, puis le sujet \all{die Bauern}.

\textbf{Solution :} \all{Wenn das Wetter gut ist, arbeiten die Bauern auf dem Feld.}

Contrôle visuel : \all{... gut ist, arbeiten ...} — les deux verbes se touchent autour de la virgule. La phrase est correcte.
\end{exoresolu}

\begin{exoresolu}[Version : de l'allemand vers le français]
Traduisez en français : \all{Als mein Bruder in Douala wohnte, arbeitete er in einer Fabrik.}
\tcblower
\textbf{Analyse} : \all{als} + Präteritum (\all{wohnte}) signale un fait passé unique ou une période passée délimitée ; la principale (\all{arbeitete er}) montre l'inversion après la virgule. \all{die Fabrik} = l'usine.

\textbf{Traduction :} « Quand mon frère habitait à Douala, il travaillait dans une usine. »
\end{exoresolu}

\courssub{Entraînement : thème et version}

\begin{exercice}[Thème : traduisez en allemand]
\begin{enumerate}
\item Si j'ai le temps, je visite l'usine.
\item Quand la récolte arrive, les paysans sont très occupés.
\item Quand j'étais enfant, je visitais souvent le marché de Yaoundé. (événement unique : période de l'enfance)
\item Si tu travailles bien, tu trouves un bon poste.
\end{enumerate}
\end{exercice}

\begin{exercice}[Version : traduisez en français]
\begin{enumerate}
\item \all{Wenn die Bauern auf dem Feld arbeiten, helfen die Kinder nach der Schule.}
\item \all{Als er früher in Douala wohnte, fuhr er jeden Tag mit dem Bus zur Fabrik.}
\item \all{Wenn man in Deutschland arbeitet, arbeitet man meistens im Dienstleistungssektor.}
\item \all{Wenn die Familie Geld spart, eröffnet sie ein kleines Geschäft.}
\end{enumerate}
\end{exercice}

\begin{corrige}[Thème]
\begin{enumerate}
\item \all{Wenn ich Zeit habe, besuche ich die Fabrik.} (condition, présent ; inversion après la virgule)
\item \all{Wenn die Ernte kommt, sind die Bauern sehr beschäftigt.} (habitude présente ; \all{kommt} en fin de subordonnée)
\item \all{Als ich Kind war, besuchte ich oft den Markt in Yaoundé.} (période passée unique $\rightarrow$ \all{als} + Präteritum \all{war}, \all{besuchte})
\item \all{Wenn du gut arbeitest, findest du eine gute Stelle.} (condition ; verbe \all{arbeitest} à la fin, inversion : \all{findest du})
\end{enumerate}
\end{corrige}

\begin{corrige}[Version]
\begin{enumerate}
\item « Quand les paysans travaillent aux champs, les enfants aident après l'école. »
\item « Quand il habitait autrefois à Douala, il allait chaque jour à l'usine en bus. »
\item « Quand on travaille en Allemagne, on travaille le plus souvent dans le secteur des services. »
\item « Si la famille économise de l'argent, elle ouvre un petit commerce. »
\end{enumerate}
\end{corrige}

\begin{aretenir}[L'essentiel sur wenn et als]
\begin{itemize}
\item \all{Wenn} et \all{als} sont des conjonctions de subordination : le verbe conjugué de la subordonnée va \textbf{à la fin}.
\item Subordonnée en tête $\rightarrow$ verbe de la principale \textbf{juste après la virgule} : \all{Wenn er kommt, arbeiten wir.}
\item \all{wenn} = condition (« si »), présent, futur, habitude (même au passé).
\item \all{als} = événement passé \textbf{unique}, toujours au Präteritum.
\item Contrôle visuel : les deux verbes se « serrent » autour de la virgule.
\end{itemize}
\end{aretenir}

\coursec{Méthode du commentaire et production guidée}

Dans la section précédente, nous avons appris à construire des propositions subordonnées avec \all{wenn} : \all{Wenn die Bauern mehr Geld verdienen, leben sie besser.} Nous savons aussi placer correctement les compléments de temps. Ces deux outils grammaticaux vont maintenant nous servir à \cle{rédiger} : commenter un texte, le résumer, puis produire un court paragraphe personnel. C'est exactement ce que l'on vous demandera au baccalauréat : comprendre, résumer, réagir. Apprenons la méthode pas à pas.

\courssub{Observer d'abord : à quoi ressemble un bon commentaire ?}

Avant d'énoncer la méthode, lisons attentivement ce court commentaire rédigé par un élève de Première sur le texte étudié dans cette leçon.

\begin{textebox}[Ein Musterkommentar — un commentaire modèle]
Der Text handelt von der Wirtschaft in Deutschland. Zuerst beschreibt der Autor die Landwirtschaft, dann die Industrie und schließlich den Dienstleistungssektor. Die meisten Menschen arbeiten im Dienstleistungssektor. Wenn man in Kamerun arbeitet, arbeitet man oft in der Landwirtschaft. Meiner Meinung nach sind alle drei Sektoren wichtig für ein Land.

\emph{Traduction :} Le texte traite de l'économie en Allemagne. D'abord, l'auteur décrit l'agriculture, puis l'industrie et enfin le secteur des services. La plupart des gens travaillent dans le secteur des services. Quand on travaille au Cameroun, on travaille souvent dans l'agriculture. À mon avis, les trois secteurs sont importants pour un pays.
\end{textebox}

\textbf{Observons.} Ce commentaire de cinq phrases contient déjà tout ce qu'il faut :
\begin{itemize}
\item une phrase d'\cle{introduction} qui présente le thème : \all{Der Text handelt von...} ;
\item un \cle{résumé} du contenu avec les connecteurs \all{zuerst}, \all{dann}, \all{schließlich} ;
\item une \cle{analyse} rapide : une information-clé du texte (\all{Die meisten Menschen arbeiten im Dienstleistungssektor}) et une comparaison avec le Cameroun grâce à une subordonnée en \all{wenn} ;
\item un \cle{avis personnel} introduit par \all{Meiner Meinung nach}.
\end{itemize}

Remarquez aussi ce que l'élève ne fait \emph{pas} : il ne recopie pas le texte, il ne traduit pas phrase par phrase, il ne donne pas son avis dès la première ligne. C'est toute la différence entre un commentaire et une simple paraphrase.

\courssub{La méthode en quatre étapes}

\begin{popfigure}
\begin{tikzpicture}[
  box/.style={draw=popblue, fill=popblueL, rounded corners, thick, align=center, minimum width=3.2cm, minimum height=1.1cm, font=\small},
  arr/.style={-{Stealth[length=3mm]}, thick, popdark}
]
\node[box, align=center] (e1) at (0,0){\textbf{1. Einleitung}\\ Présenter le texte\\ \all{Der Text handelt von...}};
\node[box, fill=popgreenL, draw=popgreen, align=center] (e2) at (4.6,0){\textbf{2. Inhalt}\\ Résumer le contenu\\ \all{zuerst, dann, schließlich}};
\node[box, fill=poporangeL, draw=poporange, align=center] (e3) at (9.2,0){\textbf{3. Analyse}\\ Vocabulaire et structures\\ mots-clés, \all{wenn}};
\node[box, fill=poppinkL, draw=poppink, align=center] (e4) at (13.8,0){\textbf{4. Meinung}\\ Avis personnel\\ \all{Meiner Meinung nach...}};
\draw[arr] (e1) -- (e2);
\draw[arr] (e2) -- (e3);
\draw[arr] (e3) -- (e4);
\end{tikzpicture}
\legende{Les quatre étapes du commentaire de texte : Einleitung, Inhalt, Analyse, eigene Meinung.}
\end{popfigure}

\begin{methode}[Rédiger un commentaire de texte en 4 étapes]
\begin{enumerate}
\item \textbf{Einleitung (introduction).} Commence toujours par présenter le thème du texte : \all{Der Text handelt von den drei Wirtschaftssektoren in Deutschland.} Une seule phrase suffit. Tu peux ajouter le type de texte si tu le connais : \all{Der Text ist ein Artikel.} (Le texte est un article.)
\item \textbf{Inhalt (contenu).} Résume le texte en 2 ou 3 phrases \emph{avec tes propres mots}, sans recopier. Utilise les connecteurs chronologiques : \all{zuerst} (d'abord), \all{dann} (puis), \all{danach} (ensuite), \all{schließlich} (enfin). Attention : ces connecteurs sont en position 1, donc le verbe vient juste après : \all{Zuerst beschreibt der Autor die Landwirtschaft.}
\item \textbf{Analyse.} Choisis 2 ou 3 mots-clés du texte (par exemple \all{die Landwirtschaft}, \all{der Dienstleistungssektor}) et commente-les : que signifient-ils ? Peux-tu comparer avec la situation du Cameroun ? C'est ici qu'une subordonnée en \all{wenn} est la bienvenue : \all{Wenn man in Kamerun arbeitet, arbeitet man oft in der Landwirtschaft.}
\item \textbf{Eigene Meinung (avis personnel).} Termine par ton avis en une ou deux phrases : \all{Meiner Meinung nach sind alle drei Sektoren wichtig.} ou \all{Ich finde es wichtig, dass junge Leute einen Beruf lernen.}
\end{enumerate}
\end{methode}

\begin{aretenir}[Les expressions utiles du commentaire]
\begin{itemize}
\item \all{Der Text handelt von...} — Le texte traite de... (+ datif)
\item \all{Der Autor sagt, dass...} — L'auteur dit que... (subordonnée, verbe à la fin)
\item \all{Der Autor beschreibt / erklärt / zeigt...} — L'auteur décrit / explique / montre...
\item \all{Meiner Meinung nach...} — À mon avis...
\item \all{Ich finde es wichtig, dass...} — Je trouve important que...
\item \all{Ich denke, dass...} — Je pense que...
\end{itemize}
\end{aretenir}

\begin{center}
\begin{tabularx}{\linewidth}{|X|X|X|}
\hline
\rowcolor{popblueL} \textbf{Étape} & \textbf{Expression allemande} & \textbf{Français} \\
\hline
Einleitung & Der Text handelt von der Wirtschaft. & Le texte traite de l'économie. \\
\hline
Inhalt & Zuerst beschreibt der Autor die Landwirtschaft. & D'abord, l'auteur décrit l'agriculture. \\
\hline
Inhalt & Dann erklärt er die Industrie. & Puis il explique l'industrie. \\
\hline
Analyse & Der Autor sagt, dass viele Menschen in Büros arbeiten. & L'auteur dit que beaucoup de gens travaillent dans des bureaux. \\
\hline
Meinung & Meiner Meinung nach ist die Landwirtschaft wichtig. & À mon avis, l'agriculture est importante. \\
\hline
Meinung & Ich finde es wichtig, dass alle Kinder zur Schule gehen. & Je trouve important que tous les enfants aillent à l'école. \\
\hline
\end{tabularx}
\end{center}

\begin{exemplebox}[Des phrases modèles à réutiliser]
\all{Meiner Meinung nach ist die Landwirtschaft in Kamerun sehr wichtig.} « À mon avis, l'agriculture est très importante au Cameroun. »

\all{Wenn die Bauern mehr Geld verdienen, leben sie besser.} « Si les paysans gagnent plus d'argent, ils vivent mieux. »

\all{Der Autor sagt, dass der Dienstleistungssektor in Deutschland sehr groß ist.} « L'auteur dit que le secteur des services est très grand en Allemagne. »

\all{Ich finde es wichtig, dass Kamerun mehr Fabriken baut.} « Je trouve important que le Cameroun construise plus d'usines. »
\end{exemplebox}

\begin{attention}[Le piège de « Meiner Meinung nach »]
Dans l'expression \all{Meiner Meinung nach}, la préposition \all{nach} est \cle{postposée} : elle se place \emph{après} le nom, et non avant. On écrit donc : \all{Meiner Meinung nach ist das wichtig.} La forme \all{nach meiner Meinung} existe dans la langue familière parlée, mais elle est à \textbf{éviter à l'écrit} dans ce sens. Deuxième piège : après \all{Meiner Meinung nach} en tête de phrase, le verbe vient immédiatement (position 2) : \all{Meiner Meinung nach \textbf{sind} alle Sektoren wichtig}, et non *\all{Meiner Meinung nach alle Sektoren sind wichtig}.
\end{attention}

\courssub{Exercice résolu : construire un commentaire}

\begin{exoresolu}[Remettre le commentaire en ordre]
Voici les cinq phrases d'un commentaire dans le désordre. Remets-les dans le bon ordre (Einleitung → Inhalt → Analyse → Meinung) :

a) \all{Meiner Meinung nach ist die Industrie wichtig für Kamerun.}

b) \all{Zuerst beschreibt der Autor die Landwirtschaft, dann die Industrie.}

c) \all{Der Text handelt von den Wirtschaftssektoren in Deutschland.}

d) \all{Der Autor sagt, dass die meisten Menschen im Dienstleistungssektor arbeiten.}

e) \all{Schließlich erklärt er den Dienstleistungssektor.}

\tcblower
\textbf{Solution détaillée.} On applique la méthode en quatre étapes :

\textbf{1. Einleitung :} phrase (c) — elle présente le thème avec \all{Der Text handelt von...}.

\textbf{2. Inhalt :} phrases (b) puis (e) — les connecteurs nous guident : \all{zuerst} (d'abord) → \all{dann} (puis) → \all{schließlich} (enfin). La phrase (b) doit donc précéder la phrase (e).

\textbf{3. Analyse :} phrase (d) — elle reprend une information-clé du texte avec \all{Der Autor sagt, dass...} (verbe \all{arbeiten} à la fin de la subordonnée).

\textbf{4. Meinung :} phrase (a) — l'avis personnel vient toujours en dernier, introduit par \all{Meiner Meinung nach}.

\textbf{Ordre correct : c – b – e – d – a.}
\end{exoresolu}

\courssub{Production guidée 1 : le résumé en 3-4 phrases}

Le résumé (\all{die Zusammenfassung}) est plus court que le commentaire : on ne donne \emph{pas} son avis, on reformule seulement le contenu essentiel avec ses propres mots et des connecteurs.

\begin{methode}[Résumer un texte en 4 phrases]
\begin{enumerate}
\item Repère les 3 ou 4 idées principales du texte (une par paragraphe).
\item Reformule chaque idée avec des mots simples que tu maîtrises — ne recopie jamais une phrase entière.
\item Relie les phrases avec \all{zuerst}, \all{dann}, \all{danach}, \all{schließlich} (verbe en position 2 après le connecteur).
\item Relis : verbe conjugué en position 2, verbe à la fin dans les subordonnées, majuscules aux noms.
\end{enumerate}
\end{methode}

\begin{exemplebox}[Un résumé modèle du texte « Deutschlands Wirtschaft »]
\all{Zuerst erklärt der Text die drei Wirtschaftssektoren. Dann beschreibt er die Landwirtschaft und die Industrie in Deutschland. Danach zeigt er, dass die meisten Menschen im Dienstleistungssektor arbeiten. Schließlich vergleicht er Deutschland und Kamerun.}

« D'abord, le texte explique les trois secteurs économiques. Puis il décrit l'agriculture et l'industrie en Allemagne. Ensuite, il montre que la plupart des gens travaillent dans le secteur des services. Enfin, il compare l'Allemagne et le Cameroun. »
\end{exemplebox}

\courssub{Production guidée 2 : « Die Wirtschaft in Kamerun »}

À ton tour ! Rédige un court paragraphe de 5 à 6 phrases sur l'économie du Cameroun. Consignes obligatoires :
\begin{itemize}
\item utilise au moins \textbf{trois mots du vocabulaire du thème} : \all{die Landwirtschaft}, \all{die Industrie}, \all{der Dienstleistungssektor}, \all{der Bauer (-n)}, \all{die Fabrik (-en)}, \all{der Beruf (-e)}, \all{der Markt („e)}, \all{arbeiten}, \all{verdienen} ;
\item écris au moins \textbf{une subordonnée en} \all{wenn} (verbe à la fin) ;
\item place au moins \textbf{un complément de temps} correctement (en position 1 avec inversion, ou en position 3/4) ;
\item soigne l'orthographe des articles et les majuscules des noms.
\end{itemize}

\begin{exemplebox}[Une production modèle]
\all{In Kamerun arbeiten viele Menschen in der Landwirtschaft. Die Bauern produzieren Kakao, Kaffee und Bananen. Am Montag verkaufen sie ihre Produkte auf dem Markt. Die Industrie ist klein, aber es gibt Fabriken in Douala und Yaoundé. Wenn die Bauern mehr Geld verdienen, leben ihre Familien besser. Meiner Meinung nach ist die Landwirtschaft sehr wichtig für unser Land.}

« Au Cameroun, beaucoup de gens travaillent dans l'agriculture. Les paysans produisent du cacao, du café et des bananes. Le lundi, ils vendent leurs produits au marché. L'industrie est petite, mais il y a des usines à Douala et à Yaoundé. Si les paysans gagnent plus d'argent, leurs familles vivent mieux. À mon avis, l'agriculture est très importante pour notre pays. »

\emph{Analyse du modèle :} trois mots du thème (\all{Landwirtschaft}, \all{Bauern}, \all{Fabriken}, \all{Markt}), un complément de temps en position 1 avec inversion (\all{Am Montag verkaufen sie...}), une subordonnée en \all{wenn} avec le verbe à la fin (\all{verdienen} / \all{leben}), et un avis final avec \all{Meiner Meinung nach}.
\end{exemplebox}

\begin{exercice}[À toi d'écrire]
Rédige ton propre paragraphe \all{Die Wirtschaft in Kamerun} (5-6 phrases) en respectant les quatre consignes ci-dessus. Tu peux parler de ta région : le cacao du Centre, le port de Douala, les marchés de Yaoundé, les plantations de bananes du Littoral.
\end{exercice}

\begin{corrige}[Production guidée 2 — éléments de correction]
Il n'existe pas une seule bonne réponse, mais vérifie les points suivants :
\begin{itemize}
\item \textbf{Vocabulaire :} au moins 3 mots du thème, avec le bon article (\all{die Landwirtschaft}, \all{der Bauer}, \all{die Fabrik}, \all{der Markt}).
\item \textbf{Subordonnée en wenn :} le verbe conjugué est bien à la fin : \all{Wenn die Bauern mehr Geld verdienen, ...} Si la subordonnée ouvre la phrase, le verbe de la principale suit immédiatement : \all{..., leben sie besser.}
\item \textbf{Complément de temps :} en position 1 avec inversion (\all{Am Montag verkaufen die Bauern...}) ou après le verbe (\all{Die Bauern verkaufen am Montag...}).
\item \textbf{Orthographe :} majuscule à tous les noms, Umlaute corrects (\all{Märkte}, \all{Büros}).
\end{itemize}
\end{corrige}

\courssub{Grille d'auto-évaluation}

Avant de rendre ta production, évalue-toi toi-même avec cette grille. Coche chaque critère seulement si tu peux le montrer dans ton texte.

\begin{center}
\begin{tabularx}{\linewidth}{|X|l|}
\hline
\rowcolor{popblueL} \textbf{Critère} & \textbf{Oui / Non} \\
\hline
J'ai utilisé au moins 3 mots du vocabulaire du thème (secteurs, métiers). & \\
\hline
J'ai écrit une subordonnée en \all{wenn} avec le verbe à la fin. & \\
\hline
J'ai placé un complément de temps correctement (inversion si position 1). & \\
\hline
J'ai vérifié les articles (\all{der / die / das}) et les majuscules des noms. & \\
\hline
J'ai terminé par mon avis avec \all{Meiner Meinung nach} ou \all{Ich finde...}. & \\
\hline
Je me suis relu au moins une fois (2 minutes). & \\
\hline
\end{tabularx}
\end{center}

\begin{savaistu}[La production écrite au baccalauréat]
Au Bac, la production écrite est évaluée sur trois critères : la \cle{richesse du vocabulaire} (utiliser les mots du thème, pas seulement \all{gut} et \all{schön}), la \cle{correction grammaticale} (place du verbe, terminaisons, articles) et le \cle{respect de la consigne} (nombre de phrases, sujet traité). Un texte simple mais correct vaut toujours mieux qu'un texte ambitieux plein de fautes. Dernier conseil de correcteur : relis-toi toujours 2 minutes — la moitié des points perdus viennent d'étourderies que tu aurais vues toi-même !
\end{savaistu}

\courssub{Production orale : le métier d'un membre de ma famille}

\begin{experience}[Présentation en 1 minute : « Der Beruf in meiner Familie »]
Prépare une courte présentation orale (environ 1 minute, 6 à 8 phrases) sur le métier d'un membre de ta famille (père, mère, oncle, tante, grand-parent...). Structure imposée :
\begin{enumerate}
\item \textbf{Einleitung :} \all{Mein Vater / Meine Mutter ist ... von Beruf.} (Mon père / ma mère est ... de métier.)
\item \textbf{Der Sektor :} \all{Er / Sie arbeitet in der Landwirtschaft / in der Industrie / im Dienstleistungssektor.}
\item \textbf{Die Tätigkeit :} \all{Er / Sie arbeitet auf dem Feld / in einer Fabrik / in einem Büro / auf dem Markt.}
\item \textbf{Un complément de temps :} \all{Am Morgen fährt er / sie nach Douala.} ou \all{Jeden Tag arbeitet er / sie acht Stunden.}
\item \textbf{Une phrase en wenn :} \all{Wenn er / sie viel arbeitet, verdient er / sie mehr Geld.}
\item \textbf{Meinung :} \all{Meiner Meinung nach ist dieser Beruf wichtig / interessant / schwer.}
\end{enumerate}
Travaille d'abord par écrit (5 minutes), puis présente à ton voisin sans lire, enfin devant la classe. Le voisin écoute et vérifie avec la grille d'auto-évaluation : a-t-on entendu un mot du thème, une phrase en \all{wenn}, un complément de temps ?
\end{experience}

\begin{exemplebox}[Une présentation orale modèle]
\all{Meine Tante ist Verkäuferin von Beruf. Sie arbeitet im Dienstleistungssektor. Jeden Tag verkauft sie Gemüse auf dem Markt in Yaoundé. Am Morgen steht sie um sechs Uhr auf. Wenn sie viele Kunden hat, verdient sie mehr Geld. Meiner Meinung nach ist ihr Beruf schwer, aber sehr wichtig.}

« Ma tante est vendeuse de métier. Elle travaille dans le secteur des services. Chaque jour, elle vend des légumes au marché de Yaoundé. Le matin, elle se lève à six heures. Quand elle a beaucoup de clients, elle gagne plus d'argent. À mon avis, son métier est dur, mais très important. »
\end{exemplebox}

\begin{aretenir}[L'essentiel de la section]
\begin{itemize}
\item Un commentaire de texte suit toujours 4 étapes : \textbf{Einleitung} (\all{Der Text handelt von...}) → \textbf{Inhalt} (résumé avec \all{zuerst, dann, danach, schließlich}) → \textbf{Analyse} (mots-clés, comparaison, phrase en \all{wenn}) → \textbf{Meinung} (\all{Meiner Meinung nach...}).
\item On résume avec ses propres mots, sans jamais recopier le texte.
\item \all{Meiner Meinung nach} : la préposition \all{nach} est postposée, et le verbe suit immédiatement.
\item Toute production (résumé, paragraphe, présentation orale) doit contenir le vocabulaire du thème, au moins une subordonnée en \all{wenn} et un complément de temps bien placé — puis être relue avec la grille d'auto-évaluation.
\end{itemize}
\end{aretenir}

\newpage

\coursec{Activité d'intégration}

\coursec{Activité d'intégration : „Die Berufe unserer Familien“}

\courssub{Objectifs de l'activité}

Cette activité d'intégration clôt la leçon sur les secteurs d'activité. Elle vise à faire réutiliser, dans une situation de communication réelle, tout ce qui a été appris dans la leçon :

\begin{itemize}
\item \cle{Compréhension orale} : écouter et comprendre les réponses des camarades sur les métiers de leurs familles.
\item \cle{Lexique thématique} : employer correctement le vocabulaire des secteurs (\all{die Landwirtschaft, die Industrie, der Dienstleistungssektor}) et des métiers (\all{der Bauer, der Lehrer, der Verkäufer, der Arbeiter...}).
\item \cle{Grammaire} : placer correctement les compléments de temps dans la phrase et construire une subordonnée en \all{wenn} avec le verbe conjugué en fin de proposition.
\item \cle{Expression écrite} : rédiger un court paragraphe de synthèse de 5 à 6 phrases.
\item \cle{Expression orale} : présenter les résultats devant la classe de manière claire et audible.
\item \cle{Compétence sociale} : travailler en équipe, répartir les rôles, évaluer le travail des autres.
\end{itemize}

\courssub{Situation}

Votre classe de Première A du lycée de Yaoundé participe à la \all{Woche der Berufe} (« la semaine des métiers »), une journée portes ouvertes organisée avec le club d'allemand. Des élèves de Seconde et de Terminale visiteront votre salle. Votre professeur vous propose de préparer une grande affiche intitulée \all{„Die Berufe unserer Familien“} (« Les métiers de nos familles ») : chaque groupe de 4 élèves enquête sur les métiers des parents, oncles et tantes de ses membres, classe chaque métier dans l'un des trois secteurs d'activité, puis rédige un court texte de synthèse. Un porte-parole par groupe présentera les résultats à la classe. À la fin, la classe vote pour le groupe le plus précis grammaticalement.

Pour vous aider à démarrer, voici le modèle d'interview que le professeur a affiché au tableau :

\begin{textebox}[Modèle d'interview — „Die Berufe unserer Familien“]
\textbf{A:} Hallo Marie! Was macht dein Vater?\\
\textbf{B:} Mein Vater ist Bauer. Er arbeitet auf dem Feld in der Nähe von Yaoundé.\\
\textbf{A:} Wann beginnt seine Arbeit?\\
\textbf{B:} Er beginnt um 6 Uhr morgens und arbeitet bis 15 Uhr.\\
\textbf{A:} Und deine Mutter? Wo arbeitet sie?\\
\textbf{B:} Sie ist Verkäuferin. Sie arbeitet in einem Geschäft am Markt. Wenn viele Kunden kommen, arbeitet sie auch am Samstag.\\
\textbf{A:} In welchem Sektor arbeitet dein Vater?\\
\textbf{B:} Er arbeitet im primären Sektor, in der Landwirtschaft. Meine Mutter arbeitet im tertiären Sektor, im Dienstleistungssektor.
\end{textebox}

\begin{definition}[Glossaire du dialogue]
\all{der Bauer, -n} : le paysan, l'agriculteur ; \all{das Feld, -er} : le champ ; \all{in der Nähe von} (+ datif) : près de ; \all{die Verkäuferin, -nen} : la vendeuse ; \all{das Geschäft, -e} : le magasin ; \all{der Markt, die Märkte} : le marché ; \all{der Kunde, -n} : le client ; \all{der primäre Sektor} : le secteur primaire ; \all{der tertiäre Sektor} : le secteur tertiaire ; \all{der Dienstleistungssektor} : le secteur des services.
\end{definition}

\courssub{Consignes}

Travaillez par groupes de 4. Répartissez d'abord les rôles : un \cle{interviewer} (il pose les questions), un \cle{secrétaire} (il note les réponses), un \cle{correcteur} (il vérifie la grammaire du texte final) et un \cle{porte-parole} (il présente les résultats). Suivez les étapes dans l'ordre.

\begin{enumerate}
\item \textbf{Préparation de l'interview (5 min).} Dans votre groupe, lisez le modèle d'interview affiché au tableau. Repérez les questions utiles (\all{Was macht...? Wo arbeitet...? Wann beginnt...?}) et recopiez-les sur votre cahier. Ajoutez une question de votre choix, par exemple : \all{Arbeitet dein Onkel in einer Fabrik?} « Ton oncle travaille-t-il dans une usine ? »

\item \textbf{L'interview (10 min).} Chaque membre du groupe répond aux questions sur les métiers de deux membres de sa famille (père, mère, oncle, tante, grand-parent). Le secrétaire note les réponses en allemand, en phrases complètes : \all{Maries Vater ist Bauer. Er arbeitet auf dem Feld.} « Le père de Marie est paysan. Il travaille au champ. »

\item \textbf{Le classement par secteurs (5 min).} Classez chaque métier dans l'un des trois secteurs sur votre feuille, puis venez inscrire vos résultats dans le grand tableau affiché au mur :

\begin{center}
\begin{tabularx}{\linewidth}{|X|X|X|}
\hline
\rowcolor{popblueL} \all{Primärer Sektor} (Landwirtschaft) & \all{Sekundärer Sektor} (Industrie) & \all{Tertiärer Sektor} (Dienstleistungen) \\
\hline
\all{der Bauer} & \all{der Arbeiter in der Fabrik} & \all{der Lehrer} \\
\hline
 & & \\
\hline
\end{tabularx}
\end{center}

\item \textbf{La rédaction de la synthèse (15 min).} Rédigez ensemble un paragraphe de \textbf{5 à 6 phrases} qui présente les résultats de votre enquête. Votre texte doit contenir obligatoirement :
\begin{itemize}
\item au moins \textbf{trois métiers} avec leur secteur d'activité ;
\item au moins \textbf{deux compléments de temps} (\all{am Morgen, um 6 Uhr, jeden Tag, von Montag bis Freitag...}) placés correctement dans la phrase ;
\item au moins \textbf{une subordonnée en \all{wenn}} avec le verbe conjugué à la fin, par exemple : \all{Wenn mein Onkel auf dem Feld arbeitet, beginnt er um 6 Uhr.} « Quand mon oncle travaille au champ, il commence à 6 heures. »
\end{itemize}

\item \textbf{La vérification (5 min).} Le correcteur relit le texte à voix haute. Vérifiez ensemble : la place du verbe (2\textsuperscript{e} position dans la phrase principale, dernière position après \all{wenn}), la majuscule à tous les noms, les articles (\all{der, die, das}) et la place des compléments de temps.

\item \textbf{La présentation orale et le vote (10 min).} Le porte-parole de chaque groupe lit la synthèse devant la classe, lentement et clairement. Les autres groupes écoutent et notent sur une fiche : les métiers cités, les secteurs, les erreurs éventuelles. À la fin, la classe vote à main levée pour le groupe le plus précis grammaticalement. Le groupe gagnant colle sa synthèse sur l'affiche de la \all{Woche der Berufe}.
\end{enumerate}

\courssub{Ressources à mobiliser}

\begin{aretenir}[Le lexique des secteurs et des métiers]
\begin{itemize}
\item Les trois secteurs : \all{der primäre Sektor} (\all{die Landwirtschaft} : l'agriculture), \all{der sekundäre Sektor} (\all{die Industrie} : l'industrie), \all{der tertiäre Sektor} (\all{der Dienstleistungssektor} : les services).
\item Les métiers : \all{der Bauer / die Bäuerin} (paysan/paysanne), \all{der Arbeiter / die Arbeiterin} (ouvrier/ouvrière), \all{der Lehrer / die Lehrerin} (professeur), \all{der Verkäufer / die Verkäuferin} (vendeur/vendeuse), \all{der Arzt / die Ärztin} (médecin), \all{der Fahrer / die Fahrerin} (chauffeur/chauffeuse), \all{der Koch / die Köchin} (cuisinier/cuisinière).
\item Les lieux de travail : \all{auf dem Feld} (au champ), \all{in der Fabrik} (à l'usine), \all{im Büro} (au bureau), \all{in der Schule} (à l'école), \all{am Markt} (au marché), \all{im Krankenhaus} (à l'hôpital).
\item Règle utile : le féminin des métiers se forme presque toujours en ajoutant \all{-in} au masculin : \all{der Lehrer} $\rightarrow$ \all{die Lehrerin}, \all{der Verkäufer} $\rightarrow$ \all{die Verkäuferin}.
\end{itemize}
\end{aretenir}

\begin{propriete}[Rappel 1 : les compléments de temps]
Le complément de temps se place juste après le verbe conjugué (ou en tête de phrase, mais alors le verbe reste en 2\textsuperscript{e} position et le sujet se place après le verbe : c'est l'inversion). Ordre habituel : \textbf{sujet + verbe + temps + lieu}.
\begin{itemize}
\item \all{Mein Vater arbeitet \textbf{jeden Tag} auf dem Feld.} « Mon père travaille chaque jour au champ. »
\item \all{\textbf{Am Montag} beginnt die Arbeit um 7 Uhr.} « Lundi, le travail commence à 7 heures. » (complément de temps en tête $\rightarrow$ verbe en 2\textsuperscript{e} position, sujet après le verbe)
\item Quelques expressions utiles : \all{am Morgen / am Nachmittag / am Abend} (le matin / l'après-midi / le soir), \all{um 6 Uhr} (à 6 heures), \all{von Montag bis Freitag} (du lundi au vendredi), \all{am Wochenende} (le week-end), \all{jeden Tag} (chaque jour).
\end{itemize}
\end{propriete}

\begin{propriete}[Rappel 2 : la subordonnée en wenn]
\all{Wenn} (« quand, si ») introduit une subordonnée : le verbe conjugué va \textbf{à la fin} de la subordonnée. Si la subordonnée est en tête de phrase, le verbe de la principale vient juste après la virgule (schéma : \all{Wenn ..., verbe + sujet ...}).
\begin{itemize}
\item \all{Wenn mein Onkel auf dem Feld \textbf{arbeitet}, \textbf{beginnt} er um 6 Uhr.} « Quand mon oncle travaille au champ, il commence à 6 heures. »
\item \all{Meine Mutter arbeitet am Samstag, wenn viele Kunden \textbf{kommen}.} « Ma mère travaille le samedi quand beaucoup de clients viennent. »
\end{itemize}
\end{propriete}

\begin{attention}[Pièges à éviter]
\begin{itemize}
\item Ne dites pas \all{*Mein Vater er ist Bauer} : un seul sujet ! Écrivez \all{Mein Vater ist Bauer.}
\item Après \all{wenn}, le verbe ne reste pas en 2\textsuperscript{e} position : pas \all{*Wenn mein Onkel arbeitet auf dem Feld...} mais \all{Wenn mein Onkel auf dem Feld arbeitet...}
\item N'oubliez pas la majuscule aux noms de métiers et aux noms en général : \all{der Lehrer}, \all{die Fabrik}, \all{der Markt}.
\item Attention au genre : \all{die Mutter arbeitet}, \all{der Vater arbeitet} ; et au féminin des métiers : \all{die Lehrerin}, \all{die Verkäuferin}.
\item Faux ami : \all{die Fabrik} signifie « l'usine », et non « la fabrique » au sens de manufacture artisanale ; \all{der Fabrikarbeiter} est « l'ouvrier d'usine ».
\item Ne confondez pas \all{der Beruf} (le métier, la profession) et \all{die Arbeit} (le travail) : \all{Was ist dein Vater von Beruf?} « Quel est le métier de ton père ? »
\end{itemize}
\end{attention}

\courssub{Proposition de production et corrigé commenté}

\begin{exoresolu}[Synthèse modèle du groupe 3]
Rédigez la synthèse de votre enquête (5 à 6 phrases) avec au moins trois métiers, deux compléments de temps et une phrase en \all{wenn}.
\tcblower
\textbf{Production modèle :}

\all{In unserer Gruppe arbeiten viele Familien in verschiedenen Sektoren. Der Vater von Marie ist Bauer und arbeitet im primären Sektor: er beginnt jeden Tag um 6 Uhr auf dem Feld. Die Mutter von Paul ist Verkäuferin am Markt von Yaoundé; sie arbeitet von Montag bis Samstag im tertiären Sektor. Wenn viele Kunden kommen, bleibt sie am Abend länger im Geschäft. Der Onkel von Amina ist Arbeiter in einer Fabrik in Douala; er arbeitet im sekundären Sektor. Am Wochenende besuchen unsere Familien oft den Markt.}

« Dans notre groupe, beaucoup de familles travaillent dans différents secteurs. Le père de Marie est paysan et travaille dans le secteur primaire : il commence chaque jour à 6 heures au champ. La mère de Paul est vendeuse au marché de Yaoundé ; elle travaille du lundi au samedi dans le secteur tertiaire. Quand beaucoup de clients viennent, elle reste plus longtemps au magasin le soir. L'oncle d'Amina est ouvrier dans une usine à Douala ; il travaille dans le secteur secondaire. Le week-end, nos familles vont souvent au marché. »

\textbf{Commentaire des choix de langue :}
\begin{itemize}
\item \textbf{Lexique} : les trois secteurs sont cités (\all{primär, sekundär, tertiär}) avec trois métiers explicites (\all{Bauer, Verkäuferin, Arbeiter}) et quatre lieux (\all{auf dem Feld, am Markt, im Geschäft, in einer Fabrik}).
\item \textbf{Compléments de temps} : \all{jeden Tag um 6 Uhr} et \all{von Montag bis Samstag} sont placés juste après le verbe conjugué ; \all{am Wochenende} est en tête de phrase, donc le verbe \all{besuchen} reste en 2\textsuperscript{e} position et le sujet \all{unsere Familien} suit le verbe : inversion correcte.
\item \textbf{Subordonnée en wenn} : \all{Wenn viele Kunden kommen, bleibt sie...} — le verbe \all{kommen} est bien à la fin de la subordonnée, et la principale commence par le verbe \all{bleibt} suivi du sujet \all{sie}.
\item \textbf{Accords et genres} : \all{der Vater von Marie}, \all{die Mutter von Paul}, \all{der Onkel von Amina} (construction avec \all{von} + datif, correcte à l'oral et acceptée à ce niveau) ; le métier féminin est correctement formé : \all{Verkäuferin}.
\item \textbf{Contexte} : les lieux \all{Yaoundé} et \all{Douala} ancrent le texte dans la réalité camerounaise des élèves.
\end{itemize}
\end{exoresolu}

\courssub{Bilan de l'activité}

À l'issue de cette enquête, chaque élève a mobilisé l'ensemble des savoir-faire de la leçon dans une communication authentique : comprendre des questions et des réponses orales sur les métiers, employer le vocabulaire des trois secteurs d'activité et des professions, placer correctement les compléments de temps et construire une subordonnée en \all{wenn}, rédiger un paragraphe cohérent et le présenter à l'oral.

\begin{experience}[Grille d'auto-évaluation]
Cochez ce que vous avez réussi :
\begin{itemize}
\item J'ai posé au moins trois questions sur les métiers en allemand.
\item J'ai classé correctement les métiers dans les trois secteurs.
\item Mon texte contient au moins deux compléments de temps bien placés.
\item Ma phrase en \all{wenn} a le verbe conjugué à la fin.
\item J'ai mis une majuscule à tous les noms et le bon article devant chaque métier.
\item J'ai écouté les autres groupes et participé au vote.
\end{itemize}
\end{experience}

\begin{savaistu}[Et au Cameroun ?]
Au Cameroun comme en Allemagne, on peut classer les métiers en trois secteurs : l'agriculture (le cacao, le café, le manioc) relève du secteur primaire, les usines et la construction du secteur secondaire, et les écoles, les banques, les transports — comme les bendskins de Douala — du secteur tertiaire. Dans les pays industrialisés comme l'Allemagne, la majorité des actifs travaillent dans le secteur tertiaire ; au Cameroun, l'agriculture occupe encore une large part des familles. Votre affiche de classe reflète sans doute cette réalité : comparez vos résultats avec ceux d'une classe allemande !
\end{savaistu}

L'affiche \all{„Die Berufe unserer Familien“} restera exposée dans la salle pendant la semaine des métiers : elle servira de support de révision du vocabulaire et des deux règles de grammaire avant l'évaluation du chapitre.

\newpage

\coursec{Méthodes et exercices résolus}

\courssub{Fiches méthode et exercices résolus}

\begin{methode}[Comprendre un texte sur les secteurs d'activité]
Pour exploiter un texte économique en allemand, procède toujours en quatre étapes :
\begin{enumerate}
\item \textbf{Lecture globale} : lis le titre et le premier paragraphe. Identifie le thème (\all{die Wirtschaft}, \all{die Landwirtschaft}, \all{die Industrie}, \all{der Dienstleistungssektor}) et le type de texte (article, reportage, interview).
\item \textbf{Repérage des mots-clés} : souligne le vocabulaire du thème (\all{der Bauer}, \all{die Fabrik}, \all{der Beruf}, \all{der Arbeiter}, \all{die Firma}) ainsi que les chiffres et les noms de lieux.
\item \textbf{Lecture détaillée} : pour chaque question, localise le passage du texte qui contient la réponse. Ne réponds jamais de mémoire : appuie-toi sur une phrase précise du texte.
\item \textbf{Rédaction de la réponse} : réponds par une phrase complète en allemand, en reprenant les mots de la question. Une réponse commence par une majuscule et finit par un point ; le verbe conjugué est en deuxième position.
\end{enumerate}
\textbf{Réflexes} : un nom allemand prend toujours une majuscule ; attention aux faux amis (\all{die Fabrik} = l'usine, pas « la fabrique » au sens abstrait ; \all{der Beruf} = le métier, la profession).
\end{methode}

\begin{exoresolu}[Questions de compréhension sur un texte]
\textbf{Texte :} \all{Deutschland hat drei wichtige Wirtschaftssektoren. Im ersten Sektor, der Landwirtschaft, arbeiten die Bauern. Sie produzieren Getreide, Obst und Gemüse. Im zweiten Sektor, der Industrie, arbeiten viele Menschen in Fabriken. Sie produzieren Autos und Maschinen. Der dritte Sektor ist der Dienstleistungssektor. Hier arbeiten Lehrer, Ärzte und Verkäufer. Heute arbeiten die meisten Deutschen in diesem Sektor.}\par
\textbf{Questions :} 1) \all{Wie viele Wirtschaftssektoren hat Deutschland?} 2) \all{Wo arbeiten die Bauern?} 3) \all{Was produzieren die Arbeiter in den Fabriken?}
\tcblower
\textbf{Raisonnement} : la question 1 porte sur le nombre de secteurs (\all{wie viele} = « combien de ») : la première phrase du texte donne la réponse (\all{drei}). La question 2 demande un lieu (\all{wo} = « où ») : on cherche le passage sur les \all{Bauern} : \all{in der Landwirtschaft}, c'est-à-dire \all{im ersten Sektor}. La question 3 demande un objet (\all{was} = « que / quoi ») : le texte dit \all{Sie produzieren Autos und Maschinen}.\par
\textbf{Réponses rédigées :}\par
1) \all{Deutschland hat drei wichtige Wirtschaftssektoren.} « L'Allemagne a trois secteurs économiques importants. »\par
2) \all{Die Bauern arbeiten in der Landwirtschaft, im ersten Sektor.} « Les paysans travaillent dans l'agriculture, dans le premier secteur. »\par
3) \all{In den Fabriken produzieren die Arbeiter Autos und Maschinen.} « Dans les usines, les ouvriers produisent des voitures et des machines. »\par
\textbf{Conclusion} : chaque réponse reprend les mots de la question, avec le verbe conjugué en deuxième position (attention à la phrase 3 : le complément de lieu en tête provoque l'inversion sujet-verbe : \all{produzieren die Arbeiter}).
\end{exoresolu}

\begin{methode}[Employer les compléments de temps dans une phrase]
Le complément de temps répond à la question \all{wann?} (quand ?). Deux points à maîtriser : sa forme et sa place.
\begin{enumerate}
\item \textbf{Les compléments de temps simples} : \all{am Montag} (le lundi), \all{im Juli} (en juillet), \all{um 8 Uhr} (à 8 heures), \all{heute} (aujourd'hui), \all{jetzt} (maintenant), \all{früher} (autrefois).
\item \textbf{La place dans la phrase} : au milieu de la phrase, le complément de temps se place juste après le verbe conjugué, avant le complément de lieu : \all{Ich arbeite heute in der Fabrik.} « Je travaille aujourd'hui à l'usine. » En tête de phrase, il provoque l'inversion du sujet et du verbe : \all{Heute arbeite ich in der Fabrik.}
\item \textbf{Les prépositions de temps à connaître} : \all{am} + jour ou date ; \all{im} + mois, saison ou année ; \all{um} + heure ; \all{seit} + point de départ (depuis) ; \all{vor} + durée écoulée (il y a).
\item \textbf{Vérification} : demande-toi toujours « quand ? » ; si tu mets le temps en tête, vérifie que le verbe reste en deuxième position et que le sujet vient après le verbe.
\end{enumerate}
\begin{center}
\begin{tabularx}{\linewidth}{|X|X|X|}\hline
\rowcolor{popblueL} Préposition & Emploi & Exemple \\ \hline
\all{am} (+ datif) & jours et dates & \all{am Montag}, \all{am 1. Mai} \\ \hline
\all{im} (+ datif) & mois, saisons, années & \all{im Juli}, \all{im Sommer}, \all{im Jahr 2024} \\ \hline
\all{um} (+ accusatif) & heures & \all{um 8 Uhr} \\ \hline
\all{seit} (+ datif) & depuis (point de départ) & \all{seit 2020}, \all{seit einem Jahr} \\ \hline
\all{vor} (+ datif) & il y a (durée écoulée) & \all{vor zwei Jahren} \\ \hline
\end{tabularx}
\end{center}
\end{methode}

\begin{exoresolu}[Thème : phrases avec compléments de temps]
Traduisez en allemand : 1) « Aujourd'hui, beaucoup de Camerounais travaillent dans l'agriculture. » 2) « Autrefois, mon grand-père travaillait dans une usine. » 3) « Le lundi, les magasins ouvrent à 8 heures. »
\tcblower
\textbf{Raisonnement} : chaque phrase commence par un complément de temps en tête : il faut donc appliquer l'inversion sujet-verbe (le verbe reste en deuxième position). « Aujourd'hui » = \all{heute} ; « autrefois » = \all{früher} ; « le lundi » = \all{am Montag} ; « à 8 heures » = \all{um 8 Uhr}. Pour la phrase 2, « travaillait » est un passé : au niveau A2, on utilise le Perfekt ou le Präteritum ; le Präteritum de \all{arbeiten} est \all{arbeitete} (avec un -e- intercalaire après -t-).\par
\textbf{Traductions :}\par
1) \all{Heute arbeiten viele Kameruner in der Landwirtschaft.} (inversion : \all{arbeiten} avant \all{viele Kameruner} ; \all{der Kameruner, -} = le Camerounais)\par
2) \all{Früher arbeitete mein Großvater in einer Fabrik.} (Präteritum de \all{arbeiten} ; \all{in einer Fabrik} : datif après \all{in} pour indiquer le lieu où l'on est)\par
3) \all{Am Montag öffnen die Geschäfte um 8 Uhr.} (\all{am} + jour ; \all{um} + heure ; \all{das Geschäft, -e} = le magasin)\par
\textbf{Conclusion} : le temps en tête n'est jamais suivi d'une virgule en allemand, et il entraîne toujours l'inversion du sujet et du verbe.
\end{exoresolu}

\begin{methode}[Construire une subordonnée en wenn]
La conjonction \all{wenn} (quand / si) introduit une subordonnée de temps ou de condition.
\begin{enumerate}
\item \textbf{Règle d'or} : dans la subordonnée en \all{wenn}, le verbe conjugué va à la \textbf{fin} de la proposition : \all{Wenn ich Zeit habe, besuche ich meinen Onkel.} « Quand j'ai le temps, je rends visite à mon oncle. »
\item \textbf{Subordonnée en tête} : si la phrase commence par la subordonnée, le verbe de la principale vient juste après la virgule (il reste en deuxième position de la phrase entière) : \all{Wenn es regnet, bleiben die Bauern zu Hause.} « Quand il pleut, les paysans restent à la maison. »
\item \textbf{Subordonnée après la principale} : pas d'inversion dans la principale : \all{Die Bauern bleiben zu Hause, wenn es regnet.}
\item \textbf{Distinction wenn / als} : \all{wenn} = « quand / si » (répétition, condition, présent ou futur) ; \all{als} = « quand » pour un événement unique au passé : \all{Als ich ein Kind war, ...} « Quand j'étais enfant, ... »
\end{enumerate}
\begin{center}
\begin{tabularx}{\linewidth}{|X|X|X|}\hline
\rowcolor{popblueL} Français & Deutsch & Remarque \\ \hline
si / quand (présent, futur, habitude) & \all{wenn} & verbe conjugué à la fin \\ \hline
quand (événement unique au passé) & \all{als} & verbe conjugué à la fin \\ \hline
\end{tabularx}
\end{center}
\end{methode}

\begin{exoresolu}[Transformation : relier deux phrases par wenn]
Reliez les deux phrases avec \all{wenn} : 1) \all{Es regnet. Die Bauern können nicht auf dem Feld arbeiten.} 2) \all{Die Industrie produziert viele Waren. Die Wirtschaft wächst.} 3) Traduisez : « Si le temps est beau, les paysans travaillent dans les champs. »
\tcblower
\textbf{Raisonnement} : on identifie d'abord la condition (elle devient la subordonnée en \all{wenn}) et la conséquence (la principale). Dans la subordonnée, le verbe conjugué passe à la fin. Si la subordonnée est en tête, la principale commence par son verbe.\par
\textbf{Solutions :}\par
1) \all{Wenn es regnet, können die Bauern nicht auf dem Feld arbeiten.} « Quand il pleut, les paysans ne peuvent pas travailler dans les champs. » (verbe \all{können} en tête de principale après la virgule ; l'infinitif \all{arbeiten} reste à la fin)\par
2) \all{Wenn die Industrie viele Waren produziert, wächst die Wirtschaft.} « Quand l'industrie produit beaucoup de marchandises, l'économie croît. » (\all{produziert} à la fin de la subordonnée ; inversion dans la principale : \all{wächst die Wirtschaft})\par
3) \all{Wenn das Wetter schön ist, arbeiten die Bauern auf den Feldern.} (attention : « si » de condition = \all{wenn}, jamais \all{als} ; « les champs » = \all{die Felder}, pluriel de \all{das Feld, -er})\par
\textbf{Conclusion} : la virgule entre les deux propositions est obligatoire en allemand, et le verbe conjugué de la subordonnée se place toujours en dernière position.
\end{exoresolu}

\begin{methode}[Version : traduire un passage sur la vie économique]
Pour traduire de l'allemand vers le français un texte sur les secteurs d'activité :
\begin{enumerate}
\item \textbf{Repère la structure} de chaque phrase : sujet, verbe conjugué (2\up{e} position), compléments. Dans une subordonnée, cherche le verbe à la fin.
\item \textbf{Traduis d'abord le sens}, mot à mot si nécessaire, puis reformule en français naturel : \all{die Landwirtschaft} = « l'agriculture », \all{der Dienstleistungssektor} = « le secteur des services / le secteur tertiaire », \all{die Ware, -n} = « la marchandise ».
\item \textbf{Méfie-toi des faux amis} : \all{die Fabrik} = « l'usine » ; \all{der Bauer, -n} = « le paysan, l'agriculteur » ; \all{der Betrieb, -e} = « l'entreprise ».
\item \textbf{Relis} ta traduction : le français doit être correct, avec les bons articles et un vocabulaire économique précis.
\end{enumerate}
\end{methode}

\begin{exoresolu}[Version : un paragraphe sur l'économie]
Traduisez en français : \all{In Kamerun spielt die Landwirtschaft eine große Rolle. Viele Familien arbeiten auf den Feldern und produzieren Kakao und Kaffee. Wenn die Ernte gut ist, verdienen die Bauern viel Geld. Aber heute arbeiten immer mehr junge Leute im Dienstleistungssektor, zum Beispiel in Banken oder in Geschäften in Douala und Yaoundé.}
\tcblower
\textbf{Raisonnement} : \all{eine große Rolle spielen} = « jouer un grand rôle » ; \all{auf den Feldern} = « dans les champs » ; \all{die Ernte, -n} = « la récolte » ; \all{Geld verdienen} = « gagner de l'argent » ; \all{immer mehr} = « de plus en plus » ; \all{zum Beispiel} = « par exemple ». La subordonnée \all{Wenn die Ernte gut ist} a son verbe \all{ist} à la fin : « quand la récolte est bonne » ; la principale qui suit commence par le verbe \all{verdienen} : cette inversion disparaît en français.\par
\textbf{Traduction :} « Au Cameroun, l'agriculture joue un grand rôle. Beaucoup de familles travaillent dans les champs et produisent du cacao et du café. Quand la récolte est bonne, les paysans gagnent beaucoup d'argent. Mais aujourd'hui, de plus en plus de jeunes travaillent dans le secteur des services, par exemple dans des banques ou dans des magasins à Douala et à Yaoundé. »\par
\textbf{Conclusion} : une bonne version respecte le sens exact du texte allemand et s'exprime dans un français économique naturel (« le secteur des services », et non un calque mot à mot).
\end{exoresolu}

\begin{methode}[Rédiger un court paragraphe sur les secteurs d'activité]
Pour produire 5 à 8 phrases sur ce thème (expression écrite guidée) :
\begin{enumerate}
\item \textbf{Plan simple} : introduction (le pays et son économie) → secteur primaire (agriculture) → secteur secondaire (industrie) → secteur tertiaire (services) → conclusion personnelle.
\item \textbf{Structures utiles} : \all{In Kamerun arbeiten viele Menschen in der Landwirtschaft.} « Au Cameroun, beaucoup de gens travaillent dans l'agriculture. » ; \all{Der wichtigste Sektor ist ...} « Le secteur le plus important est ... » ; \all{Mein Vater / Meine Mutter arbeitet als ...} « Mon père / ma mère travaille comme ... » ; \all{Wenn ich groß bin, möchte ich ... werden.} « Quand je serai grand, je voudrais devenir ... »
\item \textbf{Connecteurs} : \all{zuerst} (d'abord), \all{dann} (ensuite), \all{auch} (aussi), \all{aber} (mais), \all{deshalb} (c'est pourquoi).
\item \textbf{Vérification finale} : verbe en 2\up{e} position, majuscule à tous les noms, articles corrects (\all{der Bauer}, \all{die Fabrik}, \all{das Feld}), subordonnées en \all{wenn} avec le verbe à la fin.
\end{enumerate}
\end{methode}

\begin{exoresolu}[Rédaction guidée : l'économie de mon pays]
Rédigez un court paragraphe (5 à 7 phrases) en allemand : « Les secteurs d'activité au Cameroun ». Utilisez au moins un complément de temps et une subordonnée en \all{wenn}.
\tcblower
\textbf{Raisonnement} : on suit le plan de la méthode : introduction, les trois secteurs, conclusion. On place un complément de temps en tête (avec inversion) et une phrase en \all{wenn} (verbe à la fin).\par
\textbf{Production modèle :}\par
\all{Kamerun ist ein Land mit einer reichen Wirtschaft. Heute arbeiten viele Menschen in der Landwirtschaft: Sie produzieren Kakao, Kaffee und Bananen. Die Industrie ist auch wichtig, zum Beispiel in Douala. Wenn die Ernte gut ist, verdienen die Bauern viel Geld. Immer mehr junge Leute arbeiten aber im Dienstleistungssektor, in Banken, Schulen und Geschäften. Ich möchte später auch in diesem Sektor arbeiten.}\par
« Le Cameroun est un pays à l'économie riche. Aujourd'hui, beaucoup de gens travaillent dans l'agriculture : ils produisent du cacao, du café et des bananes. L'industrie est aussi importante, par exemple à Douala. Quand la récolte est bonne, les paysans gagnent beaucoup d'argent. Mais de plus en plus de jeunes travaillent dans le secteur des services, dans des banques, des écoles et des magasins. Moi aussi, je voudrais travailler plus tard dans ce secteur. »\par
\textbf{Conclusion} : le paragraphe contient un temps en tête avec inversion (\all{Heute arbeiten ...}), une subordonnée en \all{wenn} correcte (\all{Wenn die Ernte gut ist, verdienen ...}), et le vocabulaire du thème avec les bons articles.
\end{exoresolu}

\begin{attention}[Les 5 erreurs qui coûtent des points au Bac]
\begin{enumerate}
\item \textbf{Oublier la majuscule aux noms} : on écrit \all{die Landwirtschaft}, \all{die Fabrik}, \all{der Beruf}. Un nom sans majuscule est une faute d'orthographe sanctionnée.
\item \textbf{Verbe mal placé dans la subordonnée en wenn} : jamais une principale sans inversion après la subordonnée, jamais un verbe conjugué au milieu de la subordonnée. Correct : \all{Wenn es regnet, arbeiten die Bauern nicht auf dem Feld.} Le verbe conjugué de la subordonnée va à la fin, et la principale commence par son verbe.
\item \textbf{Pas d'inversion après un complément de temps en tête} : une phrase comme « aujourd'hui beaucoup de gens travaillent... » calquée du français est fausse. Correct : \all{Heute arbeiten viele Menschen in der Landwirtschaft.} Le verbe reste en deuxième position.
\item \textbf{Faux amis et calques du français} : \all{die Fabrik} = l'usine (pas « la fabrique ») ; \all{der Bauer} = le paysan ; « travailler dans l'agriculture » se dit \all{in der Landwirtschaft arbeiten} (datif de lieu), jamais avec l'accusatif.
\item \textbf{Confondre wenn et als} : pour « si / quand » au présent ou au futur, on utilise \all{wenn} ; \all{als} est réservé à un événement unique du passé : \all{Als ich klein war, ...} « Quand j'étais petit, ... » mais \all{Wenn ich Zeit habe, ...} « Quand j'ai le temps, ... »
\end{enumerate}
\end{attention}

\newpage

\coursec{Exercices}

\courssub{Je vérifie mes connaissances}

\begin{exercice}[QCM : le vocabulaire des secteurs]
Entoure la bonne réponse.
\begin{enumerate}
\item \all{Der Bauer arbeitet im \ldots\ Sektor.} a) \all{primären} \quad b) \all{sekundären} \quad c) \all{tertiären}
\item \all{Die Fabrik gehört zur \ldots} a) \all{Landwirtschaft} \quad b) \all{Industrie} \quad c) \all{Dienstleistung}
\item \all{Ein Lehrer arbeitet im \ldots} a) \all{Dienstleistungssektor} \quad b) \all{primären Sektor} \quad c) \all{Bergbau}
\item Quel est le pluriel de \all{der Beruf} ? a) \all{Berufs} \quad b) \all{Berufe} \quad c) \all{Berüfe}
\end{enumerate}
\end{exercice}

\begin{exercice}[Vrai ou faux ?]
Réponds par \all{richtig} ou \all{falsch} et justifie ta réponse par une courte phrase en allemand.
\begin{enumerate}
\item \all{Die Landwirtschaft gehört zum sekundären Sektor.}
\item \all{In Deutschland arbeiten viele Menschen im Dienstleistungssektor.}
\item \all{Der Bauer produziert Autos.}
\item \all{Kakao ist ein wichtiges Produkt der Landwirtschaft in Kamerun.}
\end{enumerate}
\end{exercice}

\begin{exercice}[Vocabulaire : articles et pluriels]
Complète le tableau avec l'article et le pluriel de chaque nom.
\begin{center}
\begin{tabularx}{\linewidth}{|X|X|X|}
\hline
\rowcolor{popblueL} Nom & Article & Pluriel \\
\hline
\all{Bauer} & \ldots & \ldots \\
\hline
\all{Fabrik} & \ldots & \ldots \\
\hline
\all{Beruf} & \ldots & \ldots \\
\hline
\all{Industrie} & \ldots & \ldots \\
\hline
\all{Markt} & \ldots & \ldots \\
\hline
\end{tabularx}
\end{center}
\end{exercice}

\begin{exercice}[Conjugaison à compléter]
Conjugue les verbes entre parenthèses au Präsens.
\begin{enumerate}
\item \all{Die Bauern (arbeiten) jeden Tag auf dem Feld.}
\item \all{Mein Vater (gehen) um 7 Uhr in die Fabrik.}
\item \all{Wenn das Wetter schön (sein), ernten wir den Kakao.}
\item \all{Ihr (besuchen) am Samstag den Markt in Douala.}
\item \all{Die Industrie (produzieren) viele Waren.}
\end{enumerate}
\end{exercice}

\courssub{Je m'entraîne}

\begin{exercice}[Les compléments de temps]
Complète les phrases avec \all{am}, \all{im}, \all{um}, \all{seit} ou un accusatif temporel (\all{jeden Tag}, \all{jede Woche}).
\begin{enumerate}
\item \all{\ldots\ Montag fährt mein Bruder nach Yaoundé.}
\item \all{Die Schule beginnt \ldots\ 8 Uhr.}
\item \all{\ldots\ Sommer arbeiten die Bauern viel auf dem Feld.}
\item \all{Meine Schwester arbeitet \ldots\ zwei Jahren in einer Bank.}
\item \all{\ldots\ besucht er seine Großeltern auf dem Land.} (chaque semaine)
\item \all{\ldots\ trinken wir Kaffee zum Frühstück.} (chaque jour)
\end{enumerate}
\end{exercice}

\begin{exercice}[Transformation : le complément de temps en tête]
Réécris chaque phrase en plaçant le complément de temps souligné en tête. Attention à l'inversion du sujet et du verbe.

Exemple : \all{Die Bauern arbeiten \textbf{im Sommer} viel.} $\rightarrow$ \all{Im Sommer arbeiten die Bauern viel.}
\begin{enumerate}
\item \all{Mein Vater fährt \textbf{um 6 Uhr} zur Arbeit.}
\item \all{Wir besuchen \textbf{am Wochenende} den Markt.}
\item \all{Die Schüler lernen \textbf{jeden Tag} Deutsch.}
\item \all{Meine Tante verkauft \textbf{im Dezember} Kakao.}
\item \all{Die Arbeiter gehen \textbf{am Morgen} in die Fabrik.}
\end{enumerate}
\end{exercice}

\begin{textebox}[Texte à trous : Die Wirtschaft in Kamerun]
In Kamerun ist die \ldots\ sehr wichtig. Viele \ldots\ produzieren Kakao und Kaffee. In Douala gibt es die \ldots : Dort arbeiten die Menschen in einer \ldots\ und produzieren Waren. Der \ldots\ wächst auch: Banken, Schulen und Krankenhäuser bieten viele Jobs. Mein Vater hat einen guten \ldots : Er ist Lehrer.
\end{textebox}

\begin{exercice}[Texte à trous]
Complète le texte ci-dessus avec les mots suivants : \all{Landwirtschaft}, \all{Fabrik}, \all{Dienstleistungssektor}, \all{Bauern}, \all{Industrie}, \all{Beruf}.
\end{exercice}

\begin{exercice}[Appariement : métiers et secteurs]
Associe chaque métier à son secteur (\all{primär}, \all{sekundär}, \all{tertiär}), puis forme le féminin de chaque métier.
\begin{center}
\begin{tabularx}{\linewidth}{|X|X|X|}
\hline
\rowcolor{popblueL} Beruf & Sektor & Féminin \\
\hline
\all{der Bauer} & \ldots & \ldots \\
\hline
\all{der Arbeiter} & \ldots & \ldots \\
\hline
\all{der Lehrer} & \ldots & \ldots \\
\hline
\all{der Fischer} & \ldots & \ldots \\
\hline
\all{der Verkäufer} & \ldots & \ldots \\
\hline
\all{der Arzt} & \ldots & \ldots \\
\hline
\all{der Bäcker} & \ldots & \ldots \\
\hline
\all{der Mechaniker} & \ldots & \ldots \\
\hline
\all{der Koch} & \ldots & \ldots \\
\hline
\all{der Busfahrer} & \ldots & \ldots \\
\hline
\end{tabularx}
\end{center}
\end{exercice}

\courssub{Je me prépare au Bac}

\begin{textebox}[Die Landwirtschaft in Kamerun]
Die Landwirtschaft ist der wichtigste Sektor der Wirtschaft in Kamerun. Viele Familien auf dem Land arbeiten jeden Tag auf den Feldern. Sie produzieren Kakao, Kaffee, Bananen und Maniok. Der Kakao aus Kamerun ist in der ganzen Welt bekannt.

Am Morgen gehen die Bauern früh auf die Felder, oft schon um 6 Uhr. Wenn die Ernte kommt, sind alle sehr beschäftigt: Die ganze Familie hilft. Im Oktober und November verkaufen die Bauern den Kakao an Händler aus Douala oder Yaoundé.

Aber das Leben der Bauern ist nicht leicht. Die Preise für Kakao sind oft niedrig, und die Straßen auf dem Land sind schlecht. Wenn es viel regnet, können die Lastwagen die Produkte nicht transportieren.

Seit einigen Jahren helfen Programme der Regierung den Bauern: Sie bekommen bessere Pflanzen und moderne Maschinen. Viele junge Leute wollen aber nicht auf dem Land bleiben. Sie gehen in die Städte und suchen Arbeit in der Industrie oder im Dienstleistungssektor.
\end{textebox}

\begin{exercice}[Compréhension écrite type Bac]
Lis le texte ci-dessus, puis réponds aux questions en allemand, par des phrases complètes.

\textbf{Questions globales :}
\begin{enumerate}
\item Worum geht es in diesem Text? (2--3 Sätze)
\item Welche Probleme haben die Bauern in Kamerun?
\end{enumerate}
\textbf{Questions de détail :}
\begin{enumerate}
\setcounter{enumi}{2}
\item Welche Produkte produzieren die Familien auf dem Land?
\item Wann verkaufen die Bauern den Kakao?
\item Was machen viele junge Leute? Warum?
\end{enumerate}
\textbf{Vrai ou faux (justifie avec une phrase du texte) :}
\begin{enumerate}
\setcounter{enumi}{5}
\item \all{Die Bauern gehen erst um 10 Uhr auf die Felder.}
\item \all{Die Straßen auf dem Land sind sehr gut.}
\end{enumerate}
\end{exercice}

\begin{exercice}[Thème et version]
\textbf{A. Thème} — Traduis en allemand :
\begin{enumerate}
\item Si le temps est beau, les paysans travaillent aux champs.
\item Quand j'étais enfant, je visitais la ferme de mon grand-père.
\item Chaque jour, les ouvriers vont à l'usine à 7 heures.
\item En décembre, ma famille voyage au village.
\end{enumerate}
\textbf{B. Version} — Traduis en français :
\begin{enumerate}
\item \all{Wenn die Ernte kommt, sind die Bauern sehr beschäftigt.}
\item \all{Seit zwei Jahren arbeitet meine Schwester in einer Bank in Douala.}
\item \all{Am Wochenende verkaufen die Händler ihre Produkte auf dem Markt.}
\end{enumerate}
\end{exercice}

\begin{exercice}[Production écrite guidée type Bac]
Rédige un paragraphe de 6 à 8 phrases (environ 60 à 80 mots) intitulé \all{„Die Wirtschaft in meinem Land“}.

Consignes obligatoires :
\begin{itemize}
\item utilise au moins cinq mots du vocabulaire de la leçon (\all{die Landwirtschaft}, \all{die Industrie}, \all{der Dienstleistungssektor}, \all{der Bauer}, \all{die Fabrik}, \all{der Beruf}, \all{der Markt}\ldots) ;
\item emploie au moins deux compléments de temps différents (\all{am}, \all{im}, \all{um}, \all{seit}, accusatif temporel) ;
\item écris au moins une subordonnée avec \all{wenn} ;
\item respecte la place du verbe dans chaque phrase.
\end{itemize}
\end{exercice}

\begin{methode}[Réussir la production écrite]
\begin{enumerate}
\item \textbf{Planifie} : note d'abord 5 mots de vocabulaire et 2 compléments de temps que tu veux employer.
\item \textbf{Rédige} des phrases courtes : sujet -- verbe en position 2 -- compléments.
\item \textbf{Insère} une subordonnée en \all{wenn} : verbe conjugué à la fin (\all{Wenn die Ernte kommt, \ldots}).
\item \textbf{Relis} : vérifie la majuscule des noms, la place du verbe et les articles (\all{der}, \all{die}, \all{das}).
\end{enumerate}
\textbf{Exemple de production :} \all{In Kamerun ist die Landwirtschaft sehr wichtig. Viele Bauern arbeiten jeden Tag auf dem Feld und produzieren Kakao. Am Morgen gehen sie schon um 6 Uhr zur Arbeit. Wenn die Ernte kommt, hilft die ganze Familie. In Douala gibt es auch viele Fabriken: Die Industrie wächst. Seit einigen Jahren arbeiten immer mehr Menschen im Dienstleistungssektor, zum Beispiel in Banken und Schulen. Mein Bruder hat einen guten Beruf: Er ist Mechaniker.}
\end{methode}

\newpage

\coursec{Corrigés des exercices}

\coursec{Corrigés détaillés — Les secteurs d'activité (ALA07)}

\begin{corrige}[Exercice 1 — QCM : le vocabulaire des secteurs]
\begin{enumerate}
\item \all{primären} — le paysan (\all{der Bauer}) travaille dans le secteur primaire, c'est-à-dire l'agriculture (\all{die Landwirtschaft}).
\item \all{Industrie} — l'usine (\all{die Fabrik}) appartient au secteur secondaire, celui de la transformation.
\item \all{Dienstleistungssektor} — l'enseignant travaille dans le secteur tertiaire (services).
\item \all{Berufe} — pluriel régulier en \all{-e}, sans Umlaut : \all{der Beruf, die Berufe}.
\end{enumerate}
\textbf{Point de vigilance :} ne pas confondre \all{der Beruf} (le métier) et \all{die Arbeit} (le travail) ; et attention au faux ami : \all{die Fabrik} = l'usine, pas « la fabrique » au sens abstrait.
\end{corrige}

\begin{corrige}[Exercice 2 — Vrai ou faux ?]
\begin{enumerate}
\item \all{Falsch. Die Landwirtschaft gehört zum primären Sektor.} (L'agriculture appartient au secteur primaire.)
\item \all{Richtig. In Deutschland arbeiten die meisten Menschen im Dienstleistungssektor.} (En Allemagne, la plupart des gens travaillent dans le secteur des services.)
\item \all{Falsch. Der Bauer produziert Lebensmittel, zum Beispiel Getreide, Kartoffeln und Milch.} (Le paysan produit des denrées alimentaires, par exemple du blé, des pommes de terre et du lait.)
\item \all{Richtig. Kamerun produziert viel Kakao.} (Le Cameroun produit beaucoup de cacao.)
\end{enumerate}
\textbf{Point de vigilance :} dans la justification, le verbe conjugué reste en deuxième position ; \all{zum} = \all{zu} + \all{dem} (datif après \all{gehören zu}).
\end{corrige}

\begin{savaistu}[Contexte camerounais et allemand]
En Allemagne, les agriculteurs produisent principalement des céréales (\all{Getreide}), des pommes de terre (\all{Kartoffeln}), du lait (\all{Milch}), de la viande et des légumes. Le cacao (\all{Kakao}) et le café (\all{Kaffee}) sont des produits tropicaux importés ; le Cameroun, la Côte d'Ivoire et le Ghana en sont les grands exportateurs. C'est pourquoi l'exercice 7 et le texte de compréhension (exercice 9) portent sur la réalité camerounaise, tandis que le texte support principal \all{Die Wirtschaft in Deutschland} décrit la structure allemande (secteur tertiaire dominant).
\end{savaistu}

\begin{corrige}[Exercice 3 — Vocabulaire : articles et pluriels]
\begin{center}
\begin{tabularx}{\linewidth}{|X|X|X|}
\hline
\rowcolor{popgreenL} Nom & Article & Pluriel \\
\hline
\all{Bauer} & \all{der} & \all{die Bauern} \\
\hline
\all{Fabrik} & \all{die} & \all{die Fabriken} \\
\hline
\all{Beruf} & \all{der} & \all{die Berufe} \\
\hline
\all{Industrie} & \all{die} & \all{die Industrien} \\
\hline
\all{Markt} & \all{der} & \all{die Märkte} \\
\hline
\end{tabularx}
\end{center}
\textbf{Points de vigilance :}
\begin{itemize}
\item \all{der Bauer} est un nom faible : il prend \all{-n} à toutes les formes sauf au nominatif singulier : \all{Ich sehe den Bauern.} (Je vois le paysan.)
\item \all{der Markt} prend l'Umlaut au pluriel : \all{die Märkte} — comme \all{der Stadt} $\rightarrow$ \all{die Städte}.
\item Les noms en \all{-ie} sont féminins et font leur pluriel en \all{-n} : \all{die Industrie, die Industrien}.
\end{itemize}
\end{corrige}

\begin{corrige}[Exercice 4 — Conjugaison à compléter]
\begin{enumerate}
\item \all{Die Bauern \textbf{arbeiten} jeden Tag auf dem Feld.} (3\textsuperscript{e} personne du pluriel, terminaison \all{-en}.)
\item \all{Mein Vater \textbf{geht} um 7 Uhr in die Fabrik.} (3\textsuperscript{e} personne du singulier, verbe \all{gehen}, terminaison \all{-t}.)
\item \all{Wenn das Wetter schön \textbf{ist}, ernten wir den Kakao.} (Dans la subordonnée en \all{wenn}, le verbe conjugué va à la fin.)
\item \all{Ihr \textbf{besucht} am Samstag den Markt in Douala.} (2\textsuperscript{e} personne du pluriel : \all{ihr besucht}.)
\item \all{Die Industrie \textbf{produziert} viele Waren.} (Singulier : terminaison \all{-t}.)
\end{enumerate}
\textbf{Point de vigilance :} après \all{ihr}, la terminaison est \all{-t} (\all{ihr besucht}), et non \all{-en} ; ne pas écrire \all{ihr besuchen}.
\end{corrige}

\begin{corrige}[Exercice 5 — Les compléments de temps]
\begin{enumerate}
\item \all{\textbf{Am} Montag fährt mein Bruder nach Yaoundé.} (Jour de la semaine : \all{am} + datif.)
\item \all{Die Schule beginnt \textbf{um} 8 Uhr.} (Heure précise : \all{um}.)
\item \all{\textbf{Im} Sommer arbeiten die Bauern viel auf dem Feld.} (Saison : \all{im} = \all{in} + \all{dem}.)
\item \all{Meine Schwester arbeitet \textbf{seit} zwei Jahren in einer Bank.} (Durée commencée dans le passé et qui continue : \all{seit} + datif + Präsens.)
\item \all{\textbf{Jede Woche} besucht er seine Großeltern auf dem Land.} (Répétition : accusatif temporel, \all{die Woche} est féminin, donc \all{jede}.)
\item \all{\textbf{Jeden Tag} trinken wir Kaffee zum Frühstück.} (\all{der Tag} est masculin : accusatif \all{jeden}.)
\end{enumerate}
\textbf{Point de vigilance :} le français « depuis deux ans » se rend par \all{seit} + Präsens en allemand, jamais par le Perfekt : \all{Sie arbeitet seit zwei Jahren hier.}
\end{corrige}

\begin{corrige}[Exercice 6 — Transformation : le complément de temps en tête]
\begin{enumerate}
\item \all{Um 6 Uhr fährt mein Vater zur Arbeit.}
\item \all{Am Wochenende besuchen wir den Markt.}
\item \all{Jeden Tag lernen die Schüler Deutsch.}
\item \all{Im Dezember verkauft meine Tante Kakao.}
\item \all{Am Morgen gehen die Arbeiter in die Fabrik.}
\end{enumerate}
\textbf{Rappel de la règle :} quand le complément de temps occupe la première position, le verbe conjugué reste en deuxième position et le sujet passe après le verbe (inversion). Schéma : complément + verbe + sujet + reste.
\end{corrige}

\begin{corrige}[Exercice 7 — Texte à trous]
In Kamerun ist die \all{\textbf{Landwirtschaft}} sehr wichtig. Viele \all{\textbf{Bauern}} produzieren Kakao und Kaffee. In Douala gibt es die \all{\textbf{Industrie}} : Dort arbeiten die Menschen in einer \all{\textbf{Fabrik}} und produzieren Waren. Der \all{\textbf{Dienstleistungssektor}} wächst auch: Banken, Schulen und Krankenhäuser bieten viele Jobs. Mein Vater hat einen guten \all{\textbf{Beruf}} : Er ist Lehrer.

Traduction : « Au Cameroun, l'agriculture est très importante. Beaucoup de paysans produisent du cacao et du café. À Douala, il y a l'industrie : là, les gens travaillent dans une usine et produisent des marchandises. Le secteur des services se développe aussi : les banques, les écoles et les hôpitaux offrent beaucoup d'emplois. Mon père a un bon métier : il est enseignant. »

\textbf{Point de vigilance :} \all{in einer Fabrik} (datif après \all{in} pour le lieu où l'on est) ; \all{einen guten Beruf} (accusatif masculin après \all{haben}).
\end{corrige}

\begin{corrige}[Exercice 8 — Appariement : métiers et secteurs]
\begin{center}
\begin{tabularx}{\linewidth}{|X|X|X|}
\hline
\rowcolor{popgreenL} Beruf & Sektor & Féminin \\
\hline
\all{der Bauer} & \all{primär} & \all{die Bäuerin} \\
\hline
\all{der Arbeiter} & \all{sekundär} & \all{die Arbeiterin} \\
\hline
\all{der Lehrer} & \all{tertiär} & \all{die Lehrerin} \\
\hline
\all{der Fischer} & \all{primär} & \all{die Fischerin} \\
\hline
\all{der Verkäufer} & \all{tertiär} & \all{die Verkäuferin} \\
\hline
\all{der Arzt} & \all{tertiär} & \all{die Ärztin} \\
\hline
\all{der Bäcker} & \all{sekundär} & \all{die Bäckerin} \\
\hline
\all{der Mechaniker} & \all{sekundär} & \all{die Mechanikerin} \\
\hline
\all{der Koch} & \all{tertiär} & \all{die Köchin} \\
\hline
\all{der Busfahrer} & \all{tertiär} & \all{die Busfahrerin} \\
\hline
\end{tabularx}
\end{center}
\textbf{Règle :} le féminin se forme avec le suffixe \all{-in}. Les noms en \all{-er} restent inchangés (\all{der Lehrer} $\rightarrow$ \all{die Lehrerin}), mais certains noms prennent l'Umlaut : \all{der Arzt} $\rightarrow$ \all{die Ärztin}, \all{der Koch} $\rightarrow$ \all{die Köchin}, \all{der Bauer} $\rightarrow$ \all{die Bäuerin}.

\textbf{Point de vigilance :} le boulanger (\all{der Bäcker}) transforme la matière première (farine) : il appartient donc au secteur secondaire, comme le mécanicien ; le cuisinier (\all{der Koch}) offre un service : secteur tertiaire.
\end{corrige}

\begin{corrige}[Exercice 9 — Compréhension écrite type Bac]
\textbf{Réponses rédigées :}
\begin{enumerate}
\item \all{Der Text handelt von der Landwirtschaft in Kamerun. Die Bauern produzieren Kakao und Kaffee, aber ihr Leben ist schwer. Viele junge Leute verlassen das Land und gehen in die Städte.} (Le texte parle de l'agriculture au Cameroun. Les paysans produisent du cacao et du café, mais leur vie est dure. Beaucoup de jeunes quittent la campagne et vont en ville.)
\item \all{Die Preise für Kakao sind oft niedrig, und die Straßen auf dem Land sind schlecht.} (Les prix du cacao sont souvent bas et les routes à la campagne sont mauvaises.)
\item \all{Sie produzieren Kakao, Kaffee, Bananen und Maniok.}
\item \all{Im Oktober und November verkaufen die Bauern den Kakao an Händler aus Douala oder Yaoundé.}
\item \all{Viele junge Leute gehen in die Städte und suchen Arbeit in der Industrie oder im Dienstleistungssektor, weil sie nicht auf dem Land bleiben wollen.}
\item \all{Falsch.} Le texte dit : « \all{Am Morgen gehen die Bauern früh auf die Felder, oft schon um 6 Uhr.} »
\item \all{Falsch.} Le texte dit : « \all{Die Straßen auf dem Land sind schlecht.} »
\end{enumerate}
\textbf{Barème indicatif (sur 10) :} question 1 : 2 pts (idée globale + deux phrases complètes) ; question 2 : 1 pt ; questions 3 à 5 : 1 pt chacune ; questions 6 et 7 : 1,5 pt chacune (0,5 pt pour \all{falsch}, 1 pt pour la citation justificative). On retire 0,25 pt par faute de grammaire ou de place du verbe.

\textbf{Points de vigilance :} répondre par des phrases complètes (sujet + verbe) ; reprendre le vocabulaire du texte sans recopier tout le paragraphe ; pour le vrai/faux, toujours citer la phrase du texte entre guillemets allemands „…“.
\end{corrige}

\begin{corrige}[Exercice 10 — Thème et version]
\textbf{A. Thème (français $\rightarrow$ allemand) :}
\begin{enumerate}
\item \all{Wenn das Wetter schön ist, arbeiten die Bauern auf dem Feld.} (Subordonnée en \all{wenn} : verbe conjugué \all{ist} à la fin ; dans la principale qui suit, le verbe \all{arbeiten} occupe la première position, puis le sujet.)
\item \all{Als ich ein Kind war, besuchte ich den Bauernhof meines Großvaters.} (Moment passé révolu : \all{als} + Präteritum : \all{war}, \all{besuchte}.)
\item \all{Jeden Tag gehen die Arbeiter um 7 Uhr in die Fabrik.} (Accusatif temporel \all{jeden Tag} ; heure : \all{um 7 Uhr}.)
\item \all{Im Dezember fährt meine Familie ins Dorf.} (Mois : \all{im} ; \all{ins} = \all{in} + \all{das}, mouvement, accusatif.)
\end{enumerate}
\textbf{B. Version (allemand $\rightarrow$ français) :}
\begin{enumerate}
\item Quand la récolte arrive, les paysans sont très occupés.
\item Depuis deux ans, ma sœur travaille dans une banque à Douala.
\item Le week-end, les commerçants vendent leurs produits au marché.
\end{enumerate}
\textbf{Barème indicatif (sur 10) :} 1,5 pt par phrase de thème (6 pts) ; environ 1,3 pt par phrase de version (4 pts). Toute erreur de place du verbe dans la subordonnée est sanctionnée.

\textbf{Points de vigilance :} « si » de condition = \all{wenn} ; « quand » au passé = \all{als}, jamais \all{wenn} ; « aux champs » = \all{auf dem Feld} (datif, lieu) ; « au village » = \all{ins Dorf} (accusatif, mouvement).
\end{corrige}

\begin{corrige}[Exercice 11 — Production écrite guidée type Bac]
\textbf{Proposition de rédaction modèle :}

\all{In Kamerun ist die Landwirtschaft sehr wichtig. Viele Bauern arbeiten jeden Tag auf dem Feld und produzieren Kakao und Kaffee. Am Morgen gehen sie schon um 6 Uhr zur Arbeit. Wenn die Ernte kommt, hilft die ganze Familie. In Douala gibt es auch viele Fabriken: Die Industrie wächst. Seit einigen Jahren arbeiten immer mehr Menschen im Dienstleistungssektor, zum Beispiel in Banken und Schulen. Mein Bruder hat einen guten Beruf: Er ist Mechaniker.}

Traduction : « Au Cameroun, l'agriculture est très importante. Beaucoup de paysans travaillent chaque jour aux champs et produisent du cacao et du café. Le matin, ils vont au travail dès 6 heures. Quand la récolte arrive, toute la famille aide. À Douala, il y a aussi beaucoup d'usines : l'industrie se développe. Depuis quelques années, de plus en plus de gens travaillent dans le secteur des services, par exemple dans les banques et les écoles. Mon frère a un bon métier : il est mécanicien. »

\textbf{Vérification des consignes :} vocabulaire de la leçon : \all{die Landwirtschaft}, \all{die Bauern}, \all{die Fabriken}, \all{die Industrie}, \all{der Dienstleistungssektor}, \all{der Beruf} (6 mots, minimum 5 exigé) ; compléments de temps : \all{jeden Tag}, \all{am Morgen}, \all{um 6 Uhr}, \all{seit einigen Jahren} (minimum 2 exigé) ; subordonnée en \all{wenn} : \all{Wenn die Ernte kommt, hilft die ganze Familie.}

\textbf{Barème indicatif (sur 10) :} respect des consignes (vocabulaire, compléments de temps, subordonnée) : 4 pts ; correction grammaticale (place du verbe, conjugaison, cas) : 3 pts ; orthographe (majuscules des noms, Umlaute) : 1,5 pt ; cohérence et longueur (60--80 mots) : 1,5 pt.

\textbf{Points de vigilance :} majuscule à tous les noms (\all{die Ernte}, \all{die Familie}) ; verbe en position 2 dans les principales, à la fin dans la subordonnée ; ne pas traduire mot à mot le français (« très important » = \all{sehr wichtig}, sans \all{viel}).
\end{corrige}

\newpage

\coursec{Fiche bilan}

\begin{popfigure}
\begin{tikzpicture}[mindmap, grow cyclic, every node/.style={concept, text=white, align=center, font=\small},
  root concept/.append style={concept color=popblue, font=\bfseries},
  level 1/.append style={level distance=3.4cm, sibling angle=60, font=\footnotesize},
  level 2/.append style={level distance=2.2cm, font=\tiny}]
\node[root concept, align=center]{Die Wirtschafts-\\sektoren}
  child[concept color=popgreen]{ node[align=center]{Landwirtschaft\\ \all{der Bauer, das Feld}} }
  child[concept color=poporange]{ node[align=center]{Industrie\\ \all{die Fabrik, der Arbeiter}} }
  child[concept color=poppurple]{ node[align=center]{Dienstleistungen\\ \all{der Beruf, das Büro}} }
  child[concept color=popgold]{ node[align=center]{Vocabulaire\\ articles + pluriels} }
  child[concept color=poppink]{ node[align=center]{Compléments\\ de temps\\ \all{am, im, um}} }
  child[concept color=popdark]{ node[align=center]{Propositions\\ en \all{wenn}\\ verbe à la fin} };
\end{tikzpicture}
\legende{Carte mentale de la leçon : les trois secteurs d'activité, le lexique et les deux points de grammaire}
\end{popfigure}

\begin{aretenir}[L'essentiel de la leçon]
\textbf{1. Lexique clé (à connaître avec article et pluriel)}

\begin{center}
\begin{tabularx}{\linewidth}{|X|X|}
\hline
\rowcolor{popblueL} Deutsch & Français \\
\hline
\all{die Landwirtschaft} (f., sg.) & l'agriculture \\
\hline
\all{die Industrie, -n} & l'industrie \\
\hline
\all{der Dienstleistungssektor, -en} & le secteur des services \\
\hline
\all{der Bauer, -n / die Bäuerin, -nen} & l'agriculteur / l'agricultrice \\
\hline
\all{die Fabrik, -en} & l'usine \\
\hline
\all{der Beruf, -e} & le métier, la profession \\
\hline
\all{der Arbeiter, - / die Arbeiterin, -nen} & l'ouvrier / l'ouvrière \\
\hline
\all{das Produkt, -e / produzieren} & le produit / produire \\
\hline
\all{arbeiten (in + Dat.)} & travailler (dans) \\
\hline
\all{die Wirtschaft, -en} & l'économie \\
\hline
\all{das Feld, -er} & le champ \\
\hline
\all{das Büro, -s} & le bureau \\
\hline
\end{tabularx}
\end{center}

\textbf{2. Grammaire à connaître par cœur}

\begin{center}
\begin{tabularx}{\linewidth}{|l|X|X|}
\hline
\rowcolor{popblueL} Règle & Forme & Exemple \\
\hline
Complément de temps & \all{am} + jour, \all{im} + mois/saison, \all{um} + heure & \all{Am Montag arbeite ich. Im Juli reisen wir. Um 8 Uhr beginnt die Arbeit.} \\
\hline
Place du complément de temps & en tête (verbe en 2\textsuperscript{e} position) ou après le verbe conjugué & \all{Heute arbeitet der Bauer auf dem Feld.} \\
\hline
Proposition en \all{wenn} & verbe conjugué à la FIN ; virgule entre les deux propositions & \all{Wenn es regnet, bleiben die Arbeiter zu Hause.} \\
\hline
\all{wenn} en tête & la proposition principale commence par le VERBE & \all{Wenn ich Zeit habe, besuche ich die Fabrik.} \\
\hline
\end{tabularx}
\end{center}

\textbf{3. Méthodes express}
\begin{itemize}
\item \textbf{Comprendre le texte} : lire le titre, repérer les trois secteurs, souligner les mots connus, puis répondre d'abord globalement, ensuite en détail.
\item \textbf{Répondre aux questions} : reprendre les mots de la question, phrase complète, verbe en 2\textsuperscript{e} position.
\item \textbf{Résumé} : 3 à 5 phrases, une idée par secteur, avec \all{zuerst, dann, schließlich}.
\item \textbf{Traduction} : attention à \all{wenn} (verbe final) et aux majuscules des noms.
\end{itemize}
\end{aretenir}

\begin{exemplebox}[Modèles de phrases à réutiliser à l'examen]
\all{In Deutschland gibt es drei Wirtschaftssektoren: die Landwirtschaft, die Industrie und den Dienstleistungssektor.} « En Allemagne, il y a trois secteurs économiques : l'agriculture, l'industrie et le secteur des services. »

\all{Der Bauer arbeitet im Sommer auf dem Feld.} « L'agriculteur travaille au champ en été. »

\all{Wenn die Industrie wächst, finden viele Menschen Arbeit.} « Quand l'industrie se développe, beaucoup de gens trouvent du travail. »

\all{Am Montag um 7 Uhr beginnt die Arbeit in der Fabrik.} « Lundi à 7 heures, le travail commence à l'usine. »
\end{exemplebox}

\begin{attention}[Pièges fréquents des francophones]
\begin{itemize}
\item Ne confonds pas \all{der Beruf} (le métier, la profession) et \all{die Arbeit} (le travail, l'activité) : \all{Er lernt einen Beruf.} « Il apprend un métier. » / \all{Er geht zur Arbeit.} « Il va au travail. »
\item \all{am, im, um} ne sont pas interchangeables : \all{am Montag} (jour), \all{im Juli} (mois), \all{im Sommer} (saison), \all{um 8 Uhr} (heure précise).
\item Après \all{wenn}, le verbe conjugué va à la fin ; si la subordonnée ouvre la phrase, la principale commence par le verbe : \all{Wenn es regnet, \textbf{bleiben} die Arbeiter zu Hause.}
\item Majuscule obligatoire à tous les noms : \all{die Fabrik, der Bauer, das Produkt}.
\item Lieu avec \all{in} : datif (où ?) : \all{in der Fabrik, in dem (im) Büro}.
\end{itemize}
\end{attention}

\begin{aretenir}[Checklist avant le Bac]
\begin{itemize}
\item[$\square$] Je connais les trois secteurs : \all{die Landwirtschaft, die Industrie, der Dienstleistungssektor}.
\item[$\square$] Je cite le lexique avec article et pluriel (\all{der Bauer, -n ; die Fabrik, -en}).
\item[$\square$] Je mets la majuscule à tous les noms allemands.
\item[$\square$] J'emploie \all{am} (jour), \all{im} (mois, saison), \all{um} (heure) sans les confondre.
\item[$\square$] Je place le verbe conjugué en 2\textsuperscript{e} position dans la phrase principale.
\item[$\square$] Dans une proposition en \all{wenn}, je mets le verbe conjugué à la fin.
\item[$\square$] Quand \all{wenn} ouvre la phrase, la principale commence par le verbe : \all{Wenn ..., kommt ...}.
\item[$\square$] Je n'oublie pas la virgule entre la subordonnée et la principale.
\item[$\square$] Je distingue \all{der Beruf} (le métier) et \all{die Arbeit} (le travail).
\item[$\square$] Je sais rédiger un résumé de 3 à 5 phrases avec des connecteurs.
\item[$\square$] Je relis ma production : accords, cas après \all{in} (+ datif de lieu), orthographe.
\end{itemize}
\end{aretenir}

\findecours

\end{document}
